\documentclass[11pt,oneside,openany]{book}
\usepackage[T1]{fontenc}
\usepackage[utf8]{inputenc}
\usepackage{lmodern,amsmath,amssymb,amsthm,mathtools}
\usepackage[letterpaper,margin=1in,headheight=15pt]{geometry}
\usepackage{microtype,graphicx,booktabs,longtable,array,enumitem}
\usepackage{xcolor,fancyhdr,listings}
\usepackage[numbers,sort&compress]{natbib}
\usepackage{url}
\usepackage[colorlinks=true,linkcolor=teal!65!black,citecolor=teal!65!black,urlcolor=blue!55!black]{hyperref}
\hypersetup{pdftitle={Fluid Dynamics, Orders and Quantum Information},pdfauthor={Research synthesis prepared for Amelie Schreiber},pdfsubject={A rigorous research synthesis with explicit proof and verification boundaries}}
\pagestyle{fancy}
\fancyhf{}
\fancyhead[L]{\small Fluid Dynamics and Quantum Information}
\fancyhead[R]{\small Research monograph}
\fancyfoot[C]{\thepage}
\renewcommand{\chaptermark}[1]{\markboth{\thechapter.\ #1}{}}
\setcounter{tocdepth}{1}
\setcounter{secnumdepth}{2}
\makeatletter
\renewcommand*\l@section{\@dottedtocline{1}{1.5em}{3.2em}}
\makeatother
\setlist{itemsep=3pt,topsep=6pt}
\newtheorem{theorem}{Theorem}[chapter]
\newtheorem{lemma}[theorem]{Lemma}
\newtheorem{proposition}[theorem]{Proposition}
\newtheorem{corollary}[theorem]{Corollary}
\newtheorem{conjecture}[theorem]{Conjecture}
\newtheorem{problem}[theorem]{Problem}
\theoremstyle{definition}
\newtheorem{definition}[theorem]{Definition}
\newtheorem{example}[theorem]{Example}
\newtheorem{construction}[theorem]{Construction}
\theoremstyle{remark}
\newtheorem{remark}[theorem]{Remark}
\newtheorem{warning}[theorem]{Caution}
\newcommand{\R}{\mathbb{R}}
\newcommand{\C}{\mathbb{C}}
\newcommand{\Q}{\mathbb{Q}}
\newcommand{\Z}{\mathbb{Z}}
\newcommand{\N}{\mathbb{N}}
\newcommand{\F}{\mathbb{F}}
\newcommand{\HH}{\mathcal{H}}
\newcommand{\OO}{\mathcal{O}}
\newcommand{\eps}{\varepsilon}
\newcommand{\llbracket}{\mathopen{[\![}}
\newcommand{\rrbracket}{\mathclose{]\!]}}
\DeclareMathOperator{\Hom}{Hom}
\DeclareMathOperator{\End}{End}
\DeclareMathOperator{\Ext}{Ext}
\DeclareMathOperator{\Tor}{Tor}
\DeclareMathOperator{\Spec}{Spec}
\DeclareMathOperator{\Tr}{Tr}
\DeclareMathOperator{\tr}{tr}
\DeclareMathOperator{\rank}{rank}
\DeclareMathOperator{\im}{im}
\DeclareMathOperator{\coker}{coker}
\DeclareMathOperator{\ind}{ind}
\DeclareMathOperator{\diag}{diag}
\DeclareMathOperator{\GL}{GL}
\DeclareMathOperator{\SL}{SL}
\DeclareMathOperator{\SU}{SU}
\DeclareMathOperator{\vol}{vol}
\DeclareMathOperator{\Pic}{Pic}
\DeclareMathOperator{\Cl}{Cl}
\DeclareMathOperator{\Gal}{Gal}
\DeclareMathOperator{\Sym}{Sym}
\DeclareMathOperator{\supp}{supp}
\DeclareMathOperator{\dist}{dist}
\DeclareMathOperator{\Span}{span}
\newcommand{\reviewstatus}[1]{\par\smallskip\noindent\textbf{Evidence status.} #1\par\smallskip}
\lstset{basicstyle=\ttfamily\scriptsize,breaklines=true,columns=fullflexible,keepspaces=true,frame=single,rulecolor=\color{black!20},showstringspaces=false}
% Character map for the three authoritative Lean source listings.
% Requires listings and amssymb; load after the main listing style.
% This changes typography only, never the authoritative Lean source.
\lstset{literate=
  {·}{{\ensuremath{\cdot}}}1
  {ε}{{\ensuremath{\varepsilon}}}1
  {μ}{{\ensuremath{\mu}}}1
  {ᴴ}{{\ensuremath{{}^{\mathsf H}}}}1
  {ᵀ}{{\ensuremath{{}^{\mathsf T}}}}1
  {ᵥ}{{\ensuremath{{}_{\mathtt v}}}}1
  {–}{{--}}1
  {•}{{\ensuremath{\bullet}}}1
  {₀}{{\ensuremath{{}_{0}}}}1
  {₁}{{\ensuremath{{}_{1}}}}1
  {₂}{{\ensuremath{{}_{2}}}}1
  {ₗ}{{\ensuremath{{}_{\mathtt l}}}}1
  {ₖ}{{\ensuremath{{}_{\mathtt k}}}}1
  {ℂ}{{\ensuremath{\mathbb C}}}1
  {ℕ}{{\ensuremath{\mathbb N}}}1
  {ℝ}{{\ensuremath{\mathbb R}}}1
  {ℤ}{{\ensuremath{\mathbb Z}}}1
  {←}{{\ensuremath{\leftarrow}}}1
  {→}{{\ensuremath{\rightarrow}}}1
  {↔}{{\ensuremath{\leftrightarrow}}}1
  {∃}{{\ensuremath{\exists}}}1
  {∀}{{\ensuremath{\forall}}}1
  {∈}{{\ensuremath{\in}}}1
  {≃}{{\ensuremath{\simeq}}}1
  {⊗}{{\ensuremath{\otimes}}}1
  {∧}{{\ensuremath{\wedge}}}1
  {∨}{{\ensuremath{\vee}}}1
  {≤}{{\ensuremath{\le}}}1
  {≠}{{\ensuremath{\ne}}}1
  {¬}{{\ensuremath{\neg}}}1
  {×}{{\ensuremath{\times}}}1
  {ν}{{\ensuremath{\nu}}}1
  {∑}{{\ensuremath{\sum}}}1
  {⟨}{{\ensuremath{\langle}}}1
  {⟩}{{\ensuremath{\rangle}}}1
  {⬝}{{\ensuremath{\mathbin{\scriptscriptstyle\blacksquare}}}}1
}

\title{\textbf{Fluid Dynamics, Orders\\and Quantum Information}\\[1em]\Large Microscopic reduction, finite complexes, quantum learning\\and a structural reading of the 2026 Navier--Stokes release}
\author{Prepared for Amelie Schreiber\\[.4em]\normalsize A standalone extension of the research monograph}
\date{13 November 2026\\\small Version 1.4: operator realizations and certified control}
\begin{document}
\frontmatter
\hypersetup{pageanchor=false}
\maketitle
\clearpage
\hypersetup{pageanchor=true}
\chapter*{Purpose and scope}
This supplement develops the fluid-dynamics extension of \emph{Orders, Complexes, Cryptography and Quantum Yang--Mills}. It contains the same seven fluid, reduction, quiver and operator-realization chapters and the fluid verification appendix as the complete monograph, with independent pagination. The bibliography identifies the primary sources, and the companion repository contains the source coverage, exact formal declarations, numerical programs and interactive laboratory.

The contribution is an explicit mathematical interpretation and a collection of proved interfaces. The released continuum singularity theorem is attributed to its source. The selected local Lean declarations and finite computations have their own verification records. We do not present a stronger unforced theorem, a fully formalized version of every chapter, or a trained fluid model where none was executed.
\tableofcontents
\mainmatter
\chapter{Microscopic information, kinetic limits, and fluid closure}
\label{ch:microscopic-fluid-limits}

\reviewstatus{The 249 timed English transcript segments of the requested
braintruffle video, from 0:00 through 32:36, were reviewed in the public
YouTube transcript interface.  The timestamped source audit records precisely
which accompanying papers, lecture decks, and book sections were accessible.
The full main texts of Stodolna, Vozzi, Itatani, Allen, the abridged English
J\"onsson article, and all six requested Volkov decks were reviewed.  Several
other requested originals were accessible only as abstracts or metadata;
these are not represented as full-text reviews.  The arguments proved below
are ordinary mathematical proofs.  Their presence here does not assert a
Lean formalization of the continuum theories.}

\section{What the microscopic perspective establishes}

The video organizes fluid simulation around successive reductions in the
information retained about a physical system: a quantum state, molecular
positions and velocities, a velocity distribution, and a few macroscopic
fields \citep{microBraintruffle2021}.  This is a productive organizing principle
provided that each reduction is supplied with its own assumptions.  An
exact change of coordinates, a statistical marginal, an asymptotic limit,
and an empirical constitutive law are mathematically different operations.
The distinction determines which conclusions can be carried from one level
to another.

For example, expanding a wave function in a complete eigenbasis retains
all information.  Keeping only finitely many modes discards information.
Taking the partial trace of a many-body density operator discards
correlations.  Replacing a quantum nuclear state by a classical trajectory
requires a localization or semiclassical approximation.  Replacing an
$N$-particle classical distribution by a one-particle distribution requires
a closure of its correlation hierarchy.  Replacing that distribution by
density, velocity, and temperature requires another closure.  None of these
steps follows merely from a count of variables.

The central mathematical lesson is therefore a sequence of conditional
statements.  The Born--Huang expansion identifies the terms discarded by an
adiabatic approximation.  Liouville's equation identifies the correlations
discarded by a kinetic approximation.  Collision invariants identify the
moments retained by a fluid approximation.  An entropy inequality restricts
admissible closures but does not determine all their coefficients.  A
hydrodynamic-limit theorem supplies compactness and convergence information
that a formal moment calculation does not supply.  The structural results
in Chapter~\ref{ch:fluid-structures} and the quantum algorithms in
Chapter~\ref{ch:quantum-fluid} fit into these particular stages.

\subsection{Five distinct mathematical state spaces}

For $N$ distinguishable spinless particles in three dimensions, a pure
quantum state is a unit vector in $L^2(\R^{3N})$, up to an overall phase.
An $N$-particle classical microstate is a point of $\R^{6N}$.  An uncertain
classical state is instead a probability measure on that same phase space;
when it has a density, that density is a function of $6N$ variables.  Thus
the video's contrast between a wave function of $3N$ coordinates and a
classical list of $6N$ numbers compares a general quantum state with a
single classical microstate.  A comparison of statistical descriptions must
also include classical ensemble distributions.

The one-particle kinetic density $f(t,x,v)$ is a function on
$\R^3_x\times\R^3_v$.  In the unit-mass convention used below, its basic
moments are
\begin{equation}
 \rho=\int f\,dv,\qquad
 \rho u=\int vf\,dv,\qquad
 3\rho\theta=\int |v-u|^2f\,dv.
 \label{eq:micro-basic-moments}
\end{equation}
Here $\theta=k_BT/m$ in physical units.  The fluid variables
$(\rho,u,\theta)$ contain five scalar fields.  They generally do not
determine $f$, its pressure tensor, or its heat flux.

In the quantum setting, mixed states are positive trace-one operators
$\varrho$.  A reduced one-particle density matrix and a classical kinetic
density are different kinds of objects.  A Wigner transform can relate
them in suitable semiclassical regimes, but a generic Wigner function is
not a nonnegative probability density.  A classical probabilistic
interpretation therefore cannot be inserted without a positivity or
limiting argument.

\section{Measurement, phase, and the limits of tomography}

The experiments cited by the video support wave interference and
model-dependent reconstruction of electronic structure.  J\"onsson's
multi-slit experiment compares electron diffraction with interference
predictions \citep{microJonsson1974}.  Stodolna and collaborators image
nodal structures in an outgoing electron distribution associated with
hydrogen Stark states; the detector distribution is related to a
propagated wave by the apparatus and its model
\citep{microStodolna2013}.  Itatani's molecular-orbital tomography uses
high-harmonic emission, alignment information, a reference atom, and a
recollision approximation \citep{microItatani2004}.  Vozzi's generalized
tomography improves phase recovery and the treatment of Coulomb
distortion while retaining finite-bandwidth and reconstruction assumptions
\citep{microVozzi2011}.

These are not measurements of a particle trajectory or unrestricted
tomography of an arbitrary many-electron wave function.  In particular, a
single orbital reconstruction does not by itself determine every
many-particle correlation.  The distinction can be expressed without
depending on an interpretation of quantum mechanics.

\begin{proposition}[Finite-dimensional informational completeness]
\label{prop:micro-tomography}
Let $E_1,\ldots,E_s$ be positive operators on $\C^d$ satisfying
$\sum_iE_i=I$.  The probabilities
$p_i=\Tr(E_i\varrho)$ uniquely determine every density operator
$\varrho$ if and only if the real linear span of the $E_i$ is the
space of Hermitian $d\times d$ matrices.
\end{proposition}
\begin{proof}
If the span is the full Hermitian space and two states give the same
probabilities, their difference $X$ is orthogonal to every Hermitian
matrix in the inner product $\Tr(XY)$.  Taking $Y=X$ gives
$\Tr(X^2)=0$, hence $X=0$.

Conversely, if the span is proper, choose a nonzero Hermitian matrix $X$
orthogonal to it.  Since $I=\sum_iE_i$ belongs to the span,
$\Tr X=0$.  For sufficiently small $a>0$, both
$\varrho_\pm=I/d\pm aX$ are positive trace-one matrices.  They are
distinct and satisfy $\Tr(E_i\varrho_+)=\Tr(E_i\varrho_-)$ for
every $i$.
\end{proof}

Position probabilities alone fail this criterion.  The qubit vectors
$(|0\rangle+|1\rangle)/\sqrt2$ and
$(|0\rangle-|1\rangle)/\sqrt2$ have identical computational-basis
probabilities and different relative phase.  Similarly,
$\psi(x)$ and $e^{i\alpha(x)}\psi(x)$ have the same $|\psi(x)|^2$
where both define admissible states, but may have different currents and
future dynamics.  The Born density for a complex wave function is
$|\psi|^2$, not its algebraic square.

\begin{proposition}[Local operations do not transmit a signal through a marginal]
\label{prop:micro-nosignalling}
Let $\varrho_{AB}$ be a density operator on a finite tensor product.
For operators $K_j$ on $B$ satisfying $\sum_jK_j^*K_j=I_B$, define
\[
 \widetilde\varrho_{AB}
 =\sum_j(I_A\otimes K_j)\varrho_{AB}(I_A\otimes K_j^*).
\]
Then $\Tr_B\widetilde\varrho_{AB}=\Tr_B\varrho_{AB}$.
\end{proposition}
\begin{proof}
For every observable $A$ on the first factor, cyclicity of the trace gives
\[
 \Tr((A\otimes I)\widetilde\varrho)
 =\sum_j\Tr((A\otimes K_j^*K_j)\varrho)
 =\Tr((A\otimes I)\varrho).
\]
Equality against every $A$ determines the reduced operator.
\end{proof}

Conditioning on a recorded outcome can change a conditional state.  The
outcome record must then be communicated before the other observer can
condition on it.  Bell correlations consequently do not justify an
assertion of controllable instantaneous communication.  This qualification
is needed when using the introductory discussion in van Dommelen's
book \citep{microVanDommelen}.

\subsection{What the multi-slit formula proves}

In the Fraunhofer approximation, suppose $s$ equally illuminated identical
slits contribute a common amplitude $a(\vartheta)$ and a phase increment
$\delta=2\pi d\sin\vartheta/\lambda$ between neighboring slits.
Linearity gives
\[
 A(\vartheta)=a(\vartheta)\sum_{j=0}^{s-1}e^{ij\delta}.
\]
The geometric-series identity and continuity at removable singularities
give
\begin{equation}
 |A(\vartheta)|^2
 =|a(\vartheta)|^2
 \frac{\sin^2(s\delta/2)}{\sin^2(\delta/2)}.
 \label{eq:micro-slits}
\end{equation}
For a uniformly illuminated slit of width $b$, the Fourier integral
$\int_{-b/2}^{b/2}e^{2\pi i x\sin\vartheta/\lambda}\,dx$
supplies a squared sinc envelope.  Equation~\eqref{eq:micro-slits}
therefore proves the interference factor under stated coherence,
illumination, and far-field assumptions.  It does not remove the need to
model the source, finite detector resolution, or partial coherence.
The reviewed 1974 article explicitly identifies itself as an abridged
translation; it cannot stand in for a complete review of the longer
1961 German original \citep{microJonsson1961,microJonsson1974}.

\section{Exterior powers are a physical state constraint}

The project's exterior algebra is directly relevant to identical
fermions.  Let $V$ be a one-particle space including spin.  The
antisymmetric $N$-particle subspace is $\bigwedge^N V$, identified
with the antisymmetric tensors inside $V^{\otimes N}$.  Exchanging two
complete spin-space labels changes a fermionic state by a sign.  Equality
of position probabilities under an exchange, by itself, does not specify
that sign or the spin dependence.

\begin{proposition}[Slater normalization and dimension]
\label{prop:micro-slater}
If $e_1,\ldots,e_M$ are an orthonormal basis of the finite-dimensional
one-particle space $V$, then for
$1\le i_1<\cdots<i_N\le M$ the tensors
\[
 \Psi_{i_1\cdots i_N}
 =\frac1{\sqrt{N!}}\sum_{\sigma\in S_N}
   \operatorname{sgn}(\sigma)
   e_{i_{\sigma(1)}}\otimes\cdots\otimes e_{i_{\sigma(N)}}
\]
form an orthonormal basis of $\bigwedge^N V$.  In particular,
$\dim\bigwedge^N\C^M=\binom MN$, with the dimension zero for
$N>M$.
\end{proposition}
\begin{proof}
For fixed distinct indices the $N!$ elementary tensors in the sum are
mutually orthonormal, so the squared norm is $N!/N!=1$.  Different
index sets have disjoint elementary-tensor supports.  Antisymmetry forces
the coefficient of any tensor with a repeated index to vanish, and relates
all coefficients belonging to permutations of a distinct index set by
the permutation sign.  Thus every antisymmetric tensor is a linear
combination of the displayed vectors.
\end{proof}

In position representation this is the normalized Slater determinant.
The sign representation of the symmetric group selects the Young diagram
$(1^N)$, while symmetric bosonic tensors select $(N)$.  More general
Schur functors organize symmetry types of tensor features and multi-index
responses.  This is a precise use of the algebraic material developed
earlier in the monograph: a representation specifies allowed states or
tensor components.  It does not, on its own, furnish a bound for a nonlinear
fluid equation.

\subsection{Number states and coherent states}

Mathur's accessible introductory modules are useful for interference,
tensor products, and the distinction between classical amplitude and
measurement probabilities \citep{microMathur}.  One teaching
identification needs correction: a state with a fixed number of identical
bosons in one mode is a number state, and is not a coherent state.
Let $a|n\rangle=\sqrt n\,|n-1\rangle$.  A coherent state is
\begin{equation}
 |\alpha\rangle=e^{-|\alpha|^2/2}
       \sum_{n=0}^{\infty}\frac{\alpha^n}{\sqrt{n!}}|n\rangle.
 \label{eq:micro-coherent}
\end{equation}
The exponential series shows that its squared norm is one.  Shifting
the summation index shows $a|\alpha\rangle=\alpha|\alpha\rangle$.
Its number probabilities are
$e^{-|\alpha|^2}|\alpha|^{2n}/n!$, whose mean and variance are both
$|\alpha|^2$.  A fixed $|n\rangle$ has number variance zero and
$\langle a\rangle=0$.  Quantization of the number observable therefore
does not discretize every complex amplitude or every field expectation.
This matters when deciding what a quantum neural model's input amplitudes
and measured outputs represent.

\section{Born--Huang expansion before Born--Oppenheimer truncation}

The video's move from electronic wave functions to molecular potential
energy surfaces is best made precise by first writing the coupling terms
that are later neglected.  The treatment in van Dommelen's Section 9.2
and derivation D.51 distinguishes electronic and nuclear quantum dynamics;
Allen's molecular-dynamics notes then describe classical motion on a chosen
potential \citep{microVanDommelen,microAllen2004}.

For an exact, domain-transparent derivation, take a finite electronic
space $\C^r$, nuclear coordinate $R\in\R^d$, and a $C^2$ family of
Hermitian matrices $H_e(R)$.  Suppose an orthonormal eigenbasis
$\phi_n(R)$ can be chosen $C^2$ on the open region being considered.
For $C^2$ nuclear coefficient functions $\chi_n$, set
\[
 H=-\frac{\hbar^2}{2M}\Delta_R+H_e(R),\qquad
 \Psi(R)=\sum_{n=1}^r\chi_n(R)\phi_n(R).
\]
Scalar nuclear repulsion can be included in $H_e$.  The corresponding
continuum-electronic formulas require common operator domains and the
regularity needed to justify these derivatives; the finite statement
does not assert those analytic hypotheses automatically.

\begin{proposition}[Exact coupled nuclear equations]
\label{prop:micro-bornhuang}
Set
\[
 \begin{aligned}
 d_{mn}&=\langle\phi_m,\nabla_R\phi_n\rangle,\\
 \tau_{mn}&=\langle\phi_m,\Delta_R\phi_n\rangle.
 \end{aligned}
\]
Then $H\Psi=E\Psi$ is equivalent to
\begin{equation}
 -\frac{\hbar^2}{2M}
  \left(\Delta\chi_m+
        2\sum_n d_{mn}\cdot\nabla\chi_n+
        \sum_n\tau_{mn}\chi_n\right)
  +E_m(R)\chi_m=E\chi_m
 \label{eq:micro-bornhuang}
\end{equation}
for each $m$.  Moreover, $d_{mn}=-\overline{d_{nm}}$ componentwise.
\end{proposition}
\begin{proof}
Apply the product rule to each
$\Delta(\chi_n\phi_n)$ and take its electronic inner product with
$\phi_m$.  Orthonormality gives the first term and diagonalizes
$H_e$.  Differentiating
$\langle\phi_m,\phi_n\rangle=\delta_{mn}$ proves the final identity.
\end{proof}

The usual single-surface approximation retains one electronic sector and
controls or omits its couplings to others.  Which diagonal geometric
terms are retained must be specified.  Rephasing
$\phi_n(R)\mapsto e^{i\gamma_n(R)}\phi_n(R)$ also rephases
$\chi_n$ and changes diagonal derivative couplings; a claim that they
vanish everywhere needs a gauge and topology argument.

\begin{proposition}[Energy gradients and the gap denominator]
\label{prop:micro-hellmann}
Let $E_n(R)$ be a simple normalized eigenpair of the finite matrix family.
For a coordinate derivative $\partial_j$,
\[
 \partial_j E_n
 =\langle\phi_n,(\partial_jH_e)\phi_n\rangle,
 \qquad
 \langle\phi_m,\partial_j\phi_n\rangle
 =\frac{\langle\phi_m,(\partial_jH_e)\phi_n\rangle}
        {E_n-E_m}\quad(m\ne n).
\]
If $|E_n-E_m|\ge\gamma>0$ for all $m\ne n$, then
\[
 \|(I-|\phi_n\rangle\langle\phi_n|)\partial_j\phi_n\|
 \le \gamma^{-1}\|\partial_jH_e\|.
\]
\end{proposition}
\begin{proof}
Differentiate $H_e\phi_n=E_n\phi_n$ and take inner products
with $\phi_n$ and with $\phi_m$.  The terms containing
$\partial_j\phi_n$ cancel in the first case.  In the second they
give the difference $E_n-E_m$.  Squaring the off-diagonal coefficients
and summing over the orthonormal eigenbasis gives the stated norm bound.
\end{proof}

Large nuclear mass reduces the prefactor in
\eqref{eq:micro-bornhuang}; it does not prevent a small electronic gap
from amplifying the derivative coupling.  Near crossings and conical
intersections, a single-surface argument requires additional work.
Likewise, an excited electronic surface may have an adiabatic nuclear
description.  The Born--Oppenheimer approximation is not a theorem that
electrons automatically dissipate into their ground state.

\begin{proposition}[Unitary evolution does not supply ground-state relaxation]
\label{prop:micro-no-relaxation}
For a time-independent finite Hermitian matrix $H$ and a spectral
projector $P$, the population
\[
 \Tr(Pe^{-itH/\hbar}\varrho_0e^{itH/\hbar})
\]
is constant.
\end{proposition}
\begin{proof}
The projector commutes with both exponentials.  Cyclicity of the trace
reduces the expression to $\Tr(P\varrho_0)$.
\end{proof}

A bath, emitted radiation, collisions with other degrees of freedom, or
a specified open-system approximation can generate relaxation.
Adiabatic following and relaxation are distinct physical mechanisms.
The video at 21:37--21:53 compresses them into a single explanation;
equations~\eqref{eq:micro-bornhuang} and
Proposition~\ref{prop:micro-no-relaxation} state what must be separated.

\section{When a quantum packet behaves approximately as a particle}

For
$H=P^2/(2M)+V(X)$, the formal Ehrenfest identities are
\begin{equation}
 \frac{d}{dt}\langle X\rangle=\frac{\langle P\rangle}{M},
 \qquad
 \frac{d}{dt}\langle P\rangle=-\langle\nabla V(X)\rangle.
 \label{eq:micro-ehrenfest}
\end{equation}
They follow from $d\langle A\rangle/dt=(i/\hbar)\langle[H,A]\rangle$
on a common invariant domain.  For instance, they hold for sufficiently
regular rapidly decaying solutions for which all integrations by parts
and differentiated moments are justified.  They are not yet Newton's
equations for the mean, because
$\langle\nabla V(X)\rangle$ need not equal
$\nabla V(\langle X\rangle)$.

\begin{proposition}[Localization controls the mean-force error]
\label{prop:micro-ehrenfest-error}
If $\nabla V$ is $L$-Lipschitz and the position distribution has a
finite second moment, then
\[
 \left|\langle\nabla V(X)\rangle-
       \nabla V(\langle X\rangle)\right|
 \le L\left\langle|X-\langle X\rangle|^2\right\rangle^{1/2}.
\]
\end{proposition}
\begin{proof}
Apply the Lipschitz inequality inside the position integral, then the
triangle inequality and Cauchy--Schwarz for its probability measure.
\end{proof}

This gives a quantitative condition rather than a universal particle-size
threshold.  To turn it into a trajectory approximation on a time interval,
one must control the packet width throughout that interval and estimate
the resulting perturbation of the classical flow.

For a free particle in one dimension with $\sigma>0$ and
\[
 \psi_0(x)=(2\pi\sigma^2)^{-1/4}
             e^{-x^2/(4\sigma^2)},
\]
the exact Heisenberg identity $X(t)=X+tP/M$ gives
\begin{equation}
 \operatorname{Var}(X(t))
 =\sigma^2+\frac{\hbar^2t^2}{4M^2\sigma^2}.
 \label{eq:micro-gaussian}
\end{equation}
Indeed $\operatorname{Var}X=\sigma^2$,
$\operatorname{Var}P=\hbar^2/(4\sigma^2)$, and the initial
symmetrized covariance vanishes.  Increasing $M$ slows free spreading
for a fixed packet, but \eqref{eq:micro-gaussian} still grows with time.
Heavy particles can also exhibit interference.  A suitable scale
comparison involves action, wavelength, temperature, environmental
coupling, and the observable being approximated.

\subsection{One-dimensional Coulomb models require a domain}

The video explicitly uses a softened one-dimensional potential to explain
electronic forces.  Such a model is pedagogically useful, but it is a
different operator from the singular expression
\[
 -\frac{\hbar^2}{2m}\frac{d^2}{dx^2}-\frac{\kappa}{|x|}.
\]
A differential expression does not specify a self-adjoint Hamiltonian
until its domain and behavior at the singular point are given.
The reviewed operator-domain analysis by de Oliveira and Verri finds a
family of self-adjoint realizations for the one-dimensional operator
initially defined away from the origin
\citep{microOliveiraVerri2009}.  We use that domain distinction, not
an unqualified assertion about the spectrum of every regularization.
The requested Loudon papers are separately registered as unavailable
full texts in this review \citep{microLoudon1959,microLoudon2016}.

The limit of a family of softened potentials must specify the topology
of operator convergence and any renormalization of energy.  It cannot
be inferred merely by drawing potentials that approach $-\kappa/|x|$
away from zero.  This is the same structural issue that recurs in
continuum limits of graph Laplacians: a formula for coefficients must be
accompanied by domains, boundary conditions, and convergence estimates.

\section{Molecular dynamics and conserved microscopic information}

A classical molecular Hamiltonian has the form
\begin{equation}
 H(q,p)=\sum_{i=1}^N\frac{|p_i|^2}{2m_i}+U(q),\qquad
 \dot q_i=\frac{p_i}{m_i},\quad
 \dot p_i=-\nabla_{q_i}U.
 \label{eq:micro-md}
\end{equation}
The force is the negative gradient of potential energy.  General
molecular force fields include angular, torsional, many-body, constrained,
or reactive terms; they need not be sums of pairwise distance functions.
Allen's notes explain these modeling choices, symplectic integrators,
constraints, thermostats, and statistical errors
\citep{microAllen2004}.  Their scope is much wider than the two-particle
illustration in the video.

\begin{proposition}[Liouville volume and fine-grained entropy]
\label{prop:micro-liouville}
Let $H\in C^2(\R^{2n})$ have a smooth Hamiltonian flow $\Phi_t$
on the interval considered.  Then $\det D\Phi_t=1$.
If a smooth probability density evolves by
$f_t=f_0\circ\Phi_{-t}$ and the entropy integral is absolutely
integrable, then
\[
 \int f_t\log f_t\,dq\,dp
 =\int f_0\log f_0\,dq\,dp.
\]
\end{proposition}
\begin{proof}
For the Hamiltonian vector field
$X_H=(\partial_pH,-\partial_qH)$, equality of mixed derivatives
gives
$\operatorname{div}X_H=\sum_i(\partial_{q_i}\partial_{p_i}H-
\partial_{p_i}\partial_{q_i}H)=0$.
Jacobi's determinant identity for the flow derivative yields
$\frac d{dt}\det D\Phi_t=(\operatorname{div}X_H)\circ\Phi_t
\det D\Phi_t=0$.  Its initial value is one.  A change of variables
in the entropy integral proves the second assertion.
\end{proof}

Thus a Hamiltonian description preserves fine-grained information.
Entropy growth in a reduced description is not the destruction of
information by an invertible microscopic flow.  It involves a
marginalization, an environment, a limiting operation, or a restriction
on what is resolved.  Quantum unitary conjugation similarly preserves
the spectrum of a density operator and hence its von Neumann entropy
when that entropy is defined.  A reduced subsystem can have changing
entropy because its correlations with the omitted subsystem change.

\subsection{A useful exact statement about velocity Verlet}

For separable $H=T(p)+U(q)$, the kick
$(q,p)\mapsto(q,p-h\nabla U(q))$ and drift
$(q,p)\mapsto(q+h\nabla T(p),p)$ are exact Hamiltonian subflows.
Their symmetric composition, kick by $h/2$, drift by $h$, kick by
$h/2$, is symplectic.  This proves preservation of the symplectic form;
it does not prove exact preservation of the original $H$.

For the scalar oscillator $H=(p^2+\omega^2q^2)/2$, the update is
\[
 q'=\left(1-\frac{h^2\omega^2}{2}\right)q+hp,
 \quad
 p'=\left(1-\frac{h^2\omega^2}{2}\right)p
       -h\omega^2\left(1-\frac{h^2\omega^2}{4}\right)q.
\]
Direct expansion proves that it preserves
\begin{equation}
 \widetilde H_h(q,p)
 =\frac12p^2+\frac12\omega^2
       \left(1-\frac{h^2\omega^2}{4}\right)q^2.
 \label{eq:micro-verlet-invariant}
\end{equation}
This quadratic form is positive definite when $0<h\omega<2$.
It is a valid special-case modified invariant.  In a general nonlinear
system, a formal or asymptotic modified Hamiltonian from backward error
analysis must not be promoted to an exact global invariant without the
additional theorem needed for that claim.

\section{Marginals, correlations, and the kinetic approximation}

For a smooth symmetric $N$-particle density $F_N$, let $F_N^{(s)}$
denote its $s$-particle marginal.  With unit masses and a pair force
$K(x)=-\nabla\Phi(x)$, integrating the Liouville equation over all
but particle 1 gives
\begin{equation}
 \partial_tF_N^{(1)}+v\cdot\nabla_xF_N^{(1)}
 +(N-1)\nabla_v\cdot
   \int K(x-y)F_N^{(2)}(x,v,y,w)\,dy\,dw=0,
 \label{eq:micro-bbgky-first}
\end{equation}
assuming the boundary terms vanish.  The two-particle marginal appears
because the force on a particle depends on others.  Its own equation
involves the three-particle marginal.  This is the BBGKY hierarchy.

\begin{example}[A first marginal cannot determine a pair correlation]
\label{ex:micro-marginal}
On $\{0,1\}^2$, let $P$ give mass $1/2$ each to $(0,0)$ and
$(1,1)$, and let $Q$ give mass $1/2$ each to $(0,1)$ and $(1,0)$.
Both first and second one-coordinate marginals are uniform.  Nevertheless,
$P(X=Y)=1$ and $Q(X=Y)=0$.  Thus no function of the one-coordinate
marginals alone can reconstruct all pair observables.
\end{example}

Replacing $F_N^{(2)}$ by a product is a substantive assumption.
For dilute collisions it is the incoming, appropriately scaled
correlations that matter.  Hard-sphere exclusion already prevents exact
independence at fixed particle diameter.  The word ``chaos'' in
propagation of chaos denotes an asymptotic factorization property of
fixed-order marginals; it is not interchangeable with deterministic
sensitivity to initial conditions.

\subsection{A rigorous classical-to-kinetic bridge and its time interval}

In the hard-sphere Boltzmann--Grad scaling in dimension $d\ge2$,
\[
 N\varepsilon^{d-1}=1,
 \qquad N\varepsilon^d=\varepsilon\longrightarrow0.
\]
The first relation retains an order-one collision scale, while the
second sends the occupied volume fraction to zero.  Gallagher,
Saint-Raymond, and Texier give a detailed derivation for hard spheres
and suitable short-range repulsive potentials
\citep{microGallagher2013}.  Their hard-sphere theorem propagates
admissible asymptotic initial factorization on a positive short time
interval, with convergence of velocity-averaged observables locally
away from positional diagonals.  It uses weighted uniform bounds on
the hierarchy and quantitative control of recollisions.

The reviewed source's introductory statement assumes a continuous
probability density $f_0$ with a Gaussian velocity bound
\[
 \sup_{x,v}f_0(x,v)e^{\beta|v|^2/2}<\infty
 \quad\text{for some }\beta>0.
\]
Its precise Theorem 8 includes the admissibility and hierarchy-space
hypotheses.  The proof does not assert convergence for arbitrary
finite-$N$ initial data, arbitrarily long times, every potential, or
every boundary geometry.  This chapter reviewed the introduction,
main convergence statement, and concluding scope discussion; it does
not claim to reproduce or independently verify the monograph's complete
recollision proof.

This time restriction has a practical logical consequence.  A short-time
particle-to-Boltzmann theorem and a long macroscopic-time
Boltzmann-to-fluid theorem cannot simply be concatenated.  The former
must hold uniformly on the time scales required by the latter, or a
separate joint-limit theorem must be supplied.  No amount of formal
diagram composition replaces that analytic requirement.

\section{The collision algebra and its five invariants}

We now prove the algebraic facts behind the simplest elastic collision
operator.  They are useful both for conservative simulation and for
checking learned collision models.  Set $g=v-v_*$, let
$n\in S^{d-1}$, and define
\begin{equation}
 v'=v-(g\cdot n)n,\qquad
 v_*'=v_*+(g\cdot n)n.
 \label{eq:micro-collision}
\end{equation}

\begin{proposition}[Elastic reflection]
\label{prop:micro-collision}
The map \eqref{eq:micro-collision} is an involution with absolute
Jacobian determinant one.  It preserves total momentum and kinetic
energy:
\[
 v'+v_*'=v+v_*,\qquad
 |v'|^2+|v_*'|^2=|v|^2+|v_*|^2.
\]
\end{proposition}
\begin{proof}
The center velocity $V=(v+v_*)/2$ is unchanged and
$g'=g-2(g\cdot n)n$.  This is an orthogonal reflection, so
$|g'|=|g|$, its square is the identity, and its determinant has
absolute value one.  Since
$|v|^2+|v_*|^2=2|V|^2+|g|^2/2$, all assertions follow after the
invertible linear change from $(v,v_*)$ to $(V,g)$.
\end{proof}

A scalar function $\phi$ is a collision invariant when
$\phi(v)+\phi(v_*)=\phi(v')+\phi(v_*')$ for every elastic
collision.  The functions $1,v_1,\ldots,v_d,|v|^2$ are immediately
invariant.  The following proof makes the converse precise under a
simple regularity assumption.

\begin{proposition}[Classification of smooth collision invariants]
\label{prop:micro-invariants}
For $d\ge2$, every $C^2$ collision invariant has the form
\[
 \phi(v)=a+b\cdot v+c|v|^2.
\]
Consequently the scalar invariant space has dimension $d+2$, hence
five in three dimensions.
\end{proposition}
\begin{proof}
For fixed $V$ and radius $r$, collisions can send $g$ to any vector
of the same length: reflection in a suitable hyperplane does so for
two given vectors.  Hence
$\phi(V+r\omega)+\phi(V-r\omega)$ is independent of the unit
vector $\omega$.  Its second-order expansion at $r=0$ implies
that $\omega^TD^2\phi(V)\omega$ is constant on the unit sphere.
Taking coordinate vectors and their pairwise sums shows
$D^2\phi(V)=\lambda(V)I$.

The mixed second derivatives are zero, so $\partial_i\phi(v)$
depends only on $v_i$.  The equality of all diagonal second
derivatives then says that a function of $v_i$ alone equals a function
of $v_j$ alone for independently variable $v_i,v_j$.  Since $d\ge2$,
each is the same constant.  Integrating twice yields the displayed
quadratic polynomial.  The converse follows from
Proposition~\ref{prop:micro-collision}.
\end{proof}

Some introductory notes group the invariants into three physical
categories: particle number, momentum, and energy.  Momentum is a vector,
so those categories contain five independent scalar invariants in
three dimensions, not three scalar equations.

\subsection{A collision kernel contains more information than a diameter}

For a central classical interaction, the differential cross section is
obtained from the map between impact parameter $b$ and scattering
angle $\chi$.  On a branch where the map is invertible,
\[
 \frac{d\sigma}{d\Omega}
 =\frac{b}{\sin\chi}\left|\frac{db}{d\chi}\right|.
\]
If multiple branches produce the same angle, their contributions must
be added.  Angular singularities and grazing collisions require further
analysis.  The hard-sphere model and the variable hard-sphere model
used in Volkov's notes do not reproduce every angular distribution of
a general potential \citep{microVolkov02}.  Matching a viscosity
law or a mixing statistic is a calibration of particular observables,
not an equality of the underlying scattering laws.

Internal excitation and chemical reactions add variables, transitions,
and conservation relations.  Parker's requested paper concerns
rotational and vibrational relaxation, while Valentini and collaborators'
requested study concerns nitrogen dissociation on an ab initio potential
surface \citep{microParker1959,microValentini2015}.  Only their
accessible abstracts and metadata were reviewed here.  Their reported
models must not be treated as fully audited derivations or as evidence
that every gas can be represented by an elastic monatomic kernel.

\section{The weak collision identity and entropy production}

Let $B(g,n)\ge0$ be a collision kernel invariant under exchange of the
particles and under the pre/post-collision reflection.  Define
\begin{equation}
 Q(f,f)(v)=\int_{\R^d}\int_{S^{d-1}}
       B(v-v_*,n)(f'f_*'-ff_*)\,dn\,dv_*.
 \label{eq:micro-boltzmann-Q}
\end{equation}
Here all four values of $f$ have the same $(t,x)$, when those
variables are present.  The next statements assume absolute
integrability of the displayed integrals, so that Fubini's theorem
and the changes of variables are legitimate.

\begin{proposition}[Weak collision identity]
\label{prop:micro-weak-Q}
For an admissible test function $\phi$,
\begin{align}
 \int\phi Q(f,f)\,dv
 &=\frac14\iiint B(f'f_*'-ff_*)
       (\phi+\phi_* -\phi'-\phi_*')\,dn\,dv_*\,dv
       \label{eq:micro-weak-Q}\\
 &=\frac12\iiint Bff_*
       (\phi'+\phi_*'-\phi-\phi_*)\,dn\,dv_*\,dv.
       \nonumber
\end{align}
\end{proposition}
\begin{proof}
Exchange $v$ and $v_*$ in the original integral and average the
two expressions.  In a second copy, change variables from
$(v,v_*)$ to $(v',v_*')$.  The measure and $B$ are unchanged,
the gain-minus-loss factor changes sign, and the test functions
acquire primes.  Averaging once more gives the first identity.
Applying the same pre/post change to the gain term gives the second.
\end{proof}

Testing with collision invariants proves the mass, momentum, and energy
collision balances.  Testing with the logarithm supplies a stronger
constraint, subject to the additional integrability needed for that
unbounded test.

\begin{theorem}[Classical homogeneous $H$ theorem under explicit regularity]
\label{thm:micro-H}
Let $f(t,v)>0$ be a differentiable solution of
$\partial_tf=Q(f,f)$ for which $\int f\log f$ is finite,
differentiation under the integral is valid, and all terms in
\eqref{eq:micro-weak-Q} with $\phi=\log f$ are integrable.
Then
\begin{equation}
 \frac{d}{dt}\int f\log f\,dv=-D(f)\le0,
 \quad
 D(f)=\frac14\iiint B(f'f_*'-ff_*)
           \log\frac{f'f_*'}{ff_*}\,dn\,dv_*\,dv.
 \label{eq:micro-H}
\end{equation}
\end{theorem}
\begin{proof}
The derivative equals $\int(1+\log f)Q(f,f)\,dv$.
The constant part vanishes by mass conservation.  Use
\eqref{eq:micro-weak-Q} for the remaining part.  Since
$(a-b)(\log a-\log b)\ge0$ for $a,b>0$, the displayed
entropy production is nonnegative.
\end{proof}

For spatially varying smooth solutions, transport adds the entropy flux
$\nabla_x\cdot\int vf\log f\,dv$.  Boundary entropy fluxes must
then be included.  For weak or renormalized solutions, an entropy
inequality is established by approximation and lower semicontinuity;
one cannot simply insert an unjustified logarithmic test into a weak
equation.

\begin{proposition}[Finite energy does not imply finite positive entropy]
\label{prop:micro-entropy-counterexample}
There is a probability density supported in $[0,1]^3$ with finite
second moment and $\int f\log f=+\infty$.
\end{proposition}
\begin{proof}
Let $c=6/\pi^2$ and choose disjoint measurable sets $A_n$ in
$[0,1]^3$ with volumes
$|A_n|=c e^{-n^2}/n^2$.  Such sets exist because the sum of
their volumes is less than one.  Define $f=e^{n^2}$ on $A_n$
and $f=0$ elsewhere, with $0\log0=0$.  Then
\[
 \int f=c\sum_{n\ge1}n^{-2}=1,
 \qquad \int|v|^2f\,dv\le3,
 \qquad \int f\log f=c\sum_{n\ge1}1=+\infty.
\]
\end{proof}

This corrects a gap in the elementary entropy discussion in the
reviewed Volkov slides \citep{microVolkov03}.  Similarly, integrability
of a continuous density alone does not force pointwise decay to zero
at infinity: separated continuous spikes of height one and summable
widths provide a counterexample.  Decay assumptions used in
integration by parts must therefore be stated, not inferred solely
from normalization.

\subsection{Entropy minimization is not a convergence proof}

\begin{proposition}[Maxwellian minimizes $H$ at fixed first two moments]
\label{prop:micro-maxwell}
Let $f$ be a probability density with mean $u$ and
$\int|v-u|^2f\,dv=d\theta$, where $\theta>0$.  Set
\[
 M_{u,\theta}(v)=(2\pi\theta)^{-d/2}
       e^{-|v-u|^2/(2\theta)}.
\]
Whenever the entropy expressions are defined,
\[
 \int f\log f-\int M_{u,\theta}\log M_{u,\theta}
 =\int f\log\frac{f}{M_{u,\theta}}\ge0.
\]
Equality holds exactly when $f=M_{u,\theta}$ almost everywhere.
\end{proposition}
\begin{proof}
The function $\log M_{u,\theta}$ is a quadratic polynomial in
$v$, so its integral against $f$ agrees with its integral against
$M_{u,\theta}$.  This proves the identity.  With
$r=f/M_{u,\theta}$, use $r\log r-r+1\ge0$ and integrate
against $M_{u,\theta}\,dv$.  Equality in this strictly convex
inequality requires $r=1$ almost everywhere.
\end{proof}

Under a positive kernel, sufficient continuity, and $\log f\in C^2$,
vanishing $D(f)$ forces detailed balance
$f'f_*'=ff_*$.  Proposition~\ref{prop:micro-invariants} applied to
$\log f$ then gives the Maxwellian form, with a negative quadratic
coefficient required by integrability.

These facts characterize entropy minimizers and stationary collision
equilibria.  They do not prove that every trajectory converges to one.
If a scalar function $H(t)$ is nonincreasing and bounded below, it
has a limit, but its derivative need not tend to zero.  For example,
take a nonnegative smooth function $b$ supported in $(-1/4,1/4)$
with $b(0)=1$ and define
\[
 h(t)=\sum_{n\ge1}b(2^n(t-n)),\qquad
 H(t)=\int_t^\infty h(s)\,ds.
\]
The bumps have disjoint supports and summable areas.  Hence $H$ is
smooth, nonnegative, and decreasing to zero, while $H'(n)=-1$.
Even $H'(t)\to0$ alone would not supply compactness of the state
or identify a unique limit.  Convergence-to-equilibrium theorems require
the appropriate coercivity, compactness, integrability, and control of
invariants in the stated regime.

\section{Exact moment balances and the closure problem}

For a smooth unit-mass solution of
\[
 \partial_tf+v\cdot\nabla_xf=Q(f,f),
\]
assume sufficient velocity integrability to justify the following
moments.  Let $c=v-u$ and define
\[
 P=\int c\otimes c\,f\,dv,\quad
 e=\frac1{2\rho}\int|c|^2f\,dv,\quad
 q=\frac12\int|c|^2c f\,dv,
 \quad E=\frac12\rho|u|^2+\rho e.
\]
Here $P$ is a positive semidefinite pressure tensor.  The Cauchy stress
uses the opposite sign for the kinetic contribution; keeping this
convention explicit avoids sign errors in viscosity formulas.

\begin{proposition}[Moment equations before closure]
\label{prop:micro-moments}
Under these hypotheses,
\begin{align}
 \partial_t\rho+\nabla\cdot(\rho u)&=0,
 \label{eq:micro-mass}\\
 \partial_t(\rho u)+\nabla\cdot(\rho u\otimes u+P)&=0,
 \label{eq:micro-momentum}\\
 \partial_tE+\nabla\cdot(Eu+Pu+q)&=0.
 \label{eq:micro-energy}
\end{align}
\end{proposition}
\begin{proof}
Test the kinetic equation with $1,v,|v|^2/2$ and use the collision
invariants.  Since $\int cf\,dv=0$, expansion of
$v=u+c$ gives
$\int v\otimes vf\,dv=\rho u\otimes u+P$.
Likewise,
\[
 \frac12\int |v|^2vf\,dv
 =\left(\frac12\rho|u|^2+\rho e\right)u+Pu+q.
\]
These expansions give all three equations.
\end{proof}

\begin{example}[Equal five moments, different stress]
\label{ex:micro-unclosed}
Let $f_A$ and $f_B$ be centered normalized Gaussian densities on
$\R^3$ with covariance matrices
\[
 A=I,\qquad B=\operatorname{diag}(1+\delta,1-\delta,1),
 \quad0<\delta<1.
\]
They have the same mass, mean, and scalar temperature because
$\tr A=\tr B=3$.  Their pressure tensors are $A$ and $B$,
respectively.  Thus the five fluid moments do not determine the
momentum flux.
\end{example}

A local Maxwellian imposes $P=\rho\theta I$ and $q=0$, giving
the Euler fluxes.  An arbitrary local Maxwellian with space- and
time-dependent parameters is generally not a solution of the full
Boltzmann equation: its collision term vanishes, but its transport
term need not.  Euler equations arise as moment equations under
that closure or as a justified limit, rather than as a proof that
every Euler solution lifts to an exact local-Maxwellian kinetic
solution.  This qualification is stated explicitly in the rigorous
hydrodynamic-limit literature \citep{microGolseSaintRaymond2009}.

\subsection{First-order constitutive laws and scaling}

For a monatomic gas near local equilibrium, a formal
Chapman--Enskog expansion produces the constitutive form
\begin{equation}
 P=pI-2\mu S_0-\zeta(\nabla\cdot u)I,
 \qquad q=-\kappa\nabla T,
 \quad
 S_0=\tfrac12(\nabla u+\nabla u^T)
          -\tfrac13(\nabla\cdot u)I.
 \label{eq:micro-constitutive}
\end{equation}
The bulk coefficient is zero in the simplest monatomic elastic
model.  Internal modes and other models change this conclusion.
The values of the transport coefficients require solving equations
for the linearized collision operator; conservation and symmetry
alone do not fix them.

The Knudsen number is the ratio of a microscopic mean free path to
the macroscopic length scale.  Small Knudsen number, small Mach number,
and fixed Reynolds number specify distinct relationships among
scales.  A formal expansion in one does not silently impose the
others.  A global small value also does not rule out local breakdown
near a shock, boundary layer, or strongly varying rarefied region.
Volkov's first and third decks make the physical scale separation
and the failure of elementary closure in rarefied regimes explicit
\citep{microVolkov01,microVolkov03}.

\subsection{A precise historical hydrodynamic-limit theorem}

The result of Golse and Saint-Raymond that is useful here concerns
the scaled equation on $\R^3_x\times\R^3_v$
\begin{equation}
 \varepsilon^2\partial_tF_\varepsilon
 +\varepsilon v\cdot\nabla_xF_\varepsilon
 =Q(F_\varepsilon,F_\varepsilon),
 \qquad M(v)=(2\pi)^{-3/2}e^{-|v|^2/2}.
 \label{eq:micro-hydro-scale}
\end{equation}
Their collision kernels obey the exchange and reflection symmetries
and hard-cutoff bounds of the form
\[
 0<B(z,n)\le C(1+|z|)^\beta|\widehat z\cdot n|,
 \qquad
 \int_{S^2}B(z,n)\,dn\ge C^{-1}\frac{|z|}{1+|z|},
 \quad0\le\beta\le1,
\]
with the inequalities understood almost everywhere.  Initial data
satisfy
\begin{equation}
 \sup_{\varepsilon>0}\varepsilon^{-2}
 H(F^{\mathrm{in}}_\varepsilon\mid M)<\infty,
 \quad
 H(F\mid M)=\iint\left(F\log\frac FM-F+M\right)\,dx\,dv.
 \label{eq:micro-relative-scaling}
\end{equation}
Their Theorem 2.4 describes limit points of the fluctuations
$g_\varepsilon=(F_\varepsilon/M-1)/\varepsilon$ for the
corresponding renormalized solutions.  In the stated weak local
topology the limit has the form
\[
 g(t,x,v)=u(t,x)\cdot v+
            \tfrac12\theta(t,x)(|v|^2-5),
\]
where $(u,\theta)$ solves an incompressible
Navier--Stokes--Fourier system and satisfies its energy inequality
\citep[Theorem 2.4]{microGolseSaintRaymond2009}.
The initial velocity is determined by the weak limit of the Leray
projection of the scaled momentum; a corresponding thermal moment
determines the initial temperature fluctuation.  For appropriate
well-prepared initial entropy, the velocity is a Leray solution
with the corresponding initial energy bound.

This is a theorem about a specified collision class, spatial domain,
scaling, solution concept, and convergence topology.  Its proof uses
entropy bounds, coercivity of the linearized collision operator,
velocity averaging, vanishing conservation defects, and control of
nonlinear fluxes.  It yields weak fluid solutions; it does not by
itself upgrade them to every stronger regularity or uniqueness
statement.  Statements about the separate 2026 proof release are
analyzed in Chapter~\ref{ch:navier-stokes-release}.  That chapter's
precise theorem must be used with its own forcing and solution
hypotheses, rather than inferred from this historical limit theorem.

\section{DSMC as a stochastic approximation, with auditable assumptions}

The requested sixth Volkov deck explains direct simulation Monte Carlo
as a particle approximation to kinetic dynamics, with statistical
weights, free-flight and collision substeps, cell-based collision
selection, inflow and wall kernels, and sampling of moments
\citep{microVolkov06}.  These simulated particles represent weighted
samples.  Their paths are not claimed to be the paths of every physical
molecule in the gas.

Three error categories must be distinguished: the physical model
(collision law and wall interaction), numerical discretization
(time step, cells, particle count, cutoffs), and statistical estimation.
Increasing a GPU sample count primarily addresses the latter two
categories.  It cannot validate a physical collision law by itself.

\begin{proposition}[Flux sampling at an equilibrium wall]
\label{prop:micro-flux}
Let a stationary wall emit an equilibrium gas with thermal variance
$\theta>0$ and inward unit normal $n$.  The normalized velocity
density of molecules crossing the wall into the domain is
\[
 p_{\mathrm{flux}}(v)
 =\frac{(v\cdot n)_+}{2\pi\theta^2}
       e^{-|v|^2/(2\theta)}.
\]
Its positive normal component $s$ has density
$(s/\theta)e^{-s^2/(2\theta)}$, and its two tangential
components are independent centered Gaussians of variance $\theta$.
\end{proposition}
\begin{proof}
A bulk density $M_\theta(v)$ contributes a crossing flux
$(v\cdot n)_+M_\theta(v)$.  The integral of this flux is
$\sqrt{\theta/(2\pi)}$.  Dividing by that value gives the
displayed density.  In orthonormal normal/tangential coordinates
the exponential and its prefactor factor into the three stated
one-dimensional densities.
\end{proof}

Sampling the normal velocity from an unweighted half-Gaussian would
therefore give the wrong boundary ensemble.  A finite reservoir
implementation also requires an estimate of the missing high-speed
tail; the Maxwellian has no finite maximum speed.

\begin{proposition}[The exact frozen-rate majorant identity]
\label{prop:micro-majorant}
Consider $K\ge1$ possible events with fixed rates $r_j\ge0$, and
choose $R>0$ with $R\ge\max_jr_j$.  Generate a Poisson proposal process of rate
$KR$, choose a label uniformly from $\{1,\ldots,K\}$ at each
proposal, and accept label $j$ with probability $r_j/R$.
The accepted events of label $j$ form a Poisson process of rate
$r_j$; distinct accepted labels have independent counts.
\end{proposition}
\begin{proof}
The probability that a proposal is accepted with label $j$ is
$r_j/(KR)$.  For an interval of length $t$, condition on its
Poisson number of proposals and take the joint probability
generating function of the marked counts.  Summing the Poisson
series gives
\[
 \exp\!\left(t\sum_{j=1}^K r_j(z_j-1)\right)
 =\prod_{j=1}^K\exp(tr_j(z_j-1)).
\]
This is the product of the generating functions of independent
Poisson variables with the claimed means.  The same argument on
disjoint time intervals proves the process statement.
\end{proof}

This proves the sampling principle behind majorant collision schemes.
It does not alone prove convergence of a full DSMC method whose rates
change after collisions and whose transport is split in time.  In that
case the bound must remain valid for the current state and the
discretization error must be controlled.  If an alleged acceptance
probability exceeds one, silently clipping it changes the rate.
A temperature-based constant in illustrative code is not a proof of
a global bound on an unbounded velocity distribution.

\begin{proposition}[Sampling variance with correlations]
\label{prop:micro-sampling}
If $Y_1,\ldots,Y_n$ have finite second moments, then
\[
 \operatorname{Var}\!\left(\frac1n\sum_iY_i\right)
 =\frac1{n^2}\sum_{i,j}\operatorname{Cov}(Y_i,Y_j).
\]
For identically distributed independent samples this equals
$\sigma^2/n$.  For a stationary sequence with variance
$\sigma^2>0$ and lag correlations $r_k$, it equals
\[
 \frac{\sigma^2}{n}
 \left(1+2\sum_{k=1}^{n-1}(1-k/n)r_k\right).
\]
\end{proposition}
\begin{proof}
Expand the square of the centered sum and take expectations.
For a stationary sequence, group off-diagonal pairs by their
lag; there are $n-k$ pairs at lag $k$ in each direction.
\end{proof}

Thinning a time series can reduce some correlations but does not prove
independence.  A central-limit error bar requires its own assumptions;
a fixed seed provides reproducibility, not randomness certification.
For independent samples, Chebyshev's inequality gives the explicit
bound
$\Pr(|\overline Y-\mathbb EY|\ge a)\le\sigma^2/(na^2)$.
The statistics and Monte Carlo decks are useful introductions
\citep{microVolkov04,microVolkov05}; implementations should additionally
handle finite precision, zero arguments of logarithms, array bounds,
correct random-number-generator state, and reproducible restart state.
The period of a generator is constrained by its internal state space,
not merely by the number of possible values of a single output word.

\section{Reduced models and learned closures}

An orthogonal reduced basis minimizes a chosen projection loss on its
training snapshots.  That statement does not imply that the projected
trajectory is stable, that the discarded modes exert no force, or that
the basis will predict new parameter regimes.  The currently available
replacement Farhat course decks explicitly distinguish projection
error from model-dynamics error \citep{microFarhatCourse}.  Their
complete version-specific review is recorded separately in the source
audit.

For an orthonormal column matrix $V$, write a full state
$z=Va+w$ with $V^Tw=0$.  If $\dot z=F(z)$, the exact retained
coordinate satisfies $\dot a=V^TF(Va+w)$.  Replacing the right
side by $V^TF(Va)$ is a closure.  The exact memory identity in
Chapter~\ref{ch:quantum-fluid} explains why eliminated coordinates
can leave both memory and an initial-condition contribution.

\begin{proposition}[What weighted POD actually minimizes]
\label{prop:micro-pod}
Let $X=[x_1\ \cdots\ x_s]\in\R^{n\times s}$, let
$W=\operatorname{diag}(w_1,\ldots,w_s)$ with $w_j>0$, and
let $\lambda_1\ge\cdots\ge\lambda_n\ge0$ be the eigenvalues
of $C=XWX^T$.  Fix $1\le r\le n$.  Among all
$V\in\R^{n\times r}$ with $V^TV=I$,
\[
 \min_V\sum_jw_j\|x_j-VV^Tx_j\|^2
 =\sum_{i>r}\lambda_i.
\]
A minimizing range is spanned by eigenvectors for the largest
$r$ eigenvalues.
\end{proposition}
\begin{proof}
The objective is $\tr C-\tr(V^TCV)$.  In an orthonormal
eigenbasis $e_i$ of $C$, put $a_i=\|V^Te_i\|^2$.
Then $0\le a_i\le1$, $\sum_i a_i=r$, and
$\tr(V^TCV)=\sum_i\lambda_i a_i$.  Transferring any weight
from a smaller eigenvalue to an unfilled larger one cannot decrease
this sum.  Its maximum is therefore $\sum_{i\le r}\lambda_i$,
and the indicated eigenspace attains it.
\end{proof}

The proposition concerns the chosen snapshots and weights.  It contains
no assertion about an unsampled parameter or a differential equation.
Indeed the stable matrix
\[
 A=\begin{pmatrix}-1&10\\0&-1\end{pmatrix},\qquad
 V=2^{-1/2}\begin{pmatrix}1\\1\end{pmatrix}
\]
has both eigenvalues equal to $-1$, while its orthogonal Galerkin
reduction is $V^TAV=4$.  The reduced scalar solution grows
exponentially.  Nonsymmetric stable dynamics are therefore not
automatically preserved by orthogonal projection.

For Petrov--Galerkin projection with full-column-rank trial and test
matrices $V,W$, the matrix $W^TV$ must be invertible to use
$\Pi=V(W^TV)^{-1}W^T$.  Direct multiplication gives
$\Pi^2=\Pi$, $\operatorname{range}\Pi=\operatorname{range}V$,
and $\ker\Pi=\ker W^T=(\operatorname{range}W)^\perp$.
The direction of this oblique projection is the latter kernel.

\subsection{Hyperreduction adds another controlled approximation}

State reduction need not make a nonlinear residual cheap to evaluate.
Let $U\in\R^{n\times r}$ have orthonormal columns and let
$P\in\R^{n\times r}$ select $r$ coordinates.  If
$P^TU$ is invertible, the discrete empirical interpolation map is
\[
 D=U(P^TU)^{-1}P^T.
\]

\begin{proposition}[Interpolation exactness and amplification]
\label{prop:micro-deim}
For every $y\in\R^n$, $DU=U$, $P^TDy=P^Ty$, and
\[
 \|y-Dy\|
 \le\left(1+\|(P^TU)^{-1}\|\right)
       \|(I-UU^T)y\|.
\]
\end{proposition}
\begin{proof}
The first two identities follow by cancellation of $P^TU$.
Because $D$ is the identity on the range of $U$,
$(I-D)y=(I-D)(I-UU^T)y$.
Use $\|I-D\|\le1+\|D\|$ and the unit operator norms
of $U$ and $P^T$ to obtain the estimate.
\end{proof}

The bound is deliberately elementary; sharper oblique-projection
bounds are possible.  Even this version shows why a small discarded
snapshot norm is insufficient if the sampled basis is ill-conditioned.
For oversampled least squares, full column rank of $P^TU$ is the
corresponding uniqueness and exact-reproduction requirement.
An online cost independent of $n$ additionally requires that sampled
nonlinear entries and the restricted reconstruction can be evaluated
without assembling all $n$ entries.  A dense nonlocal nonlinear map
can violate that requirement.

Other reduction methods have different hypotheses.  Classical balanced
truncation concerns stable finite-dimensional linear time-invariant
systems, with reachability and observability conditions, and its usual
input--output error estimate assumes zero initial state.  An empirical
balanced-POD approximation does not inherit every exact balancing
guarantee.  Moment matching requires the relevant full and reduced
resolvents to exist at its interpolation nodes.  Grassmann-chart
interpolation requires the chart to be valid, fixed subspace rank,
and appropriate alignment before interpolating reduced operators.
The detailed Farhat review documents these distinctions
\citep{microFarhatCourse}.

Likewise, a weighted finite-element energy
$\widetilde U(q)=\sum_e\xi_e U_e(q)$ with fixed weights has
gradient $\sum_e\xi_e\nabla U_e(q)$.  If a neural model makes
$\xi_e=\xi_e(q)$, the gradient instead includes
$\sum_e U_e(q)\nabla\xi_e(q)$.  Omitting those terms loses the
claimed potential-force relation.  A preservation theorem for a
fixed-weight mechanical hyperreduction must therefore be re-proved
when the weights become state dependent or the equation changes to
a different fluid formulation.

\subsection{A provable constraint for a neural constitutive model}

A classical or quantum neural model can be used to predict constitutive
coefficients from invariant local features.  For example, a model can
output nonnegative scalars $\mu,\zeta,\kappa$ in
\eqref{eq:micro-constitutive}.  This architecture has a limited but
useful exact guarantee.

\begin{proposition}[Nonnegative constitutive entropy production]
\label{prop:micro-learned-entropy}
For $T>0$, let the viscous stress be
$\tau=2\mu S_0+\zeta(\nabla\cdot u)I$ and the heat flux
$q=-\kappa\nabla T$, with $\mu,\zeta,\kappa\ge0$.
Then
\[
 \frac{\tau:\nabla u}{T}
       -\frac{q\cdot\nabla T}{T^2}
 =\frac{2\mu|S_0|^2}{T}
  +\frac{\zeta(\nabla\cdot u)^2}{T}
  +\frac{\kappa|\nabla T|^2}{T^2}\ge0.
\]
\end{proposition}
\begin{proof}
The symmetric stress is orthogonal in the Frobenius inner product to
the antisymmetric part of $\nabla u$.  The traceless symmetric
part $S_0$ is orthogonal to $I$.  Decompose the symmetric part of
$\nabla u$ into $S_0+(\nabla\cdot u)I/3$ and substitute the
heat-flux formula.
\end{proof}

If the scalar inputs are rotational invariants, these constitutive
tensors transform correctly under rotations.  Dependence on gradients
and thermodynamic invariants also permits Galilean covariance.
Neither property proves accurate transport coefficients.  Moreover,
nonnegative entropy production does not by itself make
$P=pI-\tau$ positive semidefinite far from equilibrium.  A kinetic
moment model that requires realizable covariances can instead
parameterize $P=LL^T$ or another positive-definite form and impose
the remaining energy and flux constraints explicitly.  These are
different requirements and must be checked together.

Tensor decomposition makes these constraints transparent.  In three
dimensions a rank-two tensor splits into scalar trace, antisymmetric,
and symmetric traceless pieces, of dimensions $1,3,5$.  The
antisymmetric piece is an exterior-square representation; the
symmetric pieces arise from symmetric tensor powers and their
trace decomposition.  Schur functors and invariant contractions
therefore help build features and outputs with known transformation
laws.  Determinants test nonsingularity or volume change in specified
finite-dimensional maps.  They do not determine a turbulent closure
or imply the analytic estimates required for a continuum theorem.

For a quantum version of ToricGT, the surface or dessin specifies
connectivity and admissible local operations, while circuit parameters
specify a family of quantum channels and observables.  A conservation
law can be enforced by restricting channels or observables to a
specified invariant sector.  A measured expectation then estimates a
well-defined feature or coefficient.  State preparation, finite-shot
variance, circuit approximation, and the relation between that feature
and the physical kinetic state remain separate parts of the model.
The quantum-fluid chapter discusses these algorithmic costs and
encodings.  Learning on a surface does not identify an arbitrary
three-dimensional fluid with a two-dimensional toric code.

\section{What passes from complexes to a continuum theorem}

On a finite cell complex, the boundary-of-boundary identity supplies
exact incidence relations.  Its adjoint cochain identities, Hodge
decompositions, and positive semidefinite Laplacians can enforce
discrete conservation and isolate harmonic components.  A gentle
surface algebra or a dessin provides additional combinatorial
organization of local operators.  A variety of complexes organizes
matrices with prescribed compositions and rank conditions.  These
constructions can certify the algebraic structure of a discretization
or constrain a neural architecture.

To pass from such a model to a continuum claim requires a separate
chain of estimates: consistency of the geometric approximation,
stability uniform in mesh size, compactness in a topology suitable
for nonlinear products, identification of the limiting equation,
and preservation of the required initial and boundary data.
Spectral gaps on individual finite graphs need not have a uniform
positive lower bound as the graphs are refined.  Rank constraints
on finitely many matrices need not control high-frequency
concentration.  A learned closure with positive entropy production
need not supply the required continuum regularity.

The constructive role of the project's algebra is to make some
requirements exact and testable, then to identify which analytic
requirements are still independent.  Chapter~\ref{ch:fluid-structures}
proves the finite structural statements, and
Chapter~\ref{ch:navier-stokes-release} identifies the actual theorem
and proof obligations of the recent release.  The microscopic
derivations in this chapter provide the physical interpretation and
the correct hierarchy of approximations without changing those
theorems' hypotheses.

\chapter{The Navier--Stokes release: exact scope and proof interfaces}
\label{ch:navier-stokes-release}

\reviewstatus{The 166-page released manuscript has been read in full across the
two mathematical review passes recorded in the source ledger. The repository
revision, final theorem statements, dependency graph, and selected implementation
modules have been inspected. Reading a mathematical proof, rebuilding a Lean
target, and independently checking an exported kernel certificate are distinct
operations; their actual outcomes are recorded separately. The propositions
proved in this chapter explain mathematical interfaces of the construction.
They do not replace the analytic construction imported from the release.}

\section{What the released theorem says}

OpenAI's release of 8 September 2026 concerns a smooth, externally forced,
three-dimensional incompressible flow. It supplies a manuscript and Lean
formalization. The Clay Mathematics Institute's response of 11 September
acknowledges the apparent resolution while describing an ongoing evaluation
process. These facts must be distinguished from an award decision
\citep{NSOpenAIAnnouncement2026,NSClayAnnouncement2026}.

Let $\nu>0$ and define the momentum residual
\begin{equation}
 \mathcal R_\nu(u,p)
 =\partial_tu+(u\cdot\nabla)u-\nu\Delta u+\nabla p.
 \label{eq:ns-residual}
\end{equation}
The source's main conclusion can be recorded without changing its quantifiers:
\begin{align}
 &f\in C_c^\infty(\R^3\times(0,\infty);\R^3),\qquad
 u,p\in C^\infty(\R^3\times[0,1)),\nonumber\\
 &\mathcal R_\nu(u,p)=f,\quad \nabla\cdot u=0,\quad u(\cdot,0)=0,
 \label{eq:ns-release-equation}\\
 &\supp u(\cdot,t)\cup\supp p(\cdot,t)\subset K
 \quad(0\leq t<1),\nonumber\\
 &\sup_{t<1}\|u(t)\|_2<\infty,\qquad
 \limsup_{t\uparrow1}\|u(t)\|_\infty=\infty.
 \label{eq:ns-release-growth}
\end{align}
Here $K$ is compact. With the same datum and force there is no global smooth
whole-space solution with uniformly bounded kinetic energy. A corresponding
periodic result is also established \citep[Theorem 1.1 and Corollary 10.6]{NSOpenAI2026}.

In Fefferman's formulation, alternatives (A) and (B) impose $f=0$, whereas
(C) and (D) permit a smooth force satisfying the specified decay conditions.
The released statements address (C) and (D). They do not logically negate (A)
or (B): an existential assertion with $f\neq0$ does not supply an example with
$f=0$. Pressure periodicity is part of the periodic formulation, including its
erratum \citep{NSFefferman}. The separate Euler result concerns a different
equation, obtained by setting viscosity to zero, and cannot be transferred to
positive viscosity by deleting the viscous estimates.

This distinction changes the appropriate goal of the present extension. A
more informative presentation can expose the cancellation, support, analytic
norms, and formal dependency interfaces. A stronger theorem would require a
new mathematical conclusion with all its additional hypotheses discharged.
The material below provides the former and identifies precisely where the
source's construction supplies the latter's indispensable analytic inputs.

\section{An exact algebra for residual correction}

The same distinction between a complex and an exact complex encountered
earlier reappears here. A differential identity constrains what a correction
can do. It does not prove that a correction of the necessary size, support,
and regularity exists.

\begin{proposition}[Quadratic residual identity]
\label{prop:ns-residual-increment}
For smooth $u,w,p,\pi$, with $\nabla\cdot w=0$,
\begin{align*}
 \mathcal R_\nu(u+w,p+\pi)
 &=\mathcal R_\nu(u,p)+L_{u,\nu}(w,\pi)
       +\nabla\cdot(w\otimes w),\\
 L_{u,\nu}(w,\pi)
 &=\partial_tw+(u\cdot\nabla)w+(w\cdot\nabla)u
       -\nu\Delta w+\nabla\pi .
\end{align*}
The tensor convention is $(w\otimes w)_{ij}=w_iw_j$ and
$(\nabla\cdot S)_i=\sum_j\partial_jS_{ij}$.
\end{proposition}
\begin{proof}
All terms other than advection are linear. Expanding the latter gives
$(u+w)_j\partial_j(u_i+w_i)$, with the four products obtained by choosing
one of $u,w$ in each factor. Moreover,
$\sum_j\partial_j(w_iw_j)=\sum_jw_j\partial_jw_i+
w_i\sum_j\partial_jw_j=(w\cdot\nabla)w_i$. This proves every component
of the identity. No averaging, limiting argument, or neglected term occurs.
\end{proof}

In particular, solving $L_{u,\nu}(w,\pi)=-\mathcal R_\nu(u,p)$ leaves the
quadratic residual $\nabla\cdot(w\otimes w)$. A numerical decrease in the
linear residual alone is therefore an incomplete verification. Conversely,
if a background satisfies
$\mathcal R_\nu(u,p)=-\nabla\cdot S+E$, a fluctuation with averaged
covariance $S$ can cancel the selected mean stress while creating other
linear and oscillatory terms. Each of those terms needs a place in the next
residual estimate.

\begin{proposition}[A transverse wave and its covariance]
\label{prop:ns-transverse-wave}
Let $a,k\in\R^3$ with $a\cdot k=0$ and put
$w(x,\phi)=a\cos(k\cdot x+\phi)$. Then
\[
 \nabla\cdot w=0,\qquad (w\cdot\nabla)w=0,\qquad
 \frac1{2\pi}\int_0^{2\pi}w\otimes w\,d\phi
       =\frac12a\otimes a.
\]
The vanishing self-advection statement requires constant $a$ and $k$.
\end{proposition}
\begin{proof}
Differentiating the cosine gives
$\nabla\cdot w=-(a\cdot k)\sin(k\cdot x+\phi)$ and
$(w\cdot\nabla)w=-(a\cdot k)a\cos(k\cdot x+\phi)
\sin(k\cdot x+\phi)$. Both vanish. Integration of
$\cos^2s=(1+\cos2s)/2$ proves the tensor identity.
If the amplitude or phase gradient varies, the product rule introduces
additional derivatives, so this proof gives no cancellation for them.
\end{proof}

The covariance is an element of $\Sym^2(\R^3)$, with nonnegative quadratic
form $v\mapsto (a\cdot v)^2/2$. The positivity is a concrete use of a
symmetric tensor power. A rank-one tensor representation does not by itself
solve the vector-valued stress-divergence equation; the allowed polarization
directions and the differential constraints still matter. The exterior-power
test for independence of two stress directions is developed in
Chapter~\ref{ch:fluid-structures}.

\begin{proposition}[Averaging and nonlinear closure]
\label{prop:ns-average-defect}
Let $P$ be a linear averaging operator on a function algebra with $P1=1$.
Assume $P^2=P$ and the conditional-expectation identity
$P(ag)=aP(g)$ whenever $Pa=a$. If $v=Pu$ and $w=u-v$, then
\[
 P(u\otimes u)=v\otimes v+P(w\otimes w).
\]
If $P$ is positive, the final term is positive semidefinite pointwise in
the averaged variables.
\end{proposition}
\begin{proof}
Idempotence gives $Pw=0$. Expand $(v+w)\otimes(v+w)$. The two mixed
terms have zero average by the assumed multiplication identity, and
$P(v_i v_j)=v_iv_j$. For a constant test vector $b$,
$b^TP(w\otimes w)b=P((b\cdot w)^2)\geq0$.
\end{proof}

An orthogonal projection onto an arbitrary collection of modes need not have
the multiplication property in this proposition. Likewise, averaging an
extended function $F(x,Y)$ over independent $Y$ and then substituting a phase
$Y=\Phi(x)$ is different from first restricting and averaging the resulting
physical function. If $F(x,Y)=\cos(2\pi(Y-x))$ on one-dimensional tori,
then $\int F(x,Y)dY=0$, whereas $F(x,x)=1$. This elementary example prevents
an unjustified replacement of a pointwise physical residual by an auxiliary
mean. The phase derivatives and support geometry in a complete construction
are precisely what allow these two kinds of expression to be related.

\section{Concentration, energy, and the information lost by norms}

\begin{proposition}[Energy estimate from rest]
\label{prop:ns-force-energy}
Let $u,p$ be a smooth solution on $[0,T]$ of
$\mathcal R_\nu(u,p)=f$, $\nabla\cdot u=0$, $u(0)=0$.
Assume periodicity, or sufficient decay for the integrations by parts
below; fixed compact spatial support of $u,p$ is sufficient. If
$F(t)=\int_0^t\|f(s)\|_2\,ds$, then
\[
 \|u(t)\|_2^2+2\nu\int_0^t\|\nabla u(s)\|_2^2\,ds\leq F(t)^2.
\]
\end{proposition}
\begin{proof}
The pressure pairing is
$\int u\cdot\nabla p=-\int p\nabla\cdot u=0$.
Advection contributes
$\int u\cdot(u\cdot\nabla)u=
\frac12\int\nabla\cdot(u|u|^2)=0$. Thus
\[
 \frac12\frac d{dt}\|u\|_2^2+\nu\|\nabla u\|_2^2=\langle f,u\rangle.
\]
Write $y_\delta=(\|u\|_2^2+\delta^2)^{1/2}$, $\delta>0$.
Then $y_\delta'\leq\|f\|_2$, so
$y_\delta(t)\leq\delta+F(t)$. Passing to $\delta\downarrow0$ gives
$\|u(t)\|_2\leq F(t)$. Integrating the energy identity once more,
\[
 \|u(t)\|_2^2+2\nu\int_0^t\|\nabla u\|_2^2
 \leq2\int_0^t F(s)F'(s)\,ds=F(t)^2.
\]
The argument on each $T<1$, followed by monotone convergence of the
nonnegative dissipation integral, also applies up to a singular time.
\end{proof}

This estimate controls an integral, not the maximum of the integrand.
Information reduction becomes mathematically consequential at this exact
point: a small spatial set can carry a large amplitude while making a small
contribution to $L^2$ energy.

\begin{proposition}[Anisotropic concentration test]
\label{prop:ns-anisotropic-scaling}
Let $0<h<1/6$, $\tau>0$, and $W\in C_c^\infty(\R^3;\R^m)$ be nonzero.
Define the comparison family
\[
 V_\tau(x_1,x_2,x_3)=\tau^{-1/2-h}
 W\left(\frac{x_1}{\tau^{1/2}},
        \frac{x_2}{\tau^{1/2}},
        \frac{x_3}{\tau^{1/2-h}}\right).
\]
For $1\leq p<\infty$,
\[
 \|V_\tau\|_p^p
  =\tau^{3/2-h-p(1/2+h)}\|W\|_p^p,
 \qquad
 \|V_\tau\|_\infty=\tau^{-1/2-h}\|W\|_\infty.
\]
Its energy tends to zero, while its maximal amplitude diverges. Moreover,
$\int_0^{\tau_0}\|\nabla V_\tau\|_2^2\,d\tau<\infty$ for every finite
$\tau_0>0$. This comparison family is not asserted to solve Navier--Stokes.
\end{proposition}
\begin{proof}
The change of variables has determinant $\tau^{3/2-h}$. This proves the
norm identities. At $p=2$ the exponent is $1/2-3h>0$. Each of the first
two spatial derivatives contributes $\tau^{-1/2}$, so its squared norm has
exponent $-1/2-3h>-1$. The axial derivative contributes
$\tau^{-1/2+h}$ and has squared-norm exponent $-1/2-h>-1$.
Both powers are integrable at zero. The supremum formula follows because
the coordinate dilation is onto.
\end{proof}

For this family the critical spatial integrability exponent is
\[
 p_*(h)=\frac{3/2-h}{1/2+h}<3.
\]
At $p=3$, $\|V_\tau\|_3^3=\tau^{-4h}\|W\|_3^3$, which diverges.
Thus the scaling test does not contradict a regularity theorem requiring
a bounded critical $L^3$ norm. Such a theorem is a substantially stronger
statement than an energy bound; the endpoint theorem for the unforced
Cauchy problem is due to Escauriaza, Seregin, and \v Sver\'ak
\citep{NSEscauriazaSereginSverak2003}. Our comparison family proves only
the logical insufficiency of the weaker norm information, independently
of any existence construction.

\section{Similarity coordinates as an exact change of variables}

The geometry relevant to the release can be organized using the positive
variables
\[
 \tau=1-t,\qquad D=\tfrac12-h,\qquad
 \tau=q(1-\eta^2),\quad z=q^D\eta,\quad X=\frac{r^2}{2q},
 \quad 0<h<\tfrac12.
\]
The next calculation verifies the differential interface rather than
assuming that symbolic substitution preserves all derivatives.

\begin{proposition}[Coordinate existence and derivative weights]
\label{prop:ns-similarity-chain-rule}
For $\tau>0$ and $z\in\R$ there is a unique
$q>|z|^{1/D}$ with $q-z^2q^{2h}=\tau$. Set $\eta=zq^{-D}$,
$d=1-\eta^2$, and $L=1-2h\eta^2$. Then $|\eta|<1$, $L\geq1-2h>0$,
and for a smooth profile $G(X,\eta)$ and any real $b$,
\begin{align*}
 \partial_t(q^bG)&=q^{b-1}L^{-1}
     (-bG+D\eta\partial_\eta G+X\partial_XG),\\
 \partial_z(q^bG)&=q^{b-D}L^{-1}
     (2b\eta G+d\partial_\eta G-2\eta X\partial_XG).
\end{align*}
\end{proposition}
\begin{proof}
On $q>|z|^{1/D}$, the function $q-z^2q^{2h}$ begins at zero,
tends to infinity, and has derivative
$1-2hz^2q^{2h-1}\geq1-2h>0$. The intermediate-value theorem and strict
monotonicity give existence and uniqueness. The implicit-function theorem
gives smoothness, and differentiation yields
\[
 q_t=-L^{-1},\quad q_z=2\eta q^{1-D}/L,\quad
 \eta_t=D\eta/(qL),\quad\eta_z=d/(q^DL),
\]
\[
 X_t=X/(qL),\qquad X_z=-2\eta X/(q^DL).
\]
The two formulas follow by the product and chain rules. The endpoint $z=0$
also satisfies them, because the implicit derivative remains nonzero.
\end{proof}

On bounded profile ranges, a transverse derivative has weight $q^{-1/2}$,
an axial derivative has weight $q^{-D}$, and a time derivative has weight
$q^{-1}$. The denominators remain bounded because $L\geq1-2h$.
Constants can depend on the derivative order and fixed profile parameters.
They cannot be allowed to hide an extra loss growing with the approximation
stage if one later intends to sum an infinite correction sequence.

The apparent cylindrical singularity at $r=0$ must also be resolved.
For instance, if the azimuthal velocity is $r\,a(r^2,z,t)$, then
\[
 r\,a(r^2,z,t)e_\theta
   =a(x_1^2+x_2^2,z,t)(-x_2,x_1,0)
\]
is smooth in Cartesian coordinates. Boundedness of a cylindrical component
alone would not establish this. This is a local regularity condition on a
representation, comparable to requiring that coordinate expressions of a
section glue smoothly; it is not supplied by the topology of the domain.

\section{A quadratic relation that controls analytic derivative loss}

An especially useful connection to the project's algebraic methods occurs
in the analytic linear recursion. The relevant feature is a forbidden
composition of derivative-carrying blocks. We isolate its mechanism in a
statement independent of fluid terminology.

\begin{lemma}[Words with no adjacent derivative letters]
\label{lem:ns-forbidden-words}
A word of length $k$ in the letters $N,D$ with no consecutive $D$ has at
most $\lceil k/2\rceil$ occurrences of $D$. The number with exactly $j$
occurrences is $\binom{k-j+1}{j}$, interpreted as zero when $j>k-j+1$.
The total number $a_k$ satisfies $a_0=1$, $a_1=2$, and
$a_k=a_{k-1}+a_{k-2}$ for $k\geq2$.
\end{lemma}
\begin{proof}
Between $j$ derivative letters there are at least $j-1$ other letters, so
$2j-1\leq k$. If their positions are
$1\leq i_1<\cdots<i_j\leq k$ with $i_{r+1}\geq i_r+2$, the shifted
positions $i_r-(r-1)$ form an ordinary $j$-element subset of
$\{1,\ldots,k-j+1\}$. This is a bijection, proving the binomial formula.
For the recurrence, an admissible word ending in $N$ is an arbitrary
admissible word of length $k-1$ followed by $N$. One ending in $D$, when
$k\geq2$, ends in $ND$ and corresponds to a word of length $k-2$.
These two classes are disjoint and exhaustive.
\end{proof}

\begin{proposition}[Sparse Volterra systems]
\label{prop:ns-sparse-volterra}
Let $E=E_0\oplus E_1$ be a finite-dimensional normed complex vector space.
Fix a compact real interval $I$ and let
$S_\rho=\{z\in\C:\dist(z,I)<\rho\}$. Consider an integral equation
on $0\leq s\leq a$ and this complex parameter neighborhood,
\[
 W=Gf+G(A_0W+A_1\partial_\eta W),\qquad
 (Gg)(s)=\int_0^s K(s,r)g(r)\,dr.
\]
Assume $\|K(s,r)\|\leq1$, $K$ preserves the two summands and is independent
of $\eta$; the coefficients and kernel are continuous in their radial
variables; $A_0,A_1,f$ are uniformly bounded and holomorphic in $\eta$;
and every value and every $\eta$-derivative of $A_1$ maps $E_0$ into $E_1$
and annihilates $E_1$. On each smaller neighborhood $S_{\rho'}$ with
$0<\rho'<\rho$, the Picard series converges for every finite $a$.
It gives the unique solution in the class of solutions bounded and
holomorphic on a parameter neighborhood strictly larger than the one
on which they are compared.
\end{proposition}
\begin{proof}
Expand the $k$-fold iterate of $G(A_0+A_1\partial_\eta)$ into words.
The kernel between adjacent coefficients preserves $E_0,E_1$. Therefore
\[
 A_1(s)K(s,r)A_1(r)=0,
 \qquad A_1(s)K(s,r)\partial_\eta A_1(r)=0.
\]
The second identity is necessary: in a composition of two derivative
operators the first derivative can strike the second coefficient. It
vanishes for the same block reason. Thus a nonzero word cannot contain
two adjacent derivative operators, even after expanding the product rule.
By Lemma~\ref{lem:ns-forbidden-words}, at most
$p_k=\lceil k/2\rceil$ Cauchy radius losses are required.

Put $\Delta=\rho-\rho'>0$. Distribute that loss in $p_k$ equal parts.
The Cauchy estimate for one parameter derivative costs at most
$p_k/\Delta$; when $p_k=0$ use cost one. Operator norms absorb the
coefficient bounds, and the nested integration domain has volume
$a^{k+1}/(k+1)!$. Bounding the number of words by $2^k$ gives, after
enlarging a constant $C$,
\[
 \bigl\|[G(A_0+A_1\partial_\eta)]^kGf\bigr\|_{S_{\rho'}}
 \leq\frac{C^{k+1}a^{k+1}}{(k+1)!}
       \max\{1,p_k/\Delta\}^{p_k}.
\]
Stirling's bound shows that the $k$th root of the right-hand side is
$O(k^{-1/2})$, tending to zero. Hence the series converges uniformly on
the smaller neighborhood. Applying the same estimate after one additional
radius restriction permits differentiation in $\eta$ and termwise
substitution into the integral equation. If $W_1-W_2=G(A_0+A_1\partial_\eta)
(W_1-W_2)$, iteration and the same estimate make the difference zero on
the smaller neighborhood. Analytic continuation gives the stated uniqueness.
\end{proof}

The source's six-component linear recursion has the splitting
$E_0=\C^4$, $E_1=\C^2$, with diagonal radial kernels. Its use of the
above relation is traceable to Lemma 5.1, rather than inferred from a
superficial resemblance to a quiver \citep{NSOpenAI2026}. The relation
has three precise interpretations. It eliminates forbidden operator words;
it limits derivative loss; and, with the integration simplex, it makes a
convergent analytic solve possible on a common radial interval.

The language of admissible words is regular: a finite automaton needs only
to remember whether the preceding letter was $D$, plus a rejecting state.
The analytic conclusion nevertheless requires unbounded iteration and
quantitative estimates. A regular syntax does not make the represented
analytic task finite-dimensional.

Nor does this calculation establish a gentle algebra. The abstract
one-vertex algebra $k\langle N,D\rangle/(D^2)$ is infinite-dimensional,
and the arrow $N$ has two allowed successors. It therefore fails a usual
gentle-algebra continuation condition. The honest bridge is to a
representation with a quadratic zero relation and a controlled path
language. Any claim of a particular gentle presentation would require
a separate equivalence or construction.

\section{From formal orders to actual smooth fields}

The order of a term in a formal expansion must not be confused with a
convergence estimate for the coefficients. Suppose coefficients of order
$j$ have constants growing faster than every exponential. Positive powers
of a small parameter do not automatically make their unrestricted sum
converge. A suitable realization uses cutoffs chosen after the coefficients
are known.

\begin{lemma}[A diagonal realization estimate]
\label{lem:ns-diagonal-realization}
Let $0<q<q_0\leq1$ on an open domain, with $q$ smooth and bounded below
on every compact subset. Suppose smooth increments $Z_j$, $j\geq1$,
satisfy
\[
 |Z_j|_m\leq C_{j,m}q^{g_j-\ell_m}(1+|\log q|)^{P_{j,m}},
 \qquad 0<g_1\leq g_2\leq\cdots\longrightarrow\infty,
\]
where $|\cdot|_m$ is the pointwise maximum of coordinate derivatives through order
$m$, and $\ell_m$ is independent of $j$. Assume that differentiating a
cutoff $\chi(aq)$ costs a power of $q^{-1}$ depending only on derivative
order, uniformly in $a$ on its transition region. Then one can choose
$a_{j+1}\geq2a_j$ so that the cutoff sum
$Z=\sum_{j\geq1}\chi(a_jq)Z_j$ is locally finite and smooth, and its
tail after a fixed initial block satisfies
\[
 \left|Z-\sum_{j=1}^JZ_j\right|_m
 \leq2^{-J}q^{g_{J+1}/2-L_m}
 \quad\left(J\geq\max\{1,m\},\ q<\frac1{2a_J}\right).
\]
Here $\chi=1$ on $[0,1/2]$, $\chi=0$ on $[1,\infty)$, and $L_m$ is
independent of $J$.
\end{lemma}
\begin{proof}
The product rule and the cutoff hypothesis bound the $m$th seminorm of
the $j$th summand by
$\widehat C_{j,m}q^{g_j-L_m}(1+|\log q|)^{\widehat P_{j,m}}$.
At stage $j$ impose this finite list of requirements for $0\leq m\leq j$:
\[
 \widehat C_{j,m}q^{g_j/2}(1+|\log q|)^{\widehat P_{j,m}}
 \leq2^{-j}\quad(0<q\leq a_j^{-1}).
\]
It is possible because every positive power of $q$ dominates every fixed
power of $|\log q|$ at zero. Increase $a_j$ further to preserve its geometric
growth. On a compact subset where $q\geq\delta>0$, terms with
$a_j>\delta^{-1}$ vanish, proving local finiteness and smoothness.
If $q<(2a_J)^{-1}$, the first $J$ cutoffs equal one. For $j>J\geq m$,
the bound just imposed and $g_j\geq g_{J+1}$ give
\[
 \sum_{j>J}2^{-j}q^{g_j/2-L_m}
 \leq2^{-J}q^{g_{J+1}/2-L_m},
\]
as required.
\end{proof}

To preserve a differential constraint, apply the cutoffs to suitable
potentials. In three-dimensional Euclidean space,
\[
 \nabla\times(\chi(a_jq)A_j)
 =\chi(a_jq)\nabla\times A_j
   +a_j\chi'(a_jq)\nabla q\times A_j.
\]
The second term is a real part of the field and must be retained. Taking
divergence of the full curl yields zero. Cutting off an already
divergence-free velocity directly would instead give
$\nabla\cdot(\chi u)=\nabla\chi\cdot u$, which need not vanish.

\begin{proposition}[All-order residual realization]
\label{prop:ns-residual-flatness}
Let $U^{[J]}$ be finite approximations and $U$ their cutoff realization.
Assume their derivatives have polynomial growth in $q^{-1}$ with
exponent losses independent of $J$, and that the tail estimate of
Lemma~\ref{lem:ns-diagonal-realization} applies through every fixed order.
Let $F$ be a fixed differential polynomial of finite order and degree whose
coefficient derivatives have polynomial growth in $q^{-1}$. If
\[
 |F(U^{[J]})|_m
 \leq C_{J,m}q^{\rho_J-K_m}(1+|\log q|)^{P_{J,m}}+E_{J,m},
 \quad \rho_J\to\infty,
\]
with $K_m$ independent of $J$ and each fixed $E_{J,m}=O(q^N)$ for every
$N$, then $|F(U)|_m=O(q^N)$ for every $m,N$.
\end{proposition}
\begin{proof}
Fix $m,N$. Let $s,d$ bound the differential order and polynomial degree
of $F$. For products use the exact identity
\[
 \prod_{i=1}^d a_i-\prod_{i=1}^d b_i
 =\sum_{r=1}^d
  \left(\prod_{i<r}a_i\right)(a_r-b_r)
  \left(\prod_{i>r}b_i\right).
\]
The product rule then bounds $F(U)-F(U^{[J]})$ in seminorm $m$ by a
fixed power of $q^{-1}$ times the tail seminorm $m+s$ and a polynomial
in the two field seminorms $m+s$. All exponent losses depend on $m,s,d$,
not $J$. Choose $J$ large enough that both $\rho_J$ and $g_{J+1}/2$
exceed $N$ plus these losses and one extra unit. Keep that $J$ fixed.
Shrink the $q$-neighborhood so all its finitely many constants and logarithms
are absorbed by that extra unit, and the first $J$ cutoffs equal one.
The finite residual, its fixed flat remainder, and the residual difference
are then all $O(q^N)$.
\end{proof}

The order of these quantifiers is essential:
\[
 \forall(m,N)\ \exists J\ \exists\delta,C\ 
  \forall q\in(0,\delta):\ |F(U)|_m\leq Cq^N.
\]
It does not assert one $J$ that controls every derivative. It does not
require summing the flat remainders $E_{J,m}$ over $J$. It also does not
assert convergence of the uncut formal series. This distinction is the
analytic counterpart of keeping a filtered algebra, its associated graded
object, and a completed realization separate.

\section{The finite correction cycle as a verifiable dependency graph}

The source's construction provides the hypotheses needed above through
separate analytic modules. The following locator table records where to
check them; the table is a dependency guide, not a substitute proof
\citep{NSOpenAI2026,NSLean2026}.

\begin{longtable}{p{.20\textwidth}p{.40\textwidth}p{.30\textwidth}}
\toprule
Interface & Mathematical obligation & Source location\\
\midrule\endhead
Background & Smooth axis; exact exterior; supported compatible stress;
positive cone margins & Theorem 4.6; Appendices A--C\\
Order recursion & A common radial interval, moment repair, derivative losses
independent of formal order & Lemmas 5.1--5.4; Proposition 5.5\\
Phase geometry & Support separation and compatible finite covering coordinates
& Section 6\\
Oscillatory solve & Transverse evolution, damping, real covariance, full curl
correction & Propositions 7.2, 7.5--7.6; Lemma 7.7\\
Mean solve & Compact primitives, auxiliary inverse, five moment equations
& Section 8\\
Iteration & Full residual recomputation, fixed-domain gains, all physical
derivatives & Propositions 9.3--9.9\\
Completion & Smooth forcing, energy, comparison, all viscosities, periodization
& Section 10\\
\bottomrule
\end{longtable}

One can encode a correction state as a typed record rather than a loose
list of estimated errors. Its fields should include a velocity, pressure,
potential representation, support certificates, exact preserved moments,
residual decomposition, and derivative bounds. A correction operation
must produce a new such record. The pressure and residual are recomputed
from the updated velocity, so one cannot reuse a previous covariance after
altering a wave amplitude. This record-based reading is directly compatible
with the typed graph approach developed earlier in the monograph.

The numerical exponents also permit an elementary exact check. Put
$\kappa=10^{-5}$, $\sigma_j=1/5+j/10$,
$B_j=1/2+\sigma_j$, and $C_j=1+\sigma_j$.
All asserted gains of the source's cycle must exceed $1/10$ after
derivative losses. For example,
\[
 \frac{18}{100}-2\kappa>\frac{17}{100}>\frac1{10},\qquad
 \sigma_j-3\kappa\geq\frac15-3\kappa>\frac{17}{100},
\]
\[
 \frac12-2\kappa>\frac1{10},\qquad
 \frac9{10}-4\kappa>\frac1{10}.
\]
These rational inequalities certify arithmetic slack once the functional
estimates are supplied. They do not establish those functional estimates.
This is a useful separation for learned proof search: a model may propose
an ordering or an exponent assignment, while exact arithmetic and analytic
lemmas discharge different obligations.

\section{Localization and smooth force extension}

\begin{proposition}[Localization without a divergence defect]
\label{prop:ns-potential-localization}
Suppose $A$, $B e_\theta$ and $p$ are smooth in Cartesian coordinates,
put $u=\nabla\times A+B e_\theta$, and assume $B$ is independent of $\theta$.
Let $c$ be a smooth axisymmetric scalar.
Set
\[
 \widetilde u=\nabla\times(cA)+cBe_\theta,
 \qquad \widetilde p=cp,\qquad a=\nabla c\times A.
\]
Then $\widetilde u=cu+a$ and $\nabla\cdot\widetilde u=0$.
Its residual is exactly
\begin{align*}
 \mathcal R_\nu(\widetilde u,cp)
 ={}&c\mathcal R_\nu(u,p)
   +(\partial_tc-\nu\Delta c)u-2\nu\sum_j(\partial_jc)\partial_ju
       +p\nabla c\\
 &+(c^2-c)(u\cdot\nabla)u+c(u\cdot\nabla c)u\\
 &+\partial_ta-\nu\Delta a+c(u\cdot\nabla)a
      +(a\cdot\nabla)(cu)+(a\cdot\nabla)a.
\end{align*}
\end{proposition}
\begin{proof}
The curl product rule gives $\widetilde u=cu+a$. Divergence of a curl
is zero; the other term has divergence $r^{-1}\partial_\theta(cB)=0$
away from the axis and hence everywhere by Cartesian smoothness.
Expand $\partial_t(cu)$, $\Delta(cu)$, and $\nabla(cp)$ by the product
rule. Finally,
$(cu\cdot\nabla)(cu)=c^2(u\cdot\nabla)u+c(u\cdot\nabla c)u$.
Adding the three advection terms containing $a$ yields the formula.
\end{proof}

The displayed expression explains why a spatial cutoff requires endpoint
regularity of the actual potentials as well as the velocity. On a region
where $c=1$ the extra terms vanish. On its transition region every term
involving the uncut fields must remain smooth through the terminal time.
A proof that estimates only $c\mathcal R_\nu(u,p)$ has omitted real terms.

\begin{lemma}[Realization of terminal force jets]
\label{lem:ns-terminal-jets}
Let $f(x,t)$ be smooth for $t<T$, supported in a fixed compact set
$K\subset\R^d$. Suppose every $\partial_x^\alpha\partial_t^jf$
has a uniform limit $\partial_x^\alpha F_j(x)$ as $t\uparrow T$,
where $F_j\in C_c^\infty(\R^d)$ are supported in $K$.
Then $f$ has a smooth extension to $t\geq T$ supported in
$K\times[T,T+1]$ on that side. If it already vanishes near an initial
time, the combined extension retains that property.
\end{lemma}
\begin{proof}
Choose a smooth cutoff $\chi$, one near zero and zero for arguments
at least one, and for $s=t-T\geq0$ define
\[
 f_+(x,T+s)=\sum_{j=0}^\infty\chi(b_js)\frac{s^j}{j!}F_j(x).
\]
For fixed $\alpha,m$, the product rule and $s\leq b_j^{-1}$ on the
summand's support give a bound
\[
 C_{j,m}b_j^{m-j}\|F_j\|_{C^{|\alpha|}}.
\]
Fix $b_0=1$. For $j\geq1$, choose increasing $b_j\geq1$ so that this is at most $2^{-j}$ whenever
$|\alpha|+m\leq\lfloor j/2\rfloor$. Each step imposes finitely many
constraints with $m-j<0$, so the choice exists. Every differentiated tail
converges uniformly. At $s=0$, the $m$th time derivative of a fixed
summand is $F_m$ when $j=m$ and zero otherwise. Uniform differentiated
convergence gives the desired jets.

Uniform convergence of the original mixed derivatives makes their limits
compatible: pass to the limit in the fundamental theorem of calculus
along a coordinate segment, and then along a time segment. The two
pieces therefore glue smoothly at $T$. All summands retain support in
$K$, and vanish for $s\geq1$; uniform convergence of every derivative
also proves smoothness at that support boundary.
\end{proof}

Flatness at a single singular point is only part of this lemma's hypotheses.
One must additionally control all terminal derivatives away from that point,
including the cutoff transition region. Conversely, vanishing of every jet
does not force a smooth function to vanish near the point: the function
$e^{-1/s^2}$ for $s>0$, extended by zero for $s\leq0$, is the standard
counterexample. Substituting ``analytic'' for ``smooth'' would change the
problem and invalidate this extension method.

\section{Comparison, quantitative residual checks, and proof strength}

An explicit singular solution rules out a global competitor only after
uniqueness has been established in the competitor's class. The whole-space
class in the Clay problem permits no arbitrary decay assumption on the
pressure. That is why a periodic energy calculation is not by itself a
complete proof of alternative (C). The source's Lemma 10.5 handles this
point using a pressure-gradient argument and expanding spatial cutoffs;
its hypotheses must be preserved in any proof adapter
\citep{NSOpenAI2026}.

The following separate estimate is useful for the project's computable
fluid examples and for any learned approximate field.

\begin{proposition}[Residual-to-solution estimate on the flat torus]
\label{prop:ns-residual-comparison}
Let $u,v$ be smooth divergence-free periodic vector fields on $[0,T]$,
and let periodic $p,q$ satisfy
\[
 \mathcal R_\nu(u,p)=f,\qquad
 \mathcal R_\nu(v,q)=f+r.
\]
Write $w=v-u$ and $K(t)=\|\nabla u(t)\|_{L^\infty,\mathrm{op}}$.
If $K$ is integrable, then
\[
 \|w(t)\|_2\leq
 e^{\int_0^tK(a)\,da}\|w(0)\|_2
 +\int_0^t e^{\int_s^tK(a)\,da}\|r(s)\|_2\,ds.
\]
In particular, equal data and $r=0$ imply $v=u$.
\end{proposition}
\begin{proof}
Subtract the equations in the form
\[
 \partial_tw+(v\cdot\nabla)w+(w\cdot\nabla)u
 -\nu\Delta w+\nabla(q-p)=r.
\]
Periodic integration and zero divergence remove the transport and pressure
pairings. Thus
\[
 \tfrac12(\|w\|_2^2)'+\nu\|\nabla w\|_2^2
 \leq K\|w\|_2^2+\|r\|_2\|w\|_2.
\]
With $y_\delta=(\|w\|_2^2+\delta^2)^{1/2}$ one has
$y_\delta'\leq Ky_\delta+\|r\|_2$. Multiplication by
$e^{-\int_0^tK}$ and integration prove the bound after
$\delta\downarrow0$.
\end{proof}

This inequality gives a precise role for a quantum or classical learner:
it may produce $v$, but a certificate needs its divergence defect, momentum
residual, initial error, and an appropriate stability coefficient. The
estimate does not remain effective near blowup merely because the sampled
residual is small. A pointwise grid evaluation does not bound the continuum
$L^2$ residual without an interpolation or enclosure argument, and
$\int K$ can become large. Numerical fitting and a posteriori proof are
different stages of the same workflow.

The same issue explains why nearby classical results are not interchangeable.
Nonuniqueness in a weak class does not give blowup of a smooth solution;
a modified averaged nonlinearity is not the original Navier--Stokes
nonlinearity; and rough forcing is not smooth forcing. These distinctions
are visible respectively in the settings of Buckmaster--Vicol,
Tao, and Albritton--Bru\'e--Colombo
\citep{NSBuckmasterVicol2019,NSTao2016,NSAlbrittonBrueColombo2022}.

\section{Viscosity, periodization, and preservation of quantifiers}

\begin{proposition}[Viscosity conversion]
\label{prop:ns-viscosity-conversion}
If $\mathcal R_1(u,p)=f$ and $\nabla\cdot u=0$, then for $\nu>0$
\[
 u_\nu(x,t)=\sqrt\nu\,u(x/\sqrt\nu,t),\quad
 p_\nu(x,t)=\nu p(x/\sqrt\nu,t),\quad
 f_\nu(x,t)=\sqrt\nu\,f(x/\sqrt\nu,t)
\]
satisfy $\mathcal R_\nu(u_\nu,p_\nu)=f_\nu$ and zero divergence.
The singular time is unchanged. In three dimensions,
\[
 \|u_\nu(t)\|_2^2=\nu^{5/2}\|u(t)\|_2^2,
 \qquad
 \nu\|\nabla u_\nu(t)\|_2^2=\nu^{5/2}\|\nabla u(t)\|_2^2.
\]
\end{proposition}
\begin{proof}
Put $y=x/\sqrt\nu$. Each term of the momentum equation acquires a
factor $\sqrt\nu$: two spatial derivatives contribute $\nu^{-1}$,
one contributes $\nu^{-1/2}$, and the amplitudes contribute the stated
powers. The divergence equals $(\nabla_y\cdot u)(y,t)$. The Jacobian
of $x=\sqrt\nu y$ is $\nu^{3/2}$, while the velocity square contributes
$\nu$ and the gradient square contributes one. This proves both norm
identities. The time argument was not changed.
\end{proof}

\begin{proposition}[Disjoint-support periodization]
\label{prop:ns-periodization}
Suppose smooth whole-space fields $u,p,f$ solve Navier--Stokes with
viscosity $\nu$ and have common spatial support in a compact set
$K\Subset(-1/2,1/2)^3$. Then their sums over integer translates define
smooth periodic fields solving the same equation, with the same viscosity.
\end{proposition}
\begin{proof}
The sums are locally finite. The distinct translated compact supports have
positive separation, so at each point a neighborhood intersects at most
one support, unless the fields vanish there. Differential operators
commute with the sums. In the advection term every mixed product of
different translates vanishes, leaving the sum of the individual
advection terms. The same argument gives zero divergence and retains
periodicity of both velocity and pressure.
\end{proof}

For a larger compact support one first uses the parabolic scaling
\[
 \widetilde u(x,t)=\lambda u(\lambda x,\lambda^2(t-t_0)),\quad
 \widetilde p(x,t)=\lambda^2p(\lambda x,\lambda^2(t-t_0)),\quad
 \widetilde f(x,t)=\lambda^3f(\lambda x,\lambda^2(t-t_0)).
\]
Choosing $\lambda$ large puts the support inside the cube. If the original
fields vanish on an initial interval, they extend smoothly by zero before
$t_0$. For original singular time one, $t_0=1-\lambda^{-2}$ keeps the
new singular time equal to one. These conditions explain why periodizing
an arbitrary spatially decaying field would not suffice: overlapping copies
would create additional nonlinear interactions.

\section{Formal proof provenance and what this extension adds}

The pinned upstream revision is
\begin{center}\small
\texttt{f9e8bc5b38b6e212696e8a30e3e91517af887bbd}.
\end{center}
Its toolchain is Lean 4.34.0-rc2, with mathlib revision
\begin{center}\small
\texttt{85e3a25e006c35636f0e53b0e9296caca2685bc0}.
\end{center}
The project's earlier selected Lean declarations use a different toolchain.
They cannot be merged by silently substituting one compiler for the other.
An isolated verification script preserves the upstream pins.

The final upstream declarations are
\begin{quote}\small\ttfamily
NavierStokes.Comparator.navier\_stokes\_breakdown\_R3\\
NavierStokes.Comparator.navier\_stokes\_breakdown\_periodic.
\end{quote}
Their types quantify over every positive viscosity and retain the exact
initial-data, forcing, regularity, energy, and periodicity predicates.
The reference challenge modules intentionally contain unproved challenge
statements. They are kept separate from the solution imports. A source
scan that merely counts every occurrence of \texttt{sorry} in the whole
repository would therefore misdescribe the solution's dependency boundary.
Conversely, a source scan finding no such token cannot replace kernel
checking or comparison of the definitions with the intended mathematics.

The inspection found 816 files under the Navier--Stokes source directory
and 609 local modules in the final solution's import closure. Selected
modules were read in detail; all 609 were not manually reviewed. The
upstream metadata reports only the standard dependencies
\texttt{propext}, \texttt{Classical.choice}, and \texttt{Quot.sound} for
the final results. The accompanying verification ledger distinguishes
that upstream report from any fresh build and independent Comparator
execution performed in this extension \citep{NSLean2026}.

In this review the exact compiler distribution was downloaded and its
published SHA-256 checksum matched. Dependency cache retrieval then failed
with HTTP 429 responses after partial decompression. A direct attempt to
elaborate the final module stopped at missing local import artifacts.
Consequently this pass did not complete a fresh build of the final upstream
theorems, and did not run the independent Comparator certificate check.
The source audit found no challenge import and no unproved-placeholder token
in the 609-module local solution closure, after removing comments and string
literals. That lexical result has the narrower meaning stated here; it is
not a substitute for the uncompleted kernel verification.

The extension's new contribution is an explicit mathematical organization:
quadratic residual calculus; exact information-loss identities; norm
scaling that separates energy from regularity; a forbidden-word estimate
that explains a common analytic domain; all-order realization with
correct quantifiers; potential-based localization; and a quantitative
residual comparison for learned fields. Chapter~\ref{ch:fluid-structures}
then develops the arithmetic, exterior-power, and filtered-algebra
connections, while Chapter~\ref{ch:quantum-fluid} treats quantum
representations and learning.

These are rigorous bridges because each has a specified map or estimate.
They do not identify a finite surface-code complex with the continuum
Navier--Stokes equation, or turn a learned function into an existence
theorem without a stability certificate. They also preserve the released
theorem's force. Removing that force, imposing a different physical
constitutive law, or deriving the same singular flow from a microscopic
quantum state would be a further theorem with additional analytic content.

\chapter{Orders, complexes and the algebra of a fluid singularity}
\label{ch:fluid-structures}
\reviewstatus{This chapter gives original, explicit proofs of the finite algebraic and differential identities used to interpret the released construction. The construction of the singular Navier--Stokes solution remains the external theorem of \citet{NSOpenAI2026}, reviewed in Chapter~\ref{ch:navier-stokes-release}. The real auxiliary torus, the physical three-dimensional fluid, a Riemann surface and a quantum register are different objects. Every identification below specifies which structure it preserves.}

\section{What a useful reformulation must preserve}

The central opportunity is to reorganize the proof into mathematical interfaces that the inherited project already studies: integral endomorphisms, exterior powers, complexes, quadratic tensor maps, moment determinants and filtered operator algebras. This is more informative than attaching algebraic terminology to a picture of a vortex. It also creates obligations. A change of notation must preserve the domain, positive viscosity, zero initial data, compact smooth forcing, bounded kinetic energy and the claimed divergence of velocity. An algebraic identity must not silently replace an analytic estimate that is uniform in the concentrating scale.

In the released proof the physical equation is on $\R^3$, with a subsequent periodic construction on $\mathbb T^3$. The auxiliary $\mathbb T^2$ carries oscillatory labels and is evaluated along a prescribed phase map. It is neither the physical domain nor a quantum phase space. A surface-algebra model can encode incidence, permitted paths and local tensor types. To obtain an actual fluid discretization one must additionally choose a metric, a volume form, boundary conditions, constitutive parameters and approximation operators. Those data will remain visible throughout this chapter.

\section{The de Rham complex behind incompressible motion}

Let $M$ be an oriented flat three-dimensional manifold without boundary, or use fields on $\R^3$ for which the integrations below have no boundary terms. Write $\alpha=u^\flat$ for the velocity one-form, $\delta$ for the metric codifferential and $\Delta_H=d\delta+\delta d$ for the nonnegative Hodge Laplacian. With the usual Euclidean sign convention, $\Delta_H$ acts componentwise as minus the coordinate Laplacian. Thus incompressibility is $\delta\alpha=0$.

\begin{proposition}[The differential-form momentum equation]
\label{prop:fluid-forms}
For a smooth divergence-free velocity, the usual forced Navier--Stokes equation is equivalent to
\[
 \partial_t\alpha+\iota_u d\alpha+\nu\Delta_H\alpha+d\pi=f^\flat,
 \qquad \delta\alpha=0,\qquad \pi=p+\tfrac12|u|^2.
\]
Its vorticity two-form $\omega=d\alpha$ satisfies
\[
 \partial_t\omega+\mathcal L_u\omega+\nu\Delta_H\omega=df^\flat.
\]
\end{proposition}
\begin{proof}
The $i$th component of $\iota_u d\alpha$ is
$\sum_j u_j(\partial_j u_i-\partial_i u_j)
=(u\cdot\nabla)u_i-\partial_i(|u|^2/2)$.
Adding $d\pi$ restores the advective and pressure terms. The sign of $\Delta_H$ gives $-\nu\Delta u$ on the left. Apply $d$ to this equation. It commutes with $\partial_t$ and with $d\delta+\delta d$ because $d^2=0$. Cartan's identity gives
$d\iota_u\omega=\mathcal L_u\omega-\iota_u d\omega=\mathcal L_u\omega$.
The pressure disappears by $d^2\pi=0$.
\end{proof}

This places exterior powers in their proper role. The momentum flux $u\otimes u$ is symmetric; vorticity is alternating. The two occupy $\operatorname{Sym}^2$ and $\bigwedge^2$ tensor types, respectively. They interact through differential operators and metric contractions. Exchanging their symmetry types would change the equation.

\begin{proposition}[Energy and the work of a differential]
\label{prop:fluid-energy}
Under the hypotheses just stated,
\[
 \frac{d}{dt}\frac12\|u\|_2^2+\nu\|\nabla u\|_2^2
 =\langle f,u\rangle_{L^2}.
\]
In particular neither pressure nor transport contributes to the total kinetic-energy derivative.
\end{proposition}
\begin{proof}
The pressure contribution is $\langle d p,\alpha\rangle=\langle p,\delta\alpha\rangle=0$. Pointwise, $\langle\iota_u d\alpha,\alpha\rangle=d\alpha(u,u)=0$, and $\langle d(|u|^2/2),\alpha\rangle$ integrates to zero. Integration by parts gives $\langle-\Delta u,u\rangle=\|\nabla u\|_2^2$. Finally $\langle\partial_tu,u\rangle=\tfrac12\partial_t\|u\|_2^2$.
\end{proof}

The proposition controls an integral. It does not bound $\|u\|_\infty$. This matters especially in the newly released construction: concentration makes the energetic cost of a large speed small. The scale computation in Chapter~\ref{ch:navier-stokes-release} is therefore compatible with this identity, rather than a violation of it. Nor does dissipation imply that every individual mode decreases: transport can redistribute energy between scales while its total work vanishes.

On a curved manifold one must specify which viscous operator is intended. The Hodge Laplacian, the rough Laplacian and the divergence of a symmetric strain tensor differ by curvature contributions under standard conventions. A topological surface model cannot choose between them. The flat torus examples below fix the Hodge convention explicitly. Their familiar three-dimensional torus drawing is a visualization of periodic coordinates, not an assertion that the induced metric of the embedded doughnut is flat.

\section{A quadratic order occurs in the actual construction}

Section~6 of \citet{NSOpenAI2026} chooses the integer matrix
\[
 J=\begin{pmatrix}3&1\\1&5\end{pmatrix},\qquad
 A=J-4I=\begin{pmatrix}-1&1\\1&1\end{pmatrix}.
\]
Its arithmetic controls two different derivative scales on an auxiliary torus. This is a direct meeting point with the orders developed earlier in the monograph.

\begin{proposition}[The integral quadratic-order representation]
\label{prop:fluid-order}
One has $A^2=2I$, $J^2-8J+14I=0$ and $\det J=14$. The map
\[
 \rho:\Z[\sqrt2]\longrightarrow M_2(\Z),\qquad
 a+b\sqrt2\longmapsto aI+bA
\]
is an injective unital ring homomorphism, and
\[
 \det\rho(a+b\sqrt2)=a^2-2b^2.
\]
Consequently $J=\rho(4+\sqrt2)$ and $\det(J^k)=14^k$ for every integer $k\ge0$.
\end{proposition}
\begin{proof}
Matrix multiplication gives $A^2=2I$. Hence
$(aI+bA)(cI+dA)=(ac+2bd)I+(ad+bc)A$, the multiplication law of the quadratic order. Its matrix is
$\left(\begin{smallmatrix}a-b&b\\b&a+b\end{smallmatrix}\right)$, whose determinant is $a^2-2b^2$. If it is the zero matrix its off-diagonal entry gives $b=0$, then $a=0$, proving injectivity. Expanding $(J-4I)^2=2I$ gives the polynomial identity. Determinants are multiplicative, including the case $k=0$.
\end{proof}

The determinant is also the scalar by which the top exterior power acts:
$\bigwedge^2J=14\operatorname{id}_{\bigwedge^2\Z^2}$. Thus the same alternating-tensor construction that gave the determinant and Cauchy--Binet earlier now measures a covering multiplicity. Since $14$ is not a unit of $\Z$, this is a covering endomorphism, not an invertible integral change of basis.

\begin{proposition}[Covering multiplicity and exact averaging]
\label{prop:fluid-covering}
The map $\pi_k:\R^2/\Z^2\to\R^2/\Z^2$ given by $[y]\mapsto[J^ky]$ is a surjective covering with $14^k$ points in each fiber. For normalized Haar measure,
\[
 \int_{\mathbb T^2}F(J^ky)\,dy=\int_{\mathbb T^2}F(y)\,dy
\]
for every continuous $F$.
\end{proposition}
\begin{proof}
Invertibility over $\R$ gives surjectivity. A fiber is a translate of
$(J^{-k}\Z^2)/\Z^2$, isomorphic to $\Z^2/J^k\Z^2$. Integer row and column operations reduce the latter quotient to Smith normal form, whose order is $|\det J^k|=14^k$. Locally the real linear map is a diffeomorphism, proving the covering assertion. For a character $e_n(y)=e^{2\pi i n\cdot y}$, pullback gives $e_{(J^k)^Tn}$. Because $J^k$ is invertible, this character has nonzero integral exactly when $n=0$. Finite linear combinations of characters satisfy the integral identity. Uniform approximation by trigonometric polynomials extends it to continuous functions.
\end{proof}

This argument does not invoke a long-time ergodic limit. It is an exact identity for a finite covering. In the released proof one uses a common auxiliary torus for overlapping dyadic bands, keeping the difference of covering levels bounded. Its derivative factors are then uniformly bounded even though the absolute covering level grows. Omitting this distinction would lose the uniformity needed later in the estimates.

\begin{proposition}[Why this is not a complex elliptic-curve endomorphism]
\label{prop:fluid-no-complex}
There is no real linear complex structure $C$ on $\R^2$, meaning $C^2=-I$, that commutes with $J$.
\end{proposition}
\begin{proof}
The real symmetric matrix $J$ has distinct real eigenvalues $4-\sqrt2$ and $4+\sqrt2$. Every commuting map preserves each of its one-dimensional real eigenspaces. On each, $C$ is multiplication by a real scalar $c$, so $C^2$ is multiplication by $c^2\ge0$ and cannot be $-1$.
\end{proof}

Therefore the auxiliary covering has an arithmetic description by a quadratic order, but cannot be advertised as a holomorphic endomorphism of one fixed complex elliptic curve. The fact that its underlying real manifold is a torus supplies too little structure for that conclusion. This distinction also prevents an unsupported identification with the elliptic-curve groups used by cryptography.

\section{The number-field norm supplies an analytic estimate}

Set $b=\sqrt2-1$, $v_r=(1,-b)$ and $v_t=(b,1)$. Direct calculation gives
$Jv_r=(4-\sqrt2)v_r$ and $Jv_t=(4+\sqrt2)v_t$. Their irrational directions are useful because no nonzero integer frequency is perpendicular to either one, and the failure of perpendicularity has a quantitative lower bound.

\begin{lemma}[A quadratic small-divisor bound]
\label{lem:fluid-small-divisors}
There is an explicit constant $c=(1+(1+\sqrt2)^2)^{-1/2}>0$ such that
\[
 |v_r\cdot n|\ge c/|n|_2,\qquad |v_t\cdot n|\ge c/|n|_2
 \quad(n\in\Z^2\setminus\{0\}).
\]
\end{lemma}
\begin{proof}
For the first expression put $a=n_1+n_2$ and $d=-n_2$. Then $v_r\cdot n=a+d\sqrt2$. Its algebraic norm $a^2-2d^2$ is a nonzero integer: otherwise irrationality of $\sqrt2$ would force $a=d=0$, and hence $n=0$. Its absolute value is at least one. The conjugate is $n_1+(1+\sqrt2)n_2$, of absolute value at most $c^{-1}|n|_2$ by Cauchy--Schwarz. Dividing the norm by this bound proves the first assertion. The other direction is $(n_2-n_1)+n_1\sqrt2$, whose conjugate has the same coefficient norm.
\end{proof}

\begin{proposition}[A tame inverse of a torus directional derivative]
\label{prop:fluid-fourier-inverse}
Let $v$ be either direction above, and let $F$ be a smooth zero-mean function on $\mathbb T^2$. There is a unique smooth zero-mean $\phi$ satisfying $v\cdot\nabla\phi=F$, namely
\[
 \phi(y)=\sum_{n\ne0}\frac{\widehat F(n)}{2\pi i\,v\cdot n}e^{2\pi i n\cdot y}.
\]
For every real $s$ there is $C_s$ with
$\|\phi\|_{H^s}\le C_s\|F\|_{H^{s+1}}$. For every nonnegative integer $m$, one also has $\|\phi\|_{C^m}\le C_m\|F\|_{C^{m+4}}$.
\end{proposition}
\begin{proof}
The lemma bounds the inverse multiplier by $C|n|$. Squaring and summing the weighted Fourier coefficients proves the Sobolev estimate. For the uniform derivative estimate, integration by parts bounds a Fourier coefficient of a $C^{m+4}$ function by $C(1+|n|)^{-m-4}\|F\|_{C^{m+4}}$. Multiplying by the $m$ derivative factors and the inverse divisor leaves $C(1+|n|)^{-3}$, summable in two dimensions. This proves uniform convergence of all required derivative series. Applying the same argument for arbitrarily large $m$ proves smoothness. Uniqueness follows because every nonzero Fourier coefficient of a solution of the homogeneous equation vanishes, and the zero coefficient was prescribed.
\end{proof}

Here an integral norm identity yields a finite derivative loss, which is then paid for using smoothness at every fixed order. That is a genuine analytic contribution of the arithmetic structure. It is not an estimate uniform in derivative order, and it would not justify infinite-order flatness from a fixed finite regularity assumption. In Section~8 of the released proof, repeated integration by parts along such directions makes compact-support cutoff errors smaller than any specified power. The number of input derivatives increases with the requested power.

\section{Symmetric tensors, covariance and the stress correction}

For a real velocity perturbation $w$ averaged by a positive probability measure, $\langle w\otimes w\rangle$ is positive semidefinite: for every $a$ its quadratic form is $\langle(a\cdot w)^2\rangle\ge0$. It belongs to the symmetric tensor representation. A single $w\otimes w$ has rank at most one, so all its two-by-two minors vanish. The determinantal variety describes this algebraic rank constraint, but averaging rank-one matrices can produce a full-rank matrix. Rank should not be transported through an average without proof.

The released construction needs two specified radial cross-flux components rather than an arbitrary covariance matrix. If $b_+$ and $b_-$ have disjoint auxiliary supports, their mixed products vanish pointwise. Writing
\[
 \mathcal C(w)=\left\langle w_r(w_\theta,w_z)\right\rangle,
 \qquad H=[\mathcal C(b_+)\mid\mathcal C(b_-)],
\]
gives $\mathcal C(a_+b_++a_-b_-)=H(a_+^2,a_-^2)^T$. A cross-flux vector itself is not a positive semidefinite matrix; positivity concerns the coefficients of this two-ray representation.

\begin{proposition}[Positive realization and a signed tangent inverse]
\label{prop:fluid-covariance}
Let $H\in\GL_2(\R)$, let $s\in\R^2$ satisfy $y=H^{-1}s\in(0,\infty)^2$, and put $a_i=\sqrt{y_i}$. For $C(a)=H(a_1^2,a_2^2)^T$, define
\[
 L_a\sigma=\operatorname{diag}(2a_1,2a_2)^{-1}H^{-1}\sigma.
\]
Then $C(a)=s$, $D C(a)L_a\sigma=\sigma$, and
\[
 C(a+L_a\sigma)=s+\sigma+H\big((L_a\sigma)_1^2,(L_a\sigma)_2^2\big)^T.
\]
The last term is generally nonzero.
\end{proposition}
\begin{proof}
The first identity substitutes $a_i^2=y_i$. Differentiating the coordinate squares gives $D C(a)=H\operatorname{diag}(2a_1,2a_2)$, whose displayed right inverse exists since both amplitudes are positive. Expanding $(a_i+h_i)^2=a_i^2+2a_ih_i+h_i^2$ proves the final identity.
\end{proof}

This is the same distinction between tangent information and an exact finite deformation that appeared in the varieties-of-complexes chapter. The determinant of $H$ proves a local algebraic inverse exists. The PDE proof additionally needs a lower bound for that determinant, control of derivatives of $H^{-1}$, and flat edge factors that permit division by the leading amplitudes. It keeps those amplitudes fixed while applying signed increments and accounts for the quadratic remainder at each subsequent step. Replacing this procedure by the claim that every target stress is a positive covariance would be false.

The Schur-functor viewpoint is useful here without overreach. In characteristic zero, the flip on $V\otimes V$ yields the complementary projectors $(I+\mathrm{flip})/2$ and $(I-\mathrm{flip})/2$ onto symmetric and alternating tensors. The velocity covariance uses the first; the differential two-form uses the second. Higher tensor lifts use Young symmetries to remove redundant coordinates. Integral straightening theory controls the corresponding bases before denominators are inverted. None of these representation decompositions by itself supplies the scale estimates of the analytic construction.

\section{Moment determinants and exact compatibility}

Compactly supported primitives have obstructions. If $(\partial_r+e/r)\phi=g$ with $\phi$ compactly supported away from zero, multiplication by $r^e$ and integration gives $\int r^e g\,dr=0$. Conversely, for real $e$ and $g\in C_c^\infty((0,\infty))$ satisfying that moment condition,
$\phi(r)=r^{-e}\int_0^r s^e g(s)\,ds$ is compactly supported and solves the equation. This elementary exactness criterion is the prototype for the weighted radial inverses in the proof. Torus-dependent coefficients require the additional phase shifts and flat remainders described there.

\begin{proposition}[Independence of ordered moment functionals]
\label{prop:fluid-moments}
Let $\alpha_1,\ldots,\alpha_m$ be distinct real numbers, and let $\beta_j$ be nonzero, nonnegative smooth bumps supported inside disjoint ordered intervals $0<I_1<\cdots<I_m$. Then
\[
 B_{ij}=\int_{I_j}x^{\alpha_i}\beta_j(x)\,dx
\]
defines an invertible matrix. Smooth variation over a compact parameter set gives bounded inverse derivatives provided all the hypotheses, including distinctness of the exponents, persist there.
\end{proposition}
\begin{proof}
A nonzero linear combination of $m$ distinct powers has at most $m-1$ distinct positive zeros. For induction divide by the smallest power. If the resulting expression is a nonzero constant, it has no zeros. Otherwise its nonzero derivative is a combination of at most $m-1$ distinct powers and has at most $m-2$ zeros by the induction hypothesis; Rolle's theorem then proves the assertion. Hence $\det[x_j^{\alpha_i}]$ never vanishes for $0<x_1<\cdots<x_m$. It has constant sign on that connected chamber. Multilinearity and Fubini give
\[
 \det B=\int_{I_1\times\cdots\times I_m}
 \det[x_j^{\alpha_i}]\prod_j\beta_j(x_j)\,dx_1\cdots dx_m.
\]
The integrand has a fixed sign and is nonzero on a set of positive measure. Thus the determinant cannot vanish. Compactness bounds it away from zero for a compact family satisfying the hypotheses. The adjugate formula and repeated differentiation give the inverse-derivative bounds.
\end{proof}

For geometrically rescaled copies of one bump, the matrix becomes a Vandermonde matrix after row normalization. This connects the proof to the same alternating determinants and interpolation combinatorics used earlier. The quantitative limitation is crucial: if two exponents coalesce, the inverse can diverge. The parameter called $\lambda$ in the source is fixed before the high radial frequency is chosen. Sending it to zero inside an argument that needs a bounded inverse would alter the proof.

\begin{lemma}[A controlled quadratic moment correction]
\label{lem:fluid-quadratic-moments}
Suppose $B$ is invertible, $Q$ is bilinear, $\|B^{-1}\|\le\beta$, $\|Q\|\le\kappa$, and $\|d\|\le d_0$. If $8\beta^2\kappa d_0\le1$, the equation $Bc+Q(c,c)=d$ has a unique solution in the closed ball $\|c\|\le2\beta d_0$.
\end{lemma}
\begin{proof}
Set $r=2\beta d_0$ and $T(c)=B^{-1}(d-Q(c,c))$. For $\|c\|\le r$,
$\|T(c)\|\le\beta d_0+\beta\kappa r^2\le3r/4\le r$.
For two points in the ball, bilinearity gives
$\|T(c)-T(c')\|\le2\beta\kappa r\|c-c'\|\le\tfrac12\|c-c'\|$.
The contraction theorem gives the unique fixed point in that ball. The fixed-point equation is exactly the required moment equation.
\end{proof}

The estimate is a local solvability statement. It does not establish that a whole family of discrepancies is small; that must be proved from the profile construction and its order of parameter choices. This is the appropriate division of labor between algebraic elimination and analysis.

\section{From varieties of complexes to compatible fluid discretization}

Let a finite oriented cell decomposition of a closed surface have real cochains $C^0,C^1,C^2$ and coboundary matrices
$d_0:C^0\to C^1$, $d_1:C^1\to C^2$ with $d_1d_0=0$. The matrix entries are signed incidence numbers. These matrices give a point of the affine variety of complexes defined by the entries of $d_1d_0$. Additional rank conditions select the strata studied earlier. The underlying incidence relation is exact over the integers; a floating-point residual is only a computational check of its implementation.

Choose symmetric positive definite mass matrices $M_0,M_1,M_2$. Define
\[
 \delta_1=M_0^{-1}d_0^TM_1,\qquad
 \delta_2=M_1^{-1}d_1^TM_2,\qquad
 L_1=d_0\delta_1+\delta_2d_1.
\]
These are the chosen discrete adjoints and Hodge operator. Geometry enters through the $M_j$, not through the relation $d_1d_0=0$ alone. Stable approximation requires additional commuting projection and norm estimates; finite element exterior calculus makes these requirements precise \cite{NSFEEC2006,NSFEEC2010}.

\begin{proposition}[Discrete Hodge energy and incompressible subspace]
\label{prop:fluid-discrete-hodge}
For every $x\in C^1$,
\[
 x^TM_1L_1x=\|\delta_1x\|_{M_0}^2+\|d_1x\|_{M_2}^2\ge0.
\]
Moreover $L_1$ preserves $\ker\delta_1$. A velocity of the form $\delta_2\psi$ belongs to this kernel.
\end{proposition}
\begin{proof}
The definitions identify $d_0$ and $\delta_1$, and $d_1$ and $\delta_2$, as adjoints in the indicated inner products. Move each adjoint across the pairing to obtain the sum of squares. Taking adjoints of $d_1d_0=0$ gives $\delta_1\delta_2=0$. Therefore
$\delta_1L_1=(\delta_1d_0)\delta_1$, which is zero on $\ker\delta_1$. The last assertion is the same zero-composition identity.
\end{proof}

An arbitrary algebraic change of basis preserves the complex equation. It preserves the numerical energy only if the metric is transported as well. If $x'=Tx$, then the corresponding matrix is $M'_1=T^{-T}M_1T^{-1}$. Retaining the identity metric after a badly conditioned change of basis can change coercivity constants arbitrarily. This repeats, in the fluid setting, the distinction between homological equivalence and quantitative spectral control.

\begin{theorem}[An energy-certified finite learned evolution]
\label{thm:fluid-finite-closure}
Work in orthonormal coordinates on a finite-dimensional divergence-free subspace. Let $K(a)$ be a locally Lipschitz skew-symmetric matrix, $L$ a symmetric positive semidefinite matrix, and $B_\theta(a)$ a locally Lipschitz matrix. For continuous forcing $g(t)$, consider
\[
 \dot a=-K(a)a-\nu La-B_\theta(a)^TB_\theta(a)a+g(t),\qquad\nu>0.
\]
It satisfies
\[
 \tfrac12\frac d{dt}|a|^2
 =-\nu a^TLa-|B_\theta(a)a|^2+a\cdot g(t).
\]
Every solution exists for all finite positive times, provided the displayed coefficients are defined and locally Lipschitz on the whole finite-dimensional space.
\end{theorem}
\begin{proof}
For a skew matrix $a^TK(a)a$ equals its own negative, hence is zero. The factorized term has work $|B_\theta(a)a|^2$. This proves the identity. Young's inequality yields
$\frac d{dt}|a|^2\le |a|^2+|g(t)|^2$.
Gronwall gives a finite bound on every finite interval because $g$ is continuous. A locally Lipschitz finite-dimensional ODE can be continued while its solution stays in a compact subset of its domain. The bound therefore rules out finite-time escape and proves global continuation.
\end{proof}

The theorem is a useful design rule for a learned closure, including one whose matrix entries are produced by quantum expectation values. It is also a useful warning: every fixed finite model can be globally well behaved while the continuum family develops unbounded pointwise velocity. The equivalence constant between a grid energy norm and the maximum grid value typically diverges under refinement. Global continuation of each finite ODE cannot be passed to the continuum without uniform regularity and consistency estimates.

The dissipative ansatz is deliberately restrictive. True interscale transfer can include backscatter, while $-B_\theta^TB_\theta$ never injects resolved energy. A trained model in this class must be assessed for that modeling bias. The algebraic certificate proves stability of its own evolution, not fidelity to the exact coarse dynamics.

\section{Gentle algebras, path relations and quantitative interfaces}

The project uses gentle surface algebras to organize permitted local paths and relations. A compatible use here is to let vertices label local field types or mesh patches, arrows label permitted transfers, and quadratic relations encode compositions known to vanish. The equation $d^2=0$ is one such zero composition. The derivative-block relation in the Volterra construction, discussed in Chapter~\ref{ch:navier-stokes-release}, is another actual example. Counting admissible words then bounds how often a derivative-losing operator can appear.

One must verify the gentle conditions for the chosen quiver and relations. A free algebra with a single relation $b^2=0$ is not automatically gentle: allowed continuations at its other generator can violate the required local uniqueness. Nor does an arbitrary triangulated surface incidence quiver become gentle merely because it is drawn on a surface. The valid general statement is that a representation of a specified path algebra with relations encodes the corresponding exact operator identities. Calling it gentle requires the extra combinatorial checks already described in the earlier chapters.

An associated graded construction is likewise conditional on a specified filtration. If classes $\mathcal M_\alpha$ obey
$\mathcal M_\alpha\mathcal M_\beta\subseteq\mathcal M_{\alpha+\beta}$,
they track orders under multiplication. If a derivative maps $\mathcal M_\alpha$ into $\mathcal M_{\alpha-\kappa}$, its loss must be included in a correction cycle. Logarithmic factors and flat shell weights are part of the analytic definition. A finitely generated semigroup algebra, a flat toric degeneration, normality or a Schur-module decomposition does not follow solely from this notation. Such claims require the corresponding algebraic hypotheses and proofs.

Standard monomial theory is valuable when one introduces an actual finite coordinate ring of rank constraints, tensor symmetries or moment minors. It provides bases, relations and algebraic elimination. The Navier--Stokes construction further needs positivity margins, compact support and all-order differentiability. These conditions are not encoded by an affine scheme over $\C$. In particular positivity is real semialgebraic or analytic data, and flat nonzero functions cannot be recovered from a convergent analytic Taylor series at their flat endpoint.

\section{Determinants, spectra and what the reformulation adds}

For a finite positive operator $L$ with eigenvalues $\lambda_j$, the identities
\[
 \det(sI+L)=\prod_j(s+\lambda_j),\qquad
 \frac d{ds}\log\det(sI+L)=\operatorname{tr}(sI+L)^{-1}
\]
hold for $s>0$ by diagonalization and ordinary differentiation. Exterior powers and Schur constructions organize coefficients of the determinant. A continuum determinant requires a specified regularization and domain; a spectral zeta function requires convergence or analytic continuation. No identification with an arithmetic $L$-function, its zeros, or the inherited GRH program follows from these finite identities.

The improvement supplied by this extension is therefore explicit: the integral torus covering is expressed as an order action; its norm becomes a small-divisor estimate; covariance realization becomes a quadratic map with a retained remainder; moment matching becomes a determinant and contraction problem; incompressibility becomes a complex relation; and learned finite models receive a verifiable energy interface. The released proof's analytic construction is preserved rather than replaced. A stronger theorem, such as a result for unforced Navier--Stokes or a quantitative algorithm producing its full singular field, would require additional work beyond these reformulations.

\chapter{Quantum information, exact reduction and learned fluid models}
\label{ch:quantum-fluid}
\reviewstatus{The reduction, channel and scalar Carleman statements below are proved explicitly. The computational companion evaluates finite classical simulations of the stated models. Quantum algorithms cited here have their own input, error and readout assumptions. No quantum speedup, trained fluid ToricGT model, or numerical reproduction of the released singular solution is claimed.}

\section{Information reduction is a map, not an autonomous dynamics}

The microscopic-to-macroscopic hierarchy in Chapter~\ref{ch:microscopic-fluid-limits} is naturally expressed by maps that retain observables. Examples include taking a partial trace, a one-particle marginal, a velocity moment, an angular average, or a projection onto a finite set of modes. These maps can be exact even when the subsequent closed evolution is approximate. The distinction is elementary: knowing a projected state need not determine the derivative of that state.

\begin{proposition}[The obstruction to an exact instantaneous closure]
\label{prop:fluid-closure-obstruction}
Consider $\dot x=Ax+Bz$ on finite-dimensional vector spaces. If $By\ne0$ for some hidden vector $y$, no function $F(x)$ can satisfy $F(x)=Ax+Bz$ for every pair $(x,z)$.
\end{proposition}
\begin{proof}
For fixed $x$, compare the hidden states $z=0$ and $z=y$. The proposed equality gives both $F(x)=Ax$ and $F(x)=Ax+By$, contradicting $By\ne0$.
\end{proof}

The statement does not rule out a closure on a specified invariant manifold, a conditional expectation under a prescribed ensemble, or an approximation with controlled error. It rules out omitting those hypotheses. It applies equally to a classical neural network and a quantum learned function: an expressive predictor cannot make an insufficient input state mathematically sufficient.

\begin{theorem}[Exact memory from a linear hidden subsystem]
\label{thm:fluid-memory}
Let $A,B,C,D$ be constant compatible finite matrices and consider
\[
 \dot x=Ax+Bz,\qquad \dot z=Cx+Dz.
\]
Then
\[
 \dot x(t)=Ax(t)+Be^{tD}z(0)
       +\int_0^t Be^{(t-s)D}Cx(s)\,ds.
\]
Conversely, any continuously differentiable $x$ satisfying this equation, together with
$z(t)=e^{tD}z(0)+\int_0^t e^{(t-s)D}Cx(s)\,ds$, satisfies the original system.
\end{theorem}
\begin{proof}
Multiply the hidden equation by $e^{-tD}$ and differentiate the product. Integration gives the displayed expression for $z$. Substitution gives the equation for $x$. In the converse direction, differentiate the finite matrix-exponential integral using the fundamental theorem of calculus. This recovers $\dot z=Cx+Dz$, and substitution gives the other equation.
\end{proof}

This is a finite, fully explicit version of the memory mechanism associated with projection methods \cite{NSMori1965}. The memory kernel is $K(t)=Be^{tD}C$, and the hidden initial condition supplies a separate term. Dropping that term is not justified simply by calling it noise. Its distribution and correlations must be supplied if a stochastic description is intended.

For example, with scalar equations $\dot x=-2x+z$, $\dot z=x-3z$, the exact kernel is $e^{-3t}$. Even when $z(0)=0$, replacing the reduced equation by $\dot x=-2x$ discards a nonzero convolution. A fast-relaxation approximation can instead use $z\approx x/3$, but its error depends on the relative time scales and initial layer. Exact elimination and time-scale approximation are different steps.

\begin{proposition}[Discrete memory and the resolvent determinant]
\label{prop:fluid-discrete-memory}
For $x_{n+1}=Ax_n+Bz_n$, $z_{n+1}=Cx_n+Dz_n$,
\[
 z_n=D^nz_0+\sum_{j=0}^{n-1}D^{n-1-j}Cx_j,
\]
and substitution gives the exact resolved recurrence. If $sI-D$ is invertible, the determinant of the full block resolvent satisfies
\[
 \det\begin{pmatrix}sI-A&-B\\-C&sI-D\end{pmatrix}
 =\det(sI-D)\det\bigl(sI-A-B(sI-D)^{-1}C\bigr).
\]
\end{proposition}
\begin{proof}
The hidden-state formula follows by induction, with the sum empty at $n=0$. Multiplying by $D$ and adding $Cx_n$ advances the sum to $n+1$. For the determinant, multiply the block matrix on the left by
$\left(\begin{smallmatrix}I&B(sI-D)^{-1}\\0&I\end{smallmatrix}\right)$, whose determinant is one. The upper-right block becomes zero and the upper-left block becomes the displayed Schur complement. The determinant of the resulting block triangular matrix is the product of its diagonal-block determinants.
\end{proof}

This links the project's determinant calculus to an actual information-reduction operation. The frequency dependence of the Schur complement is the transform of a memory kernel. Replacing it by a constant is a modeling approximation. The existence of its inverse does not provide a uniform condition-number bound, especially near a spectral pole.

\section{Coarse-graining and the unresolved stress}

For a normalized nonnegative averaging kernel, write $\bar u=\langle u\rangle$ and
$R=\langle u\otimes u\rangle-\bar u\otimes\bar u$. The quadratic nonlinearity ensures that the equation for $\bar u$ contains a contribution involving $R$. A projection that retains only $\bar u$ has discarded information relevant to the momentum equation.

\begin{proposition}[Covariance as the exact quadratic defect]
\label{prop:fluid-filter-covariance}
For a probability average with finite second moments (at a fixed observation point, when the average is spatially local),
\[
 R=\langle(u-\bar u)\otimes(u-\bar u)\rangle\succeq0,
 \quad \operatorname{tr}R=\langle|u|^2\rangle-|\bar u|^2.
\]
If a linear filter commutes with spatial derivatives, filtering $\nabla\cdot(u\otimes u)$ gives
$\nabla\cdot(\bar u\otimes\bar u)+\nabla\cdot R$.
\end{proposition}
\begin{proof}
Expand the centered product, use $\langle u\rangle=\bar u$ and $\langle1\rangle=1$, and take traces. Testing the matrix against any vector gives an average of a square. The final identity follows by linearity, the definition of $R$, and the assumed derivative commutation.
\end{proof}

For a convolution evaluated at $x$, the centered integrand in this proof is $u(y)-\bar u(x)$, averaged over $y$. It is not $u(y)-\bar u(y)$ for a general non-idempotent filter.

A sharp Fourier truncation need not be a positive averaging kernel. Its quadratic defect is still defined, but the positive-semidefinite conclusion requires the probability-average hypothesis. Similarly, a filter near a wall need not commute with differentiation. Those are concrete places where otherwise plausible information-theoretic arguments fail if their maps are left unspecified.

In the released blowup proof, oscillatory perturbations generate selected averaged quadratic fluxes. The authors construct those fluxes and bound their errors; they do not infer them from an assumed generic turbulence closure. The covariance inverse in Proposition~\ref{prop:fluid-covariance} captures one finite algebraic component of that construction. Its analytic realization retains the carrier phase, exact curl correction, support separation and derivative estimates.

\section{A quantum observable can have the same closure obstruction}

For a finite quantum system with density matrix $\rho$ and Hermitian observable $O$, the retained datum $\operatorname{tr}(\rho O)$ is only one linear functional of the state. Even noiseless knowledge of it need not determine its future derivative.

\begin{example}[Identical observed value, opposite future derivative]
Take $\hbar=1$, Hamiltonian $H=X$ and observable $O=Z$, with Pauli matrices in the conventions of the earlier quantum chapters. The states
$\rho_\pm=(I\pm Y)/2$ both have $\operatorname{tr}(\rho_\pm Z)=0$. Unitary evolution gives
\[
 \frac d{dt}\langle Z\rangle=i\langle[X,Z]\rangle=2\langle Y\rangle,
\]
so their derivatives are $+2$ and $-2$.
\end{example}

The commutator is obtained by direct multiplication, $[X,Z]=-2iY$. Thus the example requires no interpretive claim about measurement disturbance. The retained statistic simply omits a degree of freedom that matters dynamically. A reduced density matrix can likewise acquire memory when a subsystem is coupled to an environment. A time-homogeneous Markov semigroup requires additional assumptions; it does not follow from partial trace alone.

The same point controls proposed quantum information explanations of viscosity. Closed unitary evolution preserves von Neumann entropy. Effective irreversible evolution arises after a specified reduction and limiting or modeling assumptions. A completely positive trace-preserving map provides a consistent finite open-system operation, but an arbitrary such map is not a derivation of the Navier--Stokes constitutive law.

\section{An explicit heat channel on a periodic surface grid}

Let the position register have orthonormal basis $\{|j,k\rangle: j,k\in\Z/L\Z\}$. Let $T_x,T_y$ be the unitary permutations that translate by one grid point. For $0\le r\le1/4$ define
\[
 \Phi_r(\rho)=(1-4r)\rho+r\bigl(T_x\rho T_x^\dagger+T_x^\dagger\rho T_x
                    +T_y\rho T_y^\dagger+T_y^\dagger\rho T_y\bigr).
\]

\begin{proposition}[A trace-preserving quantum diffusion step]
\label{prop:fluid-heat-channel}
The map $\Phi_r$ is completely positive, trace preserving and unital. Its diagonal probabilities obey
\[
 p'_{jk}=(1-4r)p_{jk}+r(p_{j-1,k}+p_{j+1,k}+p_{j,k-1}+p_{j,k+1}).
\]
They remain nonnegative, sum to one, and their Shannon entropy cannot decrease.
\end{proposition}
\begin{proof}
Use Kraus operators $\sqrt{1-4r}I$, $\sqrt r T_x$, $\sqrt r T_x^\dagger$, $\sqrt r T_y$, $\sqrt r T_y^\dagger$. Their adjoint products sum to $I$, as do the reversed products. This proves complete positivity, trace preservation and unitality. Taking diagonal entries gives the stencil. It is a convex combination of translations of the probability vector. The function $-s\log s$ is concave on $[0,1]$, with its continuous value zero at zero. Jensen's inequality at each site, followed by summation and translation invariance, proves the entropy assertion.
\end{proof}

\begin{proposition}[Discrete heat energy and eigenvalues]
\label{prop:fluid-heat-spectrum}
Let $L$ be a positive integer, $h>0$ and $\nu\ge0$. On the periodic $L\times L$ grid, define $\Delta_h p=h^{-2}(T_x+T_x^{-1}+T_y+T_y^{-1}-4I)p$. Then
\[
 -\langle p,\Delta_h p\rangle
 =h^{-2}\sum_{j,k}\bigl((p_{j+1,k}-p_{jk})^2+(p_{j,k+1}-p_{jk})^2\bigr).
\]
The Fourier mode $e^{2\pi i(aj+bk)/L}$ has eigenvalue
$-4h^{-2}[\sin^2(\pi a/L)+\sin^2(\pi b/L)]$.
The heat equation $\dot p=\nu\Delta_hp$ therefore contracts squared Euclidean norm and preserves the constant mode.
\end{proposition}
\begin{proof}
Expand each squared difference and reindex periodic sums to obtain the quadratic identity. Translation acts on the Fourier mode by $e^{\pm2\pi ia/L}$ or $e^{\pm2\pi ib/L}$. The formula follows from $2\cos\theta-2=-4\sin^2(\theta/2)$. Differentiating $\tfrac12\|p\|^2$ along the heat equation and using the first identity proves contraction. All shifts fix the constant vector.
\end{proof}

With $r=\nu\Delta t/h^2$, the channel gives an explicit heat step. Its positivity requires the displayed time-step bound. The continuous semigroup is the limit of these steps, equivalently the exponential of the random-translation generator. This is an exact finite quantum-channel interpretation of scalar diffusion. The browser computes it classically on a $16\times16$ grid. The position space has dimension $256$, representable by eight qubits, but the computation is not an eight-qubit hardware experiment or a solution of nonlinear Navier--Stokes. A signed velocity field is also not a probability distribution; any encoding of it must preserve signs and normalization by an additional construction.

\section{What the Madelung transformation does and does not derive}

On a simply connected open region where a smooth scalar wave function does not vanish, write $\psi=\sqrt\rho\,e^{iS/\hbar}$. For a particle of mass $m>0$ in a real potential $V$, substitute this into
$i\hbar\partial_t\psi=-(\hbar^2/2m)\Delta\psi+V\psi$.

\begin{proposition}[The local quantum hydrodynamic equations]
\label{prop:fluid-madelung}
With $v=\nabla S/m$ and $Q=-(\hbar^2/2m)\Delta\sqrt\rho/\sqrt\rho$, one obtains
\[
 \partial_t\rho+\nabla\cdot(\rho v)=0,
 \qquad
 \partial_tS+\frac{|\nabla S|^2}{2m}+V+Q=0.
\]
Consequently $m(\partial_tv+(v\cdot\nabla)v)=-\nabla(V+Q)$ and $\nabla\times v=0$ in that region.
\end{proposition}
\begin{proof}
Differentiate the amplitude-phase product twice. Its Laplacian is
\[
 e^{iS/\hbar}\left[\Delta\sqrt\rho+
 \frac{2i}{\hbar}\nabla\sqrt\rho\cdot\nabla S+
 \frac{i}{\hbar}\sqrt\rho\,\Delta S-
 \frac1{\hbar^2}\sqrt\rho\,|\nabla S|^2\right].
\]
Equate imaginary parts and multiply by $2\sqrt\rho$ to obtain continuity. Divide the real part by $\sqrt\rho$ to obtain the second equation. Its gradient gives the velocity equation because the Hessian of $S$ is symmetric, so $\nabla(|v|^2/2)=(v\cdot\nabla)v$. The curl of a gradient is zero.
\end{proof}

There is no viscous term in this calculation. At zeros of $\psi$ the amplitude-phase chart fails; around a closed loop with a single-valued nonzero wave function the phase winding can give circulation $2\pi\hbar n/m$. One scalar phase on a nonvanishing simply connected patch therefore cannot represent an arbitrary smooth vortical incompressible velocity. More general spinor, mixed-state or multi-component encodings may represent richer currents, but their dynamical correspondence must be separately derived. The Born--Oppenheimer and kinetic limits reviewed earlier are distinct routes, not consequences of the scalar Madelung identity alone.

\section{Carleman lifting and the symmetric algebra}

Consider a finite quadratic system $\dot x=Ax+B(x\otimes x)+c$, with $x\in\R^n$. Differentiating a tensor power gives
\[
 \frac d{dt}x^{\otimes m}
 =\sum_{j=1}^m x^{\otimes(j-1)}\otimes\dot x\otimes x^{\otimes(m-j)}.
\]
Substitution couples tensor degrees $m-1,m,m+1$. Thus the infinite collection of powers evolves linearly in a larger space. This is an identity along an existing solution, not a conclusion that the infinite system has a bounded solution for all time.

All tensors $x^{\otimes m}$ are symmetric. In commuting coordinates the degree-$m$ monomials form a basis of dimension $\binom{n+m-1}{m}$ rather than $n^m$. The symmetric algebra and its standard monomials therefore provide an exact reduction of redundant tensor coordinates. Differentiation acts as a derivation on this algebra. Exterior powers have a different role: they encode determinants and rank constraints, not powers of one commuting vector, whose repeated wedge is zero.

\begin{proposition}[An exact scalar truncation calculation]
\label{prop:fluid-carleman}
Let $N\ge1$ be an integer, $\alpha>0$, $\beta\ge0$, $x_0>0$, and consider
$\dot x=-\alpha x+\beta x^2$. Put $r=\beta x_0/\alpha$ and $z(t)=r(1-e^{-\alpha t})$. Before the denominator vanishes,
\[
 x(t)=\frac{x_0e^{-\alpha t}}{1-z(t)}.
\]
Truncate the lifted equations
$\dot y_m=-m\alpha y_m+m\beta y_{m+1}$ at $m=N$ by setting $y_{N+1}=0$ and $y_m(0)=x_0^m$. Their first component is
\[
 y_1^{(N)}(t)=x_0e^{-\alpha t}\sum_{j=0}^{N-1}z(t)^j,
 \qquad
 x(t)-y_1^{(N)}(t)=x_0e^{-\alpha t}\frac{z(t)^N}{1-z(t)}.
\]
If $r<1$, the error is bounded by $x_0r^N/(1-r)$ for all $t\ge0$. If $r>1$, the nonlinear solution blows up at $t_*=-\alpha^{-1}\log(1-1/r)$; the finite lifted system does not remove that singularity.
\end{proposition}
\begin{proof}
For the exact equation, $w=1/x$ satisfies $\dot w=\alpha w-\beta$. Solving this linear equation yields the displayed formula. The truncated triangular system is solved by downward variation of constants. More explicitly its $m$th component is
\[
 y_m^{(N)}(t)=x_0^m e^{-m\alpha t}
 \sum_{j=0}^{N-m}\binom{m+j-1}{j}z(t)^j.
\]
Differentiate this expression, use $z'=\beta x_0e^{-\alpha t}$ and
$j\binom{m+j-1}{j}=m\binom{m+j-1}{j-1}$, and check the initial values. Uniqueness for the finite linear system proves the formula. At $m=1$ it is the geometric partial sum. Subtracting from the exact rational expression gives the remainder. The bound follows from $0\le z(t)\le r<1$. When $r>1$, the denominator first vanishes at the stated positive time, while the numerator remains positive.
\end{proof}

This calculation is a small, rigorous model of the convergence issue faced by quantum linearization. The finite lift is globally solvable even when it no longer approximates the nonlinear solution. It follows that solving a truncated lifted system, classically or quantum mechanically, cannot establish continuum regularity without a truncation estimate valid in the relevant regime.

\section{What the quantum-algorithm literature actually supplies}

\citet{NSLiu2021} analyze a dissipative quadratic ODE algorithm with a parameter measuring nonlinear and forcing strength relative to linear dissipation. Their stated favorable regime requires $R<1$, together with sparse access, state preparation, error and solution-norm assumptions. The final output is a normalized quantum state encoding a vector; obtaining a full classical velocity field has an additional readout cost. Their numerical Burgers examples outside the sufficient regime are evidence about those examples, not a theorem covering every Reynolds number. The retrieved arXiv source is version~4, dated 7 July 2026. That correction adds $\|F_0\|\le\|F_2\|$ to the relevant convergence statements. It does not repair a separate eigenvalue-to-norm inference in the printed Lemma~1, which must not be imported into this project.

\begin{proposition}[A counterexample to eigenvalue-only norm decay]
\label{prop:fluid-nonnormal-counterexample}
For the real quadratic system $\dot u=F_1u+F_2(u\otimes u)+F_0$, choose
\[
 F_1=\begin{pmatrix}-1&4\\0&-2\end{pmatrix},\quad
 F_2=\begin{pmatrix}1/100&0&0&0\\0&0&0&0\end{pmatrix},\quad
 F_0=\begin{pmatrix}1/1000\\0\end{pmatrix},\quad
 u(0)=\begin{pmatrix}1/2\\1/2\end{pmatrix}.
\]
The matrix $F_1$ is diagonalizable with strictly negative eigenvalues. In the Euclidean and induced spectral norms, the source's parameter is $R=3\sqrt2/500<1$, and $\|F_0\|\le\|F_2\|$. Nevertheless
\[
 \left.\frac d{dt}\|u(t)\|^2\right|_{t=0}=\frac{1007}{2000}>0.
\]
In particular these hypotheses do not imply $\|u(t)\|<\|u(0)\|$ for every positive sufficiently small $t$.
\end{proposition}
\begin{proof}
The distinct eigenvalues are $-1,-2$, the three relevant norms are
$\|u(0)\|=1/\sqrt2$, $\|F_2\|=1/100$, $\|F_0\|=1/1000$,
and substitution gives the stated $R$. Direct evaluation gives
$\dot u(0)=(3007/2000,-1)^T$. Therefore
$2u(0)^T\dot u(0)=1007/2000$. The bound asserted by the source's equation (4.2), with spectral real part $-1$, would instead be
$-1+3\sqrt2/500<0$. The exact positive derivative contradicts that bound and yields an interval of initial norm increase by differentiability.
\end{proof}

The problem is that $u^*(F_1+F_1^*)u$ is controlled by the largest eigenvalue of the Hermitian part, not by twice the largest real part of the eigenvalues of an arbitrary nonnormal matrix. The source specifies Euclidean and spectral norms and does not first transform the system into an eigenbasis. Its scalar rescaling leaves $F_1$ unchanged. The counterexample concerns that particular printed inference; it is not a claim that every application or a suitably repaired quantum algorithm fails.

\begin{proposition}[A sufficient dissipative norm hypothesis]
\label{prop:fluid-dissipative-ball}
Suppose $(F_1+F_1^*)/2\preceq-\alpha I$ with $\alpha>0$, $\|F_2\|\le\beta$, and $\|F_0(t)\|\le\gamma$, with continuous forcing. If a radius $a>0$ satisfies $\beta a^2-\alpha a+\gamma<0$, every solution starting in $\|u(0)\|\le a$ remains in that ball and continues for all positive times.
\end{proposition}
\begin{proof}
At a nonzero state of radius $r$, differentiation and Cauchy--Schwarz give
$\dot r\le-\alpha r+\beta r^2+\gamma$. At the boundary $r=a$ the right side is strictly negative. Continuity makes it negative on a sufficiently small annulus about that radius, so a trajectory cannot first cross that annulus outwards. A boundary initial state immediately moves inward by the same inequality. Thus the ball is invariant. The polynomial vector field is locally Lipschitz in the state, and continuity of the forcing on each finite time interval and boundedness of the trajectory give global finite-dimensional continuation.
\end{proof}

This repaired condition is stronger than spectral stability and is exactly the kind of additional metric information demanded throughout the monograph. A weighted inner product can sometimes establish such a condition, but its conditioning and the transformed nonlinear norm must then enter the quantitative and algorithmic bounds. We use only the independently proved estimates here, not the unrepaired general norm-decay lemma.

Non-dissipative extensions also need hypotheses. Theorem~1.1 of the reviewed version of \citet{NSWu2025} assumes a diagonalizable $F_1$, eigenvalues with nonpositive real parts, a uniform solution bound $\|x(t)\|_1\le\mu$, and the strictly positive resonance gap
\[
 \Delta=\inf_{k,\,m\in\N^n,\,|m|\ge2}
       \left|\lambda_k-\sum_jm_j\lambda_j\right|>0.
\]
Here $\N$ includes zero. Its stated estimate uses
$R_r=4e\mu\kappa_1(W)\|F_2\|_1/\Delta$; geometric convergence requires $R_r<1$. Merely asserting that each individual resonance expression is nonzero does not establish a positive infimum over infinitely many expressions. The condition also has a particularly restrictive consequence for real systems.

\begin{proposition}[Real nonresonance excludes neutral oscillatory modes]
\label{prop:fluid-real-nonresonance}
Let $F_1$ be a real matrix, with all eigenvalues in the closed left half-plane. If no eigenvalue equals a sum of at least two eigenvalues, allowing repetition, then every eigenvalue has strictly negative real part.
\end{proposition}
\begin{proof}
A zero eigenvalue violates the hypothesis because $0=0+0$. A nonzero purely imaginary eigenvalue $i\omega$ of a real matrix has the conjugate eigenvalue $-i\omega$. Then $i\omega=2i\omega+(-i\omega)$ is a sum with three terms, again violating the hypothesis. These are all possibilities on the imaginary axis, so the nonpositive real parts are strictly negative.
\end{proof}

Consequently, the title of that paper cannot be used to assert a theorem for arbitrary real conservative or neutral fluid systems. Its numerical examples and its precise theorem assumptions must be assessed separately. This observation does not refute its nonresonant estimate; it identifies the real-system domain in which that estimate could apply. We do not use the source's stronger estimate as a dependency of the released singularity proof. A direct estimate adequate for the present finite comparison is as follows.

\begin{proposition}[A norm-controlled short-time Carleman estimate]
\label{prop:fluid-carleman-contractivity}
Let $N\ge1$ be an integer. Suppose $\dot x=F_1x+F_2x^{\otimes2}$ on $[0,T]$, with $\|x(t)\|\le M$. Use compatible Euclidean tensor norms or compatible $\ell^1$ tensor norms, and assume $\|e^{tF_1}\|\le1$ for $t\ge0$. For the degree-$N$ Carleman truncation initialized by $y_j(0)=x(0)^{\otimes j}$,
\[
 \|x(t)-y_1(t)\|\le M\bigl(M\|F_2\|t\bigr)^N,
 \qquad 0\le t\le T.
\]
In particular the right side converges geometrically in $N$ when $M\|F_2\|T<1$.
\end{proposition}
\begin{proof}
The diagonal degree-$j$ generator is
$A_{j,j}=\sum_{i=0}^{j-1}I^{\otimes i}\otimes F_1\otimes I^{\otimes(j-1-i)}$.
Its semigroup is $(e^{tF_1})^{\otimes j}$ and is contractive under the stated norms. The adjacent block is the analogous sum with $F_2$ replacing one input slot; its norm is at most $j\|F_2\|$. Write $\eta_j=x^{\otimes j}-y_j$. All initial errors vanish, and the sole terminal forcing is $A_{N,N+1}x^{\otimes(N+1)}$. Repeated variation of constants expresses $\eta_1$ as an $N$-fold ordered integral. The integrand norm is at most
$N!\|F_2\|^N M^{N+1}$. The ordered simplex has volume $t^N/N!$, giving the bound.
\end{proof}

Diagonalizability with nonpositive spectral real parts alone does not imply contractivity in the original norm. If $F_1=W\Lambda W^{-1}$, one can apply the proposition to $z=W^{-1}x$, with the transformed nonlinearity $\widetilde F_2=W^{-1}F_2(W\otimes W)$ and a bound on $z$. Returning to $x$ costs $\|W\|$. This makes the conditioning visible instead of silently discarding it. Semisimplicity of neutral eigenspaces also differs from claiming that every eigenvalue has geometric multiplicity one. These distinctions correct misleading shortcuts in the preliminary estimates of the reviewed preprint without importing them into our argument.

For our project the required interface is a bound for the particular lifted operator and initial state, on the requested time interval, with its dependence on discretization made explicit.

The September 2026 end-to-end study by \citet{NSJennings2026} is especially relevant. It accounts for Carleman truncation, time stepping, conditioning, unitary block construction and observable extraction in a lattice-Boltzmann approach. Its shifted formulation and discrete lifting address specific earlier bottlenecks, but its conclusions concern selected observables and parameter regimes. The paper reports numerical convergence and nonconvergence behavior under changing truncation order, and distinguishes proven bounds from fitted scaling. We retain that distinction. Neither its quantum encoding nor its finite Reynolds-number experiments are a proof of the released three-dimensional singularity theorem.

Two further approaches deserve separate names. Variational quantum methods can represent nonlinear functionals using several copies of a state and a chosen ansatz \cite{NSLubasch2020}. Their practical result depends on optimization, ansatz approximation, measurement variance and implementation costs. Tensor-network methods such as \citet{NSGourianov2022} can exploit interscale correlations on classical computers; their compression depends on the ranks needed by the particular flow. A tensor-network simulation is not a quantum-device execution. Both can be useful research tools without claiming a universal exponential speedup.

For every proposed fluid algorithm, the complete cost should specify at least the spatial approximation, time discretization, nonlinear lift or ansatz error, state preparation, success probability, conditioning and requested observable. Normalizing a large vector does not make its small physically relevant component free to read out. A statement polynomial in the number of qubits can still have a large or exponential dependence on physical accuracy, evolution time or a hidden condition number.

\section{Exact finite repairs for collision block encodings}
\label{sec:fluid-circuit-repairs}

The detailed algorithm of \citet{NSJennings2026} makes it possible to inspect the finite maps underlying its collision encoding. In the reviewed arXiv version~2, Appendix~D, equation~(D56), gives a purported unitary for an equality map, while Appendix~I, equation~(I65), gives a matrix dilation without its necessary contraction normalization. We supply explicit finite replacements. These proofs concern the two interfaces; they do not establish Carleman convergence, a total resource bound, or quantum advantage for the complete fluid algorithm.

Let $\mathcal H$ have orthonormal basis $\{|a\rangle:a\in\{0,1\}^n\}$, with $n\ge1$. Padding a position register to a power of two permits this notation. Define
\[
 S=\sum_a|a,0\rangle\langle a,a|,\qquad
 \Pi_{\mathrm{eq}}=\sum_a|a,a\rangle\langle a,a|,\qquad
 \Pi_0=\sum_a|a,0\rangle\langle a,0|.
\]
The register $0$ is the all-zero bit string. The map $S$ retains equal input positions and resets the second register on that subspace.

\begin{proposition}[Equality map as a partial isometry]
\label{prop:fluid-equality-partial-isometry}
One has $S^\dagger S=\Pi_{\mathrm{eq}}$, $SS^\dagger=\Pi_0$, and $\|S\|=1$. The expression
\[
 B=S\otimes I+(S-I)\otimes X
\]
on two position registers and a flag qubit is not unitary.
\end{proposition}
\begin{proof}
Multiply the two finite sums for $S$ and $S^\dagger$. Orthonormality eliminates every cross term, giving the two projectors. Hence $S$ is isometric on the equality subspace and zero on its orthogonal complement, which proves the norm statement. For any nonzero $a$,
\[
 B|a,a,0\rangle
 =|a,0,0\rangle+|a,0,1\rangle-|a,a,1\rangle.
\]
The three output vectors are pairwise orthogonal. Thus a unit input has output squared norm three, contradicting unitarity.
\end{proof}

This counterexample applies to the expression printed in (D56), with its stated square embedding $S=\sum_a|a,0\rangle\langle a,a|$. It is not an obstruction to constructing an equality block encoding.

\begin{proposition}[Reversible equality encoding and its success probability]
\label{prop:fluid-equality-unitary}
Define $U|a,b\rangle=|a,b\mathbin{\mathrm{xor}}a\rangle$. Then $U$ is unitary and $\Pi_0U=S$. A flag measurement can realize this block with success probability $\|Sx\|^2$ on any normalized input $x\in\mathcal H\otimes\mathcal H$.
\end{proposition}
\begin{proof}
The basis permutation $U$ is its own inverse. Applying $\Pi_0$ to its output retains exactly the inputs with $b=a$, proving $\Pi_0U=S$. For an explicit flag, define the reversible basis permutation
\[
 V|a,c,z\rangle=|a,c,z\mathbin{\mathrm{xor}}\mathbf1_{c\ne0}\rangle,
 \qquad W=V(U\otimes I).
\]
Both factors are unitary. Starting the flag in $|0\rangle$ and retaining outcome zero gives
\[
 (I\otimes\langle0|)W(x\otimes|0\rangle)=Sx.
\]
The Born rule yields success probability $\|Sx\|^2$. Writing $x=\sum_{a,b}x_{ab}|a,b\rangle$ makes this probability $\sum_a|x_{aa}|^2$. When it is positive, the conditional output is $Sx/\|Sx\|$; when it is zero, that outcome never occurs.
\end{proof}

No copying of an arbitrary unknown quantum state occurs here. The XOR operation copies a basis label into a zero register, and on superpositions acts coherently, generally producing correlations. Applied to $x\otimes x$, the equality branch contains squared amplitudes, but preparing two copies already presupposes repeated access to the state preparation. Its success probability remains part of the operation.

\begin{proposition}[Contraction-normalized rectangular dilation]
\label{prop:fluid-halmos-dilation}
Let $A:\mathcal H_{\mathrm{in}}\to\mathcal H_{\mathrm{out}}$ be a finite complex matrix and let $\alpha>0$ satisfy $\alpha\ge\|A\|$. Put $C=A/\alpha$ and
\[
 D_{\mathrm{out}}=(I-CC^\dagger)^{1/2},\qquad
 D_{\mathrm{in}}=(I-C^\dagger C)^{1/2}.
\]
Then
\[
 \mathcal U_C=
 \begin{pmatrix}C&D_{\mathrm{out}}\\D_{\mathrm{in}}&-C^\dagger\end{pmatrix}
 :\mathcal H_{\mathrm{in}}\oplus\mathcal H_{\mathrm{out}}
 \longrightarrow\mathcal H_{\mathrm{out}}\oplus\mathcal H_{\mathrm{in}}
\]
is unitary. On normalized input $(x,0)$, projection onto the first output summand succeeds with probability $\|Ax\|^2/\alpha^2$ and, when successful, produces $Ax/\|Ax\|$. Conversely, a block of a unitary representing $A/\alpha$ requires $\alpha\ge\|A\|$.
\end{proposition}
\begin{proof}
Take a singular-value decomposition $C=L\Sigma R^\dagger$. Every singular value $s_j$ lies in $[0,1]$. The positive square roots therefore exist, and the SVD gives
$CD_{\mathrm{in}}=D_{\mathrm{out}}C$: on each paired singular direction both sides multiply by $s_j\sqrt{1-s_j^2}$, and both vanish on the unpaired directions. Consequently
\[
 \mathcal U_C\mathcal U_C^\dagger
 =\begin{pmatrix}
 CC^\dagger+D_{\mathrm{out}}^2&CD_{\mathrm{in}}-D_{\mathrm{out}}C\\
 D_{\mathrm{in}}C^\dagger-C^\dagger D_{\mathrm{out}}&D_{\mathrm{in}}^2+C^\dagger C
 \end{pmatrix}=I.
\]
The same calculation gives $\mathcal U_C^\dagger\mathcal U_C=I$. Its action on $(x,0)$ is $(Cx,D_{\mathrm{in}}x)$, proving the measurement formulas. Finally, the norm of a compression of a unitary is at most one, so any such block must satisfy $\|A/\alpha\|\le1$.
\end{proof}

Thus (I65) must use the contraction $\Sigma_2/\alpha$, including the same scaling inside both square roots. The reviewed paper's equation~(D58a) gives $\alpha=6/\bar\tau^\star$ for its $D=2$ collision matrix. A normalization cannot be removed merely because the requested output is eventually renormalized: it determines the success amplitude and enters the resources of the larger block-encoding algorithm.

\begin{proposition}[Linear instrument, conditional normalization]
\label{prop:fluid-postselection-instrument}
The successful and unsuccessful branches of the preceding dilation define linear completely positive maps
\[
 \rho\longmapsto C\rho C^\dagger,\qquad
 \rho\longmapsto D_{\mathrm{in}}\rho D_{\mathrm{in}}.
\]
Their total trace equals $\operatorname{tr}\rho$. Normalizing the successful branch does not turn it into deterministic nonlinear quantum evolution.
\end{proposition}
\begin{proof}
Each branch has a single Kraus operator and is completely positive: after any ancillary extension it sends a positive operator to a conjugate of that operator. Since $C^\dagger C+D_{\mathrm{in}}^2=I$, cyclicity of trace proves the trace identity. The conditional state is divided by $\operatorname{tr}(C\rho C^\dagger)$ only after the corresponding random outcome occurs. Discarding that probability would change the physical operation.
\end{proof}

\section{A global compatibility test for a spin representation}
\label{sec:quantum-spin-scope}

Meng and Yang~\cite{NSMengYang2024} give a two-component spin representation of viscous flow. Their evolution is nonlinear and non-Hermitian. In the constant-density formulation, their Eq.~(15) requires an auxiliary scalar $\phi$ satisfying
\begin{equation}
 \omega\cdot\nabla\phi
 =-\frac{2\nu}{\rho_0}|\nabla\psi|^2(u\cdot\omega),
 \qquad \omega=\nabla\times u.
 \label{eq:spin-auxiliary-compatibility}
\end{equation}
The source explicitly leaves global solvability open. The following calculation makes a necessary global condition concrete. It concerns a particular global lift, not the existence of the velocity field.

\begin{proposition}[Periodic compatibility for a smooth global lift]
\label{prop:spin-periodic-compatibility}
Let $\mathbb T^3=(\mathbb R/2\pi\mathbb Z)^3$, let $\nu,\rho_0,m,\hbar>0$, and fix a time. Suppose $\psi:\mathbb T^3\to\mathbb C^2$ is a globally defined $C^3$ periodic spinor with $\psi^*\psi=\rho_0$, and $\phi:\mathbb T^3\to\mathbb R$ is a globally defined $C^2$ periodic scalar. Define
\[
 u_j=\frac{\hbar}{m\rho_0}\operatorname{Im}(\psi^*\partial_j\psi),
 \qquad |\nabla\psi|^2=\sum_{j=1}^3\|\partial_j\psi\|^2,
 \qquad \omega=\nabla\times u.
\]
If \eqref{eq:spin-auxiliary-compatibility} holds, then
\[
 \int_{\mathbb T^3}|\nabla\psi|^2(u\cdot\omega)\,dx=0,
 \qquad
 |\nabla\psi|^2\geq\frac{m^2\rho_0}{\hbar^2}|u|^2.
\]
\end{proposition}
\begin{proof}
Periodicity and $\nabla\cdot\omega=0$ give
$\int\omega\cdot\nabla\phi=-\int\phi\nabla\cdot\omega=0$.
Integrating \eqref{eq:spin-auxiliary-compatibility} and using $\nu/\rho_0>0$ proves the first assertion. For each coordinate, the complex Cauchy--Schwarz inequality gives
\[
 |\operatorname{Im}(\psi^*\partial_j\psi)|^2
 \leq |\psi^*\partial_j\psi|^2
 \leq\rho_0\|\partial_j\psi\|^2.
\]
Summing and substituting the definition of $u$ gives the second assertion.
\end{proof}

\begin{proposition}[A smooth fluid solution excluded by these global-lift assumptions]
\label{prop:spin-beltrami-obstruction}
For every $\nu>0$, the velocity and pressure
\[
 u(t,x,y,z)=e^{-\nu t}(\sin z,\cos z,0),\qquad p=0
\]
are a smooth periodic solution of
$\partial_tu+(u\cdot\nabla)u=-\nabla p+\nu\Delta u$,
$\nabla\cdot u=0$, for all $t\geq0$.
At no finite time can this velocity satisfy all the global-lift assumptions of Proposition~\ref{prop:spin-periodic-compatibility}.
\end{proposition}
\begin{proof}
The components depend only on $z$, and the third component vanishes. Hence the divergence and advective derivative are zero. Direct differentiation gives $\Delta u=-u$, $\partial_tu=-\nu u$, and $\nabla\times u=u$, establishing the equations. Set $a=e^{-\nu t}>0$. Then $|u|^2=u\cdot\omega=a^2$ everywhere. The preceding pointwise bound would imply
\[
 \int_{\mathbb T^3}|\nabla\psi|^2(u\cdot\omega)\,dx
 \geq\frac{m^2\rho_0}{\hbar^2}a^4(2\pi)^3>0,
\]
contradicting its necessary zero value.
\end{proof}

The hypotheses include a single-valued periodic auxiliary scalar and a single globally defined smooth periodic spinor. Patchwise lifts, transition functions, singularities, other boundary conditions, or a different auxiliary system require a separate analysis and are not ruled out by this proposition. The example has an explicitly smooth solution: it provides no evidence of fluid blowup. Nor does it contradict the source's acknowledged global qualification. It identifies a concrete obligation before treating a local or conditional representation as a global equivalence for all periodic flows. These two propositions have complete ordinary proofs here; they are not included among the companion's Lean-checked declarations.

\paragraph{A separate algebraic obligation.}
For a finite real matrix $A$, let $A^+$ be its Moore--Penrose pseudoinverse. Substitution shows that $f=A^+S$ solves $Af=S$ exactly if and only if $AA^+S=S$. A singular-value decomposition makes $AA^+$ the orthogonal projector onto $\operatorname{range}A$; otherwise the residual is $(I-AA^+)S$. At positive rank,
\[
 \|A^+S\|\leq\frac{\|S\|}{\sigma_{\min}^+(A)},
\]
where $\sigma_{\min}^+$ is the smallest positive singular value. The compatible example $A_\varepsilon=\operatorname{diag}(\varepsilon,0)$, $S=(1,0)^T$, $\varepsilon>0$, has $A_\varepsilon^+S=(\varepsilon^{-1},0)^T$. Thus the notation of a pseudoinverse does not remove conditioning problems near rank changes. When applying such an auxiliary solve to a spin-field model, range compatibility and uniform estimates must be established for its actual structured right-hand side. The generic example alone does not decide that structured question.


\section{A mathematically constrained quantum ToricGT extension}

The previously implemented quantum ToricGT was trained and checked on tokenization-path data. Its weights cannot be presented as a validated fluid model. The reusable component is its architecture: incidence-aware vertex and edge tokens, local quantum features, classical aggregation and an explicit certificate after prediction. This section specifies a fluid adaptation and proves the structural certificate independently of training.

For a surface mesh, tokens should include oriented incidence, the chosen mass matrices, local coordinates or invariant metric features, boundary tags, resolved velocity and any retained memory state. A surface's embedding alone is insufficient. Gauge-like changes of cochain basis should transform the fields and metric together. Under mesh relabeling, the local circuit input and classical aggregation must obey a stated equivariance rule; otherwise predictions can depend on arbitrary vertex numbering.

Let $|\psi_\theta(x)\rangle=U_\theta(x)|0\rangle$ be a normalized finite circuit and let
$q_j(x)=\langle\psi_\theta(x)|O_j|\psi_\theta(x)\rangle$ for Hermitian $O_j$ with $\|O_j\|\le1$. Then $q_j\in[-1,1]$. A classical readout can construct a matrix $B_\theta(x)$ and a skew matrix $K_\theta(x)=\tfrac12(A_\theta(x)-A_\theta(x)^T)$. In divergence-free coordinates, use the vector field of Theorem~\ref{thm:fluid-finite-closure}. The resulting dissipation identity is true for every choice of circuit parameters; it does not rely on the optimizer finding a good model.

\begin{proposition}[A parameterized positive constitutive coefficient]
\label{prop:fluid-positive-readout}
Let $0<\nu_{\min}<\nu_{\max}$ and let $\ell_\theta(x)$ be any finite real readout of circuit features. Then
\[
 \nu_\theta(x)=\nu_{\min}+
 (\nu_{\max}-\nu_{\min})\frac1{1+e^{-\ell_\theta(x)}}
\]
lies strictly between the two bounds. If the readout is locally Lipschitz, using this coefficient in a factorized finite viscosity preserves the energy sign in Theorem~\ref{thm:fluid-finite-closure}.
\end{proposition}
\begin{proof}
The exponential is positive, so the logistic factor is strictly between zero and one. Its derivative is bounded. Composition with a locally Lipschitz readout is locally Lipschitz. Multiplication of a squared norm by a positive scalar preserves nonnegativity, giving the asserted sign.
\end{proof}

\begin{proposition}[The circuit derivative and a nonlinear loss]
\label{prop:fluid-parameter-shift}
Suppose one parameter occurs in one gate $e^{-i\theta P/2}$ with $P=P^*$ and $P^2=I$. For any expectation $q(\theta)$ measured after the full circuit,
\[
 q'(\theta)=\tfrac12[q(\theta+\pi/2)-q(\theta-\pi/2)].
\]
For a differentiable loss $L(q_1,\ldots,q_m)$, its derivative is the chain-rule sum
$\sum_j(\partial L/\partial q_j)q'_j$. Shifting the entire nonlinear loss is not generally the same operation.
\end{proposition}
\begin{proof}
Expand the gate as $\cos(\theta/2)I-i\sin(\theta/2)P$. The expectation is of the form $a+b\cos\theta+c\sin\theta$ after the fixed preceding and following operations are absorbed into the state and observable. Differentiation and the two shifted trigonometric identities prove the first assertion. The second is ordinary multivariable differentiation. For example, if $q=\cos\theta$ and $L=q^2$, shifting $L$ by $\pm\pi/2$ gives equal values and an erroneous zero, whereas $L'=-\sin2\theta$ in general.
\end{proof}

Shared parameters in multiple gates require a sum of gate-occurrence derivatives. The previous ToricGT experiments already distinguish this chain rule from an incorrect whole-loss shift. For a fluid extension it is equally essential: a conservation or energy penalty is a nonlinear loss and must be differentiated correctly.

\subsection{A concrete training and validation specification}

A future trained fluid ToricGT should be assessed against the same dynamical problem solved by its classical baselines. Data must identify viscosity, forcing, domain, mesh, time step, reference solver and truncation error. Split by trajectories and physical parameter families, not by randomly mixing nearby snapshots of the same trajectory. Include independently held-out meshes and resolutions. Compare with a classical model of similar trainable size, a frozen-circuit model, the physical baseline without learned closure and an exact manufactured solution.

Report resolved trajectory error, derivative or stress error where available, divergence, energy balance, pressure residual and total wall time including data transfer and readout. A projection or repair can make a constraint pass even when a predictor is inaccurate, so pre-repair and post-repair results must be distinguished. A small training residual is not a certified PDE residual unless its sampling, quadrature and interpolation errors have also been bounded.

Memory is a plausible place to use learning: approximate a stable auxiliary realization of a kernel, with enough state to distinguish histories that share the same instantaneous resolved velocity. This is more defensible than asking a memoryless network to invert an information-losing projection. For a linear stable hidden subsystem, Theorem~\ref{thm:fluid-memory} supplies a reference target and a complete exact comparison. For a nonlinear flow, the corresponding error and stability hypotheses require new analysis.

\section{The executable examples and their proof boundaries}

The new laboratory separates four tasks. The Taylor--Green field is an exact two-dimensional periodic Navier--Stokes solution with known pressure and decay, used to test derivatives and normalization. The integer covering panel evaluates the actual order matrix from the released proof. The circuit panel computes a two-qubit state and uses its bounded features to select a positive heat-channel parameter. The concentrating-core panel displays the released scaling exponents, without claiming to compute the source's fully corrected singular field.

The heat circuit is an untrained parameterized example. It is not a fluid-trained ToricGT, and the finite heat channel does not include the nonlinear momentum term. The code records actual finite invariants and numerical discrepancies. The companion's selected Lean declarations formalize energy, projection, finite memory and order identities, with precise component mappings; they do not certify every browser operation or every theorem in these chapters.

Taken together, these examples make information reduction operational. One can see which quantities are preserved by a complex, which are lost under projection, which reappear as a memory kernel or a quadratic stress, and which remain controlled by a positive quantum-channel construction. The released Navier--Stokes proof then occupies its proper place: a continuum construction whose all-order analytic estimates establish a forced singularity, and whose algebraic components can be understood and checked within this wider framework.

\chapter{Exact reduction contracts and the information lost by tropicalization}
\label{ch:reduction-contracts}

\reviewstatus{This chapter adds exact factorization, discrete trajectory, error,
and signed polynomial results. The companion module
\texttt{ReductionContracts.lean} contains sixteen kernel-checked declarations
with a per-declaration axiom manifest. The cited analytic-forcing obstruction
is an external theorem, not a result formalized by this module.}

Information reduction becomes a mathematical claim when the retained data
determine the quantity being predicted. This chapter supplies a necessary and
sufficient condition, then a quantitative obstruction when the condition fails.
It also shows precisely why the orders of two signed terms do not determine
the order of their sum. These statements make the project's graph, order, and
tropical interfaces testable without identifying them with a continuum proof.

\section{The exact closure criterion}

Let $X$ be a state space, $Y$ the space of retained observations, and $Z$ an
output space. A map $P:X\to Y$ describes what the representation retains;
$Q:X\to Z$ describes what a model must predict. No topology or probability
distribution is needed for the following statement.

\begin{theorem}[Factorization through retained information]
\label{thm:reduction-fiber-factorization}
Suppose $P$ is surjective. There is a map $G:Y\to Z$ satisfying
$G(Px)=Qx$ for every $x$ if and only if
\[
 Px=Py\quad\Longrightarrow\quad Qx=Qy
 \qquad(x,y\in X).
\]
When it exists, $G$ is unique.
\end{theorem}
\begin{proof}
If $G$ exists, equal observations have equal images under $G$, proving the
condition. Conversely, for each $b\in Y$ choose $x_b$ with $P(x_b)=b$ and
define $G(b)=Q(x_b)$. The fiber condition makes this independent of the
choice and gives $G(Px)=Qx$. If $G,H$ both factor $Q$, surjectivity gives,
for every $b=Px$, the equality $G(b)=Qx=H(b)$.
\end{proof}

Surjectivity is a modeling convention, not a demand that a compressed
representation cover an arbitrary ambient tensor space: one can take $Y$ to
be the image of $P$. The theorem proves existence of a set-theoretic map.
It does not prove that $G$ is continuous, measurable, Lipschitz, efficiently
computable, or representable by a particular neural architecture. Those are
separate requirements on a usable model.

For a discrete update $F:X\to X$, take $Q=P\circ F$. Then the criterion
is exactly preservation of the equivalence relation $x\sim y$ defined by
$Px=Py$. This gives a precise meaning to an exact reduced dynamics.

\begin{proposition}[Propagation of an exact discrete reduction]
\label{prop:reduction-trajectory}
Suppose $P\circ F=G\circ P$, $u_{n+1}=F(u_n)$,
$v_{n+1}=G(v_n)$, and $Pu_0=v_0$. Then $Pu_n=v_n$ for every natural
number $n$.
\end{proposition}
\begin{proof}
The initial equality is given. If it holds at $n$, then
$Pu_{n+1}=PF(u_n)=G(Pu_n)=G(v_n)=v_{n+1}$. Induction proves the claim.
\end{proof}

For a differentiable observation of a continuous-time trajectory
$\dot x=F(x)$, the chain rule instead gives
\[
 \frac{d}{dt}P(x(t))=DP(x(t))F(x(t)).
\]
The quantity to factor is therefore $Q(x)=DP(x)F(x)$. Only for a linear
observation does this become $Q(x)=PF(x)$. A nonlinear tokenizer or learned
coordinate map cannot silently use the linear rule. The factorization
theorem applies to this $Q$ as a map, while differentiability and existence
of trajectories remain analytic premises. The Lean trajectory declaration
in this chapter concerns discrete time only.

For example, $P(x,y)=x$ discards the variable needed to predict $Q(x,y)=y$.
States $(0,0)$ and $(0,1)$ share their observation and have distinct outputs.
No deterministic function of the observation alone can give both outputs.
Adding a memory state or retaining $y$ can change the representation and
remove this obstruction; additional network capacity on the same input cannot.

\section{A lower bound on every approximate closure}

The preceding obstruction has a quantitative version, independent of how
$G$ is trained or parameterized.

\begin{proposition}[Unresolved diameter bounds approximation error]
\label{prop:reduction-error-diameter}
Let $Z$ be a seminormed additive commutative group. For any
$P:X\to Y$, $Q:X\to Z$, and $G:Y\to Z$, if $Px=Py$, then
\[
 \|Qx-Qy\|\leq\|Qx-G(Px)\|+\|Qy-G(Py)\|.
\]
Consequently, if each error on the right is at most $\varepsilon$,
then $\|Qx-Qy\|\leq2\varepsilon$.
\end{proposition}
\begin{proof}
Put $g=G(Px)=G(Py)$ and write
$Qx-Qy=(Qx-g)+(g-Qy)$. The triangle inequality and
$\|g-Qy\|=\|Qy-g\|$ give the first assertion. Bounding the two terms
separately gives the second.
\end{proof}

This is a worst-case bound on a specified pair. A probabilistic average-loss
claim additionally needs a distribution and a loss; a sampled dataset does
not establish a diameter bound on every unseen fiber. The theorem provides
a diagnostic for proposed reductions: find equal retained states with
separated target stresses or derivatives. A quantum expectation readout,
graph embedding, or toric feature map is subject to the same criterion once
its accessible classical observation has been specified.

\section{Signed cancellation is invisible to orders alone}

For $p\in\mathbb Q[X]$, let $\operatorname{ord}_0(p)$ be its smallest
exponent with nonzero coefficient, with $\operatorname{ord}_0(0)=\infty$.
The finite-order statements below use nonzero inputs and nonzero residuals,
so they do not exploit a convention at zero. Tropical order calculations
are useful bookkeeping, but their additive rule must retain the possibility
of cancellation; compare the positive-semiring setting of
\citet{ReductionLitvinov2005}.

\begin{proposition}[Arbitrarily large cancellation gain]
\label{prop:reduction-signed-order}
For every integer $k\geq0$, put
\[
 p=X,\qquad q_k=-X+X^{k+2}.
\]
Then $p,q_k,p+q_k$ are nonzero and
\[
 \operatorname{ord}_0(p)=\operatorname{ord}_0(q_k)=1,
 \qquad \operatorname{ord}_0(p+q_k)=k+2>1.
\]
In particular, the identity
$\operatorname{ord}_0(p+q)=\min\{\operatorname{ord}_0p,
\operatorname{ord}_0q\}$ is false in general even for nonzero summands
and a nonzero sum.
\end{proposition}
\begin{proof}
The constant coefficients vanish. The coefficient of $X$ in $p$ is one
and that in $q_k$ is minus one, since $k+2\geq2$. Each summand has order
one and is nonzero. Their sum is exactly $X^{k+2}$, which is nonzero
and has order $k+2$. This proves all assertions.
\end{proof}

The integer $k$ is arbitrary. Thus no fixed correction to the minimum of
the input orders can recover all cancellation gains. Signed leading
coefficients must first be compared; when they cancel, higher coefficients
must be examined. For complex amplitudes, phases carry analogous information.
This issue is especially consequential for a momentum residual: its
transport, pressure, temporal, and viscous terms carry signs, and smallness
of their exact sum is a different property from the size of each term.
An order graph can organize candidate balances. A proof must still establish
the coefficient identities and bounds on the remaining terms.

\section{A current analytic obstruction to removing the force}

Theorem~1.1 of \citet{ReductionCIV2026}, version~2, addresses suitable weak
solutions smooth before the candidate singular time. It assumes a uniform
spatial $C^2$ bound on the force, spatial analyticity locally uniform on
interior time intervals, the anisotropic angular-mean bounds in its
equations~(1.7a)--(1.7b) through derivative order two, and exact axisymmetry
on a positive-radius core at every earlier time. Under these hypotheses the
candidate point is regular. Zero forcing meets the force hypotheses.
Consequently, preserving this structural class while deleting the force
cannot retain the released singularity. This is an obstruction to that
specified class, not a global unforced regularity theorem and not a
disproof of the smooth forced construction. The cited analytic theorem is
used as an external result; this extension has not formalized its proof.

For the project's future developments, this gives three concrete obligations.
A proposed information reduction must state its fibers and the outputs they
determine. A tropical residual certificate must retain signed coefficients
where leading orders tie. A transformed fluid theorem must preserve its
forcing and regularity hypotheses, or separately prove the replacements.
These obligations are compatible with graph-based proof search and quantum
learning, but they cannot be discharged by an analogy between representations.

\section{Formal coverage and reproducibility}

The companion module proves the two directions and uniqueness in
Theorem~\ref{thm:reduction-fiber-factorization}, its discrete-closure
specialization and a mismatched-fiber obstruction,
Proposition~\ref{prop:reduction-trajectory}, both inequalities in
Proposition~\ref{prop:reduction-error-diameter}, and every polynomial
assertion in Proposition~\ref{prop:reduction-signed-order}.
The declarations use ordinary set-theoretic functions, actual seminorms,
and actual rational polynomials; no placeholder predicate stands in for
the missing analytic theorem. The verification manifest records the source
hash, exact compiler and mathlib pins, process exit code, and each
declaration's foundational dependencies. Its sixteen checked declarations
are additional to the earlier finite fluid and quantum modules.

\chapter{Corrected transport, information and quantum contracts for tropical quivers}
\label{ch:toric-transport-audit}
\reviewstatus{This chapter adds explicit corrections to selected claims in the pinned Tropical Quivers and ToricGT research manuscripts. Seventeen finite declarations in TransportContracts.lean are kernel checked; the information-theoretic and stochastic-limit discussions are mathematical arguments with stated hypotheses, not a formalization of all source papers. None of the statements below extends the external forced Navier--Stokes theorem to the unforced problem.}

\section{Source scope and a quantitative contract}

The source revisions are:
\begin{itemize}
\item Tropical\_Quivers\_of\_Archs \cite{ToricAuditQuivers2026}, commit\\
\texttt{2b6ab79a685e2c6233fb8892e59589951feeddad};
\item ToricGT-Program, commit\\
\texttt{f91ba1811b5a4f70d53c7120acf8ca2a0049f85b}.
\end{itemize}
The pinned primary sources are available online.\footnote{Primary sources:
\url{https://github.com/amelie-iska/Tropical_Quivers_of_Archs/tree/2b6ab79a685e2c6233fb8892e59589951feeddad}
and \url{https://github.com/amelie-iska/ToricGT-Program/tree/f91ba1811b5a4f70d53c7120acf8ca2a0049f85b}.}
The companion review gives
the exact files, inspected sections, and unresolved coverage. The first
repository contains a manuscript and its PDF. The second contains five
research-program papers, earlier versions, reference PDFs, and a ToricGT
gitlink. The gitlink points at a separate code revision. A proposed
architecture, a numerical experiment and a proved theorem remain different
kinds of evidence.

The useful common interface is a finite typed representation with a
specified metric, update map, certificate and error budget. A surface
cellulation provides incidence maps; a gentle algebra additionally imposes
path relations; a physical fluid requires constitutive laws and analytic
estimates; a quantum channel requires complete positivity and trace
preservation. Each additional structure has its own contract. The following
corrections make those contracts compositional without identifying them.

\section{Suffix transport is the correct cyclic coboundary}

Let $V_0,\ldots,V_k$ be vector spaces, let
$A_j:V_{j-1}\to V_j$ be linear maps, and let $x_j\in V_j$. Define the edge
residual $r_j=x_j-A_jx_{j-1}$ and the suffix transport
\[
 T_j=A_kA_{k-1}\cdots A_{j+1}:V_j\longrightarrow V_k,
 \qquad T_k=I_{V_k}.
\]
Composition acts right to left. In particular $T_j$ transports a residual
from the target of edge $j$ to the endpoint of the path. No commutativity,
inverse, equal dimension or flatness assumption is required.

\begin{theorem}[Typed transport telescope]
\label{thm:typed-transport-telescope}
With the preceding notation,
\[
 \sum_{j=1}^{k}T_jr_j=x_k-A_k\cdots A_1x_0.
\]
For a closed path, $V_k=V_0$ and $x_k=x_0$, this equals $(I-H)x_0$, where
$H=A_k\cdots A_1$ is the based holonomy. It vanishes for every initial
vector if and only if $H=I$.
\end{theorem}
\begin{proof}
Expand the left side as
\[
 \sum_{j=1}^{k}A_k\cdots A_{j+1}x_j
 -\sum_{j=1}^{k}A_k\cdots A_jx_{j-1}.
\]
The term in the first sum with index $j<k$ equals the term in the second
sum with index $j+1$. These cancel in their common target space. The two
remaining terms are $x_k$ and $-A_k\cdots A_1x_0$. On a closed path, the
residual is zero on every vector precisely when the linear map $I-H$ is
zero. For a particular vector it is enough that $Hx_0=x_0$.
\end{proof}

The equal-space specialization is formalized by induction in
\path{TransportContracts.suffix_transport_telescope}. Its recursive
form is
\[
 s_0=0,\qquad s_{j+1}=A_{j+1}s_j+r_{j+1}.
\]
This is also an efficient implementation: the accumulated discrepancy is
transported once per edge, rather than storing all suffix matrices.

\begin{proposition}[A flat-cycle counterexample to prefix transport]
\label{prop:transport-prefix-counterexample}
Take scalar arrows $A_1=2$, $A_2=1/2$ and states
$x_0=x_2=0$, $x_1=1$. Then $A_2A_1=1$. Nevertheless the alternative
prefix-weighted expression
\[
 A_1r_1+(A_1A_2)r_2=\frac32
\]
is nonzero. The suffix expression $A_2r_1+r_2$ is zero.
\end{proposition}
\begin{proof}
Here $r_1=1$ and $r_2=-1/2$. Substitution gives
$2-1/2=3/2$ for prefixes and $1/2-1/2=0$ for suffixes.
\end{proof}

The pinned Tropical Quivers section on elliptic complexes used forward
prefix products as transport back to a cycle basepoint. The counterexample
shows that its asserted flatness implication does not follow from that
formula. Replacing prefixes by the suffixes above repairs the direction
and multiplication order simultaneously. The source is retained unchanged;
this chapter is an explicit amendment.

\section{A positive Laplacian does not itself supply a complex}

For finite inner-product spaces, let $D_0:C^0\to C^1$ and
$D_1:C^1\to C^2$ be arbitrary linear maps. Define
$L_1=D_0D_0^*+D_1^*D_1$. Positivity holds with no complex assumption:
\[
 \langle w,L_1w\rangle=\|D_0^*w\|^2+\|D_1w\|^2\ge0,
 \qquad \ker L_1=\ker D_0^*\cap\ker D_1.
\]
However, an orthogonal exact/coexact/harmonic decomposition requires
$D_1D_0=0$.

\begin{proposition}[Exact hypothesis for orthogonal Hodge components]
The subspaces $\operatorname{im}D_0$ and $\operatorname{im}D_1^*$ are
orthogonal if and only if $D_1D_0=0$. Under that hypothesis,
\[
 C^1=\operatorname{im}D_0\ \mathbin{\oplus^{\perp}}
 \operatorname{im}D_1^*\ \mathbin{\oplus^{\perp}}\ \ker L_1.
\]
\end{proposition}
\begin{proof}
For every $x,y$,
$\langle D_0x,D_1^*y\rangle=\langle D_1D_0x,y\rangle$.
Vanishing of all these pairings is equivalent to $D_1D_0=0$ by
nondegeneracy. The orthogonal complement of the sum of the two images is
$\ker D_0^*\cap\ker D_1=\ker L_1$. Finite dimensionality gives the
orthogonal direct sum. For arbitrary maps the two images can overlap;
positivity of $L_1$ does not remove that overlap.
\end{proof}

For covariant quiver differences, Theorem~\ref{thm:typed-transport-telescope}
identifies the curvature defect of each based face as $I-H$. An
approximately flat numerical connection is therefore not an exact chain
complex. Its residual norm can be reported and penalized, but its
``cohomology'' must not be treated as the cohomology of an exact differential
unless the relation is established. For gentle algebras the corresponding
path relations must likewise be checked in the actual representation.
Continuum ellipticity, mesh-uniform coercivity and convergence remain
additional analytic requirements.

\section{What a prefix-derived certificate can add}

The mutual-information identities in ToricGT II--IV need an explicit
distinction between the entire available prefix and a lossy model state.
The source definition makes its certificate measurable from the full
prefix, so it cannot add Bayes information to that same full prefix.

\begin{proposition}[Deterministic certificates have zero extra Bayes information]
Let $B,X,Z$ be finite random variables and suppose $Z=f(X)$. Then
\[
 H(B\mid X,Z)=H(B\mid X),\qquad I(B;Z\mid X)=0.
\]
The same conditional-expectation conclusion holds when $Z$ is measurable
with respect to a sigma-field $\mathcal F$ and $B$ has a finite alphabet.
\end{proposition}
\begin{proof}
For every $x$ of positive probability, the event $X=x$ already fixes
$Z=f(x)$. Thus the conditional distribution of $B$ given $(X,Z)$ is the
conditional distribution given $X$. Their conditional entropies coincide,
and their difference is the displayed mutual information. In the sigma-field
form, $\sigma(\mathcal F,Z)=\mathcal F$, so both conditional expectations
condition on the same sigma-field.
\end{proof}

This does not imply that computable certificates are useless. Let
$S=g(X)$ be the representation actually available to a bounded predictor.
In general $I(B;Z\mid S)$ can be positive, even though
$I(B;Z\mid X)=0$. For example, take $X$ uniform on $\{0,1,2,3\}$,
$B=Z=X\bmod2$, and $S=\lfloor X/2\rfloor$. Given $S$, $B$ is a fair bit;
given $(S,Z)$, it is known. The corresponding improvement is one bit, while
the improvement beyond complete $X$ is zero.

\begin{proposition}[Actual predictor improvement includes approximation error]
Let $q_0(\cdot\mid S)$ and $q_1(\cdot\mid S,Z)$ be normalized predictors
positive on the support of the target, with finite expected log losses
$L_0,L_1$. In base-two units,
\begin{align*}
 L_0-L_1={}&I(B;Z\mid S)\\
 &+\mathbb E\,D_{\rm KL}\bigl(p(\cdot\mid S)\|q_0(\cdot\mid S)\bigr)\\
 &-\mathbb E\,D_{\rm KL}\bigl(p(\cdot\mid S,Z)\|q_1(\cdot\mid S,Z)\bigr).
\end{align*}
\end{proposition}
\begin{proof}
For each conditioning value, insert and subtract $-\log_2p(B\mid\cdot)$
in the expected log loss. This gives conditional entropy plus conditional
KL divergence for each predictor. Subtract the two decompositions and use
the defining entropy difference for conditional mutual information.
\end{proof}

A deterministic topology, toric or proof certificate can therefore reduce
the approximation or optimization error of a restricted neural predictor.
That is a measurable training benefit, not new information beyond its
complete source data. A deployable certificate should be evaluated by
held-out predictive loss and its separately accounted artifact and runtime
cost. Leakage control still requires every certificate to be available
before the predicted observation; it does not require its usefulness to
come from a positive conditional mutual information given all that data.
This distinction also prevents a finite topological invariant from being
mistaken for extra evidence about a physical system.

\section{Damping and stochastic limits require their own hypotheses}

For a map $F$ with Lipschitz bound $L$ and $0<\eta\le1$, the relaxed map
$G_\eta=(1-\eta)I+\eta F$ has the upper bound
$\operatorname{Lip}(G_\eta)\le(1-\eta)+\eta L$.
When $L>1$, that upper bound is greater than one for every positive $\eta$.
It cannot certify contraction by decreasing $\eta$.

\begin{proposition}[Damping cannot repair every expanding cycle]
\label{prop:transport-damping-obstruction}
For $F(x)=2x$ on $\mathbb R$, $G_\eta(x)=(1+\eta)x$, so its exact
Lipschitz constant is greater than one for every $\eta>0$.
\end{proposition}
\begin{proof}
For distinct $x,y$, the ratio
$|G_\eta(x)-G_\eta(y)|/|x-y|$ equals $1+\eta>1$.
\end{proof}

The pinned manuscript's assertion that the displayed bound restores
contraction for $L>1$ is therefore incorrect. Damping can still stabilize
other maps: for $F(x)=-2x$, the relaxed multiplier is $1-3\eta$, whose
absolute value is less than one for $0<\eta<2/3$. The correct test concerns
the actual map, a valid norm bound, or a justified local Jacobian criterion;
an unnecessarily large upper bound is not an exact Lipschitz constant.

\begin{proposition}[A sufficient dissipativity condition for relaxation]
\label{prop:transport-dissipative-relaxation}
Let $b$ be a map on a real Hilbert space and suppose that, for all $x,y$,
\[
 \langle b(x)-b(y),x-y\rangle\le-\mu\|x-y\|^2,
 \qquad \|b(x)-b(y)\|\le L_b\|x-y\|,
\]
where $0<\mu\le L_b$. If $0<\eta<2\mu/L_b^2$, then
$x\mapsto x+\eta b(x)$ is a contraction with bound
$\sqrt{1-2\mu\eta+L_b^2\eta^2}<1$.
\end{proposition}
\begin{proof}
For $d=x-y$ and $e=b(x)-b(y)$, expand
$\|d+\eta e\|^2=\|d\|^2+2\eta\langle d,e\rangle+\eta^2\|e\|^2$.
The assumptions bound this by
$(1-2\mu\eta+L_b^2\eta^2)\|d\|^2$.
The coefficient is nonnegative since its minimum over $\eta$ is
$1-\mu^2/L_b^2\ge0$, and it is less than one on the stated interval.
Taking square roots proves the assertion.
\end{proof}

For $G_\eta=(1-\eta)I+\eta F$, apply the proposition to $b=F-I$.
This supplies a usable sufficient condition rather than an inference from
$\operatorname{Lip}(F)>1$. A learned skew-plus-dissipative fluid closure may
preserve the energy of one trajectory while failing this two-trajectory
condition when its coefficients depend on the state; that additional
derivative control must be established separately.

The stochastic ``embedding fluid'' passage needs a second qualification.
The continuous/hybrid GFlowNet theory of \citet{ToricAuditLahlou2023}
does not, by its balance identities alone, establish a small-step limit.
Let a small update have conditional mean
$\mathbb E[\Delta Z\mid Z=z]=b(z)\Delta t+o(\Delta t)$ and conditional
second moment
$\mathbb E[\Delta Z\Delta Z^T\mid Z=z]=a(z)\Delta t+o(\Delta t)$.
For a smooth compactly supported test function $\varphi$, Taylor expansion,
with a controlled remainder, gives
\[
 \frac{\mathbb E[\varphi(Z+\Delta Z)-\varphi(Z)\mid Z=z]}{\Delta t}
 \longrightarrow b\cdot\nabla\varphi+
 \frac12\sum_{i,j}a_{ij}\partial_i\partial_j\varphi.
\]
Under the additional tightness, moment and convergence assumptions needed
for a diffusion limit, the density consequently satisfies
\[
 \partial_t\rho=-\nabla\cdot(b\rho)
 +\frac12\sum_{i,j}\partial_i\partial_j(a_{ij}\rho).
\]
This conservative Markov limit has no independently generated source or sink.
Terms $\sigma_{\rm src}-\sigma_{\rm sink}$ require separately specified
branching, killing or injection dynamics, with their own scaling and
convergence hypotheses; they do not follow from the displayed conditional
first and second moments alone.
The pure continuity equation follows only when the second-order term
vanishes in the chosen scaling, or when the current has been defined to
include that diffusive contribution. A GFlowNet balance objective by itself
does not supply this limit. Likewise an antisymmetric-Jacobian penalty is a
useful proposed regularizer, not a theorem that semantic circular reasoning
has been eliminated or that physical Navier--Stokes dynamics is reproduced.

\section{Quantum and gentle-algebra versions of the transport contract}

The transport identity extends directly to linear maps between matrix
spaces. Let $\Phi_j$ be quantum channels and let $\rho_j$ be proposed
density operators. Set $r_j=\rho_j-\Phi_j(\rho_{j-1})$. The $r_j$ are
Hermitian trace-zero discrepancies, not density operators. Applying
Theorem~\ref{thm:typed-transport-telescope} to the underlying real vector
spaces of Hermitian matrices gives
\[
 \sum_j(\Phi_k\circ\cdots\circ\Phi_{j+1})(r_j)
 =\rho_k-(\Phi_k\circ\cdots\circ\Phi_1)(\rho_0).
\]
Trace-norm contractivity of each channel on Hermitian discrepancies then
bounds the final discrepancy by $\sum_j\|r_j\|_1$. This uses actual channel
assumptions; an arbitrary learned real connector does not acquire them by
being called quantum. A noisy cycle need not have identity holonomy, and
no inverse quantum operation is assumed.

For completeness, this contractivity has a short finite-dimensional proof.
For Hermitian $Y$, spectral decomposition gives
$\|Y\|_1=\max_{-I\le K\le I}\operatorname{Tr}(KY)$.
The adjoint of a positive trace-preserving map is positive and unital, so
$-I\le K\le I$ implies $-I\le\Phi^*(K)\le I$. Therefore
$\operatorname{Tr}(K\Phi(Y))=\operatorname{Tr}(\Phi^*(K)Y)\le\|Y\|_1$.
Taking the maximum proves $\|\Phi(Y)\|_1\le\|Y\|_1$; the triangle
inequality applied to the telescope gives the stated error bound. Complete
positivity supplies the stronger physical requirement of positivity after
adjoining an arbitrary ancillary system.

\begin{proposition}[Trace preservation obstructs a zero path relation]
\label{prop:transport-trace-zero-obstruction}
Suppose $\Phi$ and $\Psi$ are composable trace-preserving linear maps on
matrix spaces containing a normalized state $\rho$. Then
$\Psi\circ\Phi$ is not the zero map. Hence a nontrivial relation
$ba=0$ of a gentle path algebra cannot be represented by making both
composable arrows $a,b$ into trace-preserving channels on those full spaces.
\end{proposition}
\begin{proof}
Trace preservation gives
$\operatorname{Tr}(\Psi(\Phi(\rho)))=\operatorname{Tr}\rho=1$.
The zero map would instead produce trace zero. This contradiction uses
only linearity, normalization and trace preservation, so it applies in
particular to completely positive trace-preserving channels.
\end{proof}

The proposition is formalized, with an abstract linear trace functional, as
\begin{quote}\small
\path{TransportContracts.trace_preserving_composite_nonzero}.
\end{quote}
A legitimate quantum interpretation of
forbidden paths can use completely positive trace-non\-increasing instrument
branches, explicit rejection flags, or linear amplitude maps whose ranges
and kernels enforce the relation. For instrument branches, a forbidden
composition has zero probability; the complete instrument, including all
outcomes, is trace preserving. Normalizing a postselected branch is a
separate nonlinear state update and its success probability must be retained.
This distinction aligns the gentle algebra with operational quantum
information without assigning impossible channel semantics to its arrows.

\section{Tropical approximation is not an automatic universality transfer}

Finite ReLU networks give continuous piecewise-affine functions. With
arbitrary real slopes they can be expressed using differences of maxima of
real affine functions. Classical lattice-based toric geometry imposes
integrality or rationality data beyond this statement. Nonzero biases also
produce polyhedral complexes rather than necessarily conical fans. These
are compatible descriptions once the category and homogenization convention
are specified.

Softmax is smooth and is not a finite piecewise-affine map on an open
nontrivial logit region. A local piecewise-affine surrogate therefore needs
an approximation error bound. Moreover, dot-product logits with both query
and key variable are bilinear. A maximum of such scores need not be
piecewise affine in the joint inputs. The finite maximum-reduction identity
used by ring evaluation remains exact, independently of whether the score
functions are tropical polynomials.

The graph-universality argument in the Tropical Quivers manuscript also
requires a measure-specific step. The cited primary transformer theorem
is Yun et al.~\cite{ToricAuditYun2020}; the separate TokenGT construction is
\citet{ToricAuditTokenGT2022}. Yun et al.'s Theorems~2--3 use
$L^p$, $1\le p<\infty$, rather than the uniform norm asserted by the
fixed-cell theorem in the program's Paper~I.
Approximation in ambient Lebesgue
$L^p$ does not in general imply approximation on a lower-dimensional
tokenized graph family. For example,
$g_\epsilon(t)=\max(1-|t|/\epsilon,0)$ has
$\|g_\epsilon\|_{L^p(\mathbb R)}^p=2\epsilon/(p+1)$, tending to zero,
while $g_\epsilon(0)=1$ for every $\epsilon>0$. A tokenized domain supported
at that point sees no convergence. A valid transfer must specify the
graph-domain measure or prove uniform approximation there, and must match
the graph-relabeling action to the token symmetry group. Continuous
detokenization, finite padded sizes and admissible output types must also
be retained. Expressive approximation under those assumptions does not
establish that gradient descent finds the approximant or that unseen
physical distributions generalize.

The bundled transformer-geometry reference is version 1 of
\texttt{2604.14727}. Its current version 2 \cite{ToricAuditSuLiu2026}
explicitly distinguishes fixed-key
top-one routing, a log-lifted value parameterization, and finite-temperature
approximations.
That distinction should be retained when importing its geometry into a
trainable quiver with jointly varying keys and queries. The separate
tropical-attention paper \cite{ToricAuditHashemi2025} proposes a specified
algebraic kernel and reports algorithmic benchmarks. Its Appendix~C also
contains the localized errors repaired next, so this extension does not
import that proof package wholesale or transfer its claims to a different
dot-product-plus-maximum operation.

\section{Independent repairs to the tropical-attention reference}

Version~2 of \citet{ToricAuditHashemi2025}, Appendix~C's proof of
Lemma~3.1, uses $v_\lambda(a)=\log a-\lambda$ for $a>0$ and claims a
non-Archimedean additive inequality. This logarithm is an Archimedean
coordinate change. Neither a constant offset nor projective normalization
turns ordinary addition into exact maximum.

\begin{proposition}[Exact logarithmic identities and dequantization error]
\label{prop:tropical-log-repair}
For $a,b>0$ and real $\lambda$,
\[
 v_\lambda(ab)=v_\lambda(a)+v_\lambda(b)+\lambda.
\]
For any $m\ge1$, $z_1,\ldots,z_m\in\mathbb R$ and $\tau>0$,
\[
 M\le\tau\log\left(\sum_{j=1}^m e^{z_j/\tau}\right)
 \le M+\tau\log m,
 \qquad M=\max_j z_j.
\]
In particular, for $a=b=1$,
$v_\lambda(a+b)=\log2-\lambda>
\max\{v_\lambda(a),v_\lambda(b)\}$.
\end{proposition}
\begin{proof}
The product identity follows by expanding $\log(ab)=\log a+\log b$.
Factoring $e^{M/\tau}$ out of the sum leaves a number between one and
$m$. Apply the increasing logarithm and multiply by $\tau>0$ to obtain
the two bounds. The displayed counterexample follows from $\log2>0$.
\end{proof}

Thus the source's product offset has the wrong sign, and its additive
inequality omits a strictly positive term even on positive inputs.
The valid limiting statement is convergence to maximum as $\tau\downarrow0$
for a fixed finite number of terms, with the explicit error above. If the
number of terms varies, $\tau\log m$ must also be controlled.
For signed or complex terms, cancellation requires additional information,
as proved in Proposition~\ref{prop:reduction-signed-order}. These
corrections do not invalidate ordinary maximum-plus dynamic programming.

The same Appendix~C, Lemma~C.1, uses a large finite penalty as if it were
the exact value $-\infty$. A finite mask can be exact on a bounded input
domain, but it is not an input-independent annihilator on unbounded values.

\begin{proposition}[When a finite tropical mask is exact]
\label{prop:tropical-finite-mask}
Suppose a desired maximum has at least one admissible candidate of score
at least $L$. Suppose every excluded unpenalized candidate has score at
most $U$. Giving every excluded candidate an additive penalty $-M$ makes
it strictly unable to attain the maximum whenever $M>U-L$.
Without an upper bound on excluded values, no fixed finite $M$ guarantees
this conclusion.
\end{proposition}
\begin{proof}
Every penalized excluded score is at most $U-M<L$, below an admissible
score. Conversely, with an admissible score zero and any finite penalty
$M$, an excluded raw score $M+1$ becomes one and wins.
\end{proof}

An explicit $-\infty$ structural mask, or a proved range and sufficient
finite margin, repairs the construction. An all-masked query needs its own
defined output. Bounds on learned input features and their perturbations
must be included in a claimed finite-margin certificate.
The source also assigns $-\infty$ to excluded value entries in that proof.
If such entries are genuinely enforced, they remove those candidates
independently of the finite score penalty. The defect identified here is
the inference from a large finite penalty alone; it is not a counterexample
to a separately and exactly masked maximum.

\begin{proposition}[Exact-step and bounded-length tropical paths]
\label{prop:tropical-horizon-repair}
Write $\oplus=\max$ and $\otimes=+$ for scalar tropical operations, with
their induced matrix operations. Let $D$ be a finite tropical matrix and
let $I$ have zero diagonal and $-\infty$ elsewhere. Then
\[
 v_{t+1}=D\otimes v_t\quad\Longrightarrow\quad
 v_t=D^{\otimes t}\otimes v_0,
\]
whereas
\[
 w_{t+1}=w_t\oplus(D\otimes w_t),\quad w_0=v_0
 \quad\Longrightarrow\quad
 w_t=\left(\bigoplus_{j=0}^{t}D^{\otimes j}\right)\otimes v_0.
\]
Thus a pure $t$-step update and a finite Kleene closure are different
operations in general.
\end{proposition}
\begin{proof}
The first identity is induction and associativity. For the second,
distribute $D$ over the inductive sum; adding the previous sum retains
precisely powers zero through $t+1$. Idempotence removes duplicate terms.
For a direct counterexample take the one-by-one matrix $D=(-1)$ and
$v_0=0$. At $t=2$ the pure update is $-2$, while the maximum of the
zeroth, first and second powers applied to $v_0$ is zero. Even the source's
convention summing powers zero through $t-1$ gives zero, not $-2$.
\end{proof}

This repairs the mismatch in Theorem~C.5 of the reference. For maximum-plus
longest-path closure, deleting a repeated cycle can only improve a path
when every cycle has nonpositive total weight. Under that hypothesis a
best finite path can be chosen with at most $n-1$ edges, with the length-zero
path allowed. For minimum-plus shortest paths the corresponding hypothesis
is absence of negative-weight cycles. The sign and semiring cannot be
interchanged. On graphs with profitable cycles, a bounded-horizon
computation is still defined, but it is not the unbounded closure.

These are complete finite arguments for the operations actually used.
They do not assert that a trained attention stack has learned the exact
parameters, nor that any method solves arbitrary NP-hard instances in
polynomial resources. The end-to-end scientific contract remains a measured
prediction together with independently checked algebraic structure.

\section{A nonnormal obstruction to a Jacobian power certificate}

The Tropical Quivers manuscript's section on Jacobian spectral-radius
estimation iterates $v\mapsto Jv/\|Jv\|$ but labels the resulting estimate
as the largest singular value. For nonnormal matrices these are different
quantities. The following exact witness isolates the error without any
floating-point approximation.

\begin{proposition}[A nilpotent Jacobian with arbitrarily large gain]
\label{prop:nonnormal-jvp}
For any real $M>1$, put
\[
 J_M=\begin{pmatrix}0&M\\0&0\end{pmatrix},\qquad
 e_2=\begin{pmatrix}0\\1\end{pmatrix}.
\]
Then $J_M^2=0$, $\rho(J_M)=0$, and $\|J_M\|_2=M$.
Nevertheless $\|J_Me_2\|_2^2=M^2>1$, while the next application of
$J_M$ is zero. In particular, a vanishing or small eigenvalue power probe
does not certify one-step Euclidean contraction.
\end{proposition}
\begin{proof}
Direct multiplication gives $J_M^2=0$ and
$J_M^TJ_M=\operatorname{diag}(0,M^2)$, whereas
$J_MJ_M^T=\operatorname{diag}(M^2,0)$. Thus $J_M$ is nonnormal and
both eigenvalues are zero. For $x=(x_1,x_2)$,
$\|J_Mx\|_2^2=M^2x_2^2\le M^2\|x\|_2^2$, with equality at
$x=e_2$. This proves the norm formula and all remaining assertions.
\end{proof}

For a general real Jacobian, singular-norm power iteration uses
$J^TJ$, requiring both a Jacobian-vector product and its adjoint product.
Even that correction does not turn a finite numerical estimate into a
certified upper bound: a Rayleigh quotient samples a direction and can
miss the maximizing singular direction. A contraction certificate needs a
proved upper bound, a validated enclosure, or an explicitly controlled
operator parametrization. The nilpotent example can have powers that
eventually vanish while still amplifying perturbations at the first step.

The companion formalizes its explicit matrix products, nonnormality and
squared coordinate gain. The spectral-radius and operator-norm conclusions
are the written finite-dimensional argument above; the corresponding Lean
components are not represented as a complete analytic norm formalization.

\section{Path multiplicity separates tropical growth from operator stability}

The max-plus cocycle in \citet{ToricAuditQuivers2026} records the largest
logarithmic path weight. An ordinary matrix product adds contributions
from many paths. Its operator growth therefore needs a multiplicity term.
This is a qualification of the source's contraction interpretation, not
a rejection of its max-plus subadditivity calculation.

Let $A_t:\mathbb C^{d_{t-1}}\to\mathbb C^{d_t}$, $1\le t\le T$ and $T\ge1$, with
each $d_t\ge1$, and set $P_T=A_T\cdots A_1$. An admissible path from $j$
to $i$ is a sequence $(i_0,\ldots,i_T)$ with $i_0=j$, $i_T=i$ and
all $A_t(i_t,i_{t-1})\ne0$. Let $C_T(i,j)$ count these paths and let
$W_T(i,j)$ be their maximum absolute product weight, taking both values
to be zero when no admissible path exists. Thus $W_T$ is a maximum-product
matrix product; its logarithm, with $\log0=-\infty$, is the corresponding
max-plus product. Let $S_t$ be the zero--one support matrix of $A_t$.

\begin{theorem}[Finite path-multiplicity bound]
\label{thm:tropical-path-multiplicity}
With the preceding notation,
\[
 |P_T(i,j)|\le (|A_T|\cdots|A_1|)(i,j)
 \le C_T(i,j)W_T(i,j),
 \qquad C_T=S_T\cdots S_1.
\]
Writing $N_T=\max_i\sum_jC_T(i,j)$ and
$m_T=\max_{i,j}W_T(i,j)$ gives the induced row-sum norm bound
\[
 \|P_T\|_\infty
 \le\max_i\sum_j C_T(i,j)W_T(i,j)
 \le N_Tm_T.
\]
For square $n$-dimensional arrows, $N_T\le n^T$. Consequently, whenever
$m_T>0$ eventually and
$\limsup_{T\to\infty}T^{-1}\log m_T+\log n<0$, the ordinary matrix
products decay exponentially in this norm.
\end{theorem}
\begin{proof}
Expand each entry of $P_T$ as the sum over the intermediate indices.
The triangle inequality bounds its absolute value by the sum of the
absolute path weights. Each nonzero summand is at most $W_T(i,j)$,
and exactly $C_T(i,j)$ such paths occur. Multiplying the support matrices
counts those paths, proving the entrywise assertions. Sum over $j$ and
maximize over $i$ to obtain the norm inequalities. In dimension $n$, a
fixed final index leaves $T$ freely chosen earlier indices, so there
are at most $n^T$ paths ending at that index. The strict limsup inequality
then gives constants $c>0$ and $T_0$ with
$n^Tm_T\le e^{-cT}$ for $T\ge T_0$.
\end{proof}

If $m_T=0$ at any time, every path of that length vanishes, so all later
products are zero and the conclusion holds directly. More precise bounds
can retain $N_T$ or the entire entrywise table instead of $n^T$.
Using $\log(\epsilon+|A_t(i,j)|)$ with $\epsilon>0$, as in the source,
increases path weights and produces a conservative maximum-product upper
bound; it does not remove the number of summed paths. Conversely, signed
or complex cancellation can make $P_T$ smaller than the positive bound.
These inequalities do not recover such cancellation from magnitudes.

\begin{proposition}[Decaying maximal paths with an expanding linear map]
\label{prop:tropical-branching-counterexample}
For $n\ge2$ and $1/n<a<1$, let $A=a\mathbf1\mathbf1^T$ on
$\mathbb R^n$. Every length-$k$ path has weight $a^k$, so its tropical
growth exponent is $\log a<0$. Nevertheless
\[
 A\mathbf1=na\mathbf1,\qquad
 A^k=n^{k-1}a^k\mathbf1\mathbf1^T,\qquad
 \|A^k\|_\infty=(na)^k>1 \quad(k\ge1).
\]
The path-multiplicity norm bound is attained.
\end{proposition}
\begin{proof}
Each entry is $a$, and each row has $n$ entries. The eigenvector identity
follows by row summation. Since
$(\mathbf1\mathbf1^T)^2=n\mathbf1\mathbf1^T$, induction gives the
power formula. Its row sum is $(na)^k$. There are $n^{k-1}$ paths between
any prescribed pair of endpoints and $n^k$ paths ending at a prescribed
final vertex, so both multiplicity inequalities are equalities.
\end{proof}

For example, $n=2$, $a=3/4$ gives individual path decay $(3/4)^T$ and
ordinary growth $(3/2)^T$. Thus a negative maximal-path exponent alone
cannot certify contraction of a learned fluid closure, a quiver update,
or a quantum-amplitude surrogate. A separately proved norm inequality
or a preserved stochastic/channel structure is required. The exact
regression computes support counts and maximum weights by dynamic
programming and compares them against exhaustive path enumeration on
small, possibly rectangular, signed rational examples. It is a validation
of finite products, not an estimate of an infinite-time physical exponent.

\section{Weighted small-gain certification at a fixed rate}

The weighted small-gain criterion in \citet{ToricAuditQuivers2026} is
useful, but its instruction to minimize the weighted row bound by a
linear program needs a qualification. Suppose $G\ge0$ is a finite
matrix of valid block Lipschitz gains. For positive weights $w_i$ define
$\|x\|_{\infty,w}=\max_i\|x_i\|/w_i$. Its row bound is
\[
 q(w)=\max_i \frac{(Gw)_i}{w_i}.
\]

\begin{proposition}[Fixed-rate feasibility is linear]
\label{prop:weighted-gain-feasibility}
For fixed $q\ge0$, $q(w)\le q$ is equivalent to the componentwise
linear inequalities $Gw\le qw$, $w_i>0$. If fixed bounds
$0<\ell_i\le w_i\le u_i$ are imposed, feasibility is a linear-program
problem. Feasibility is monotone in $q$. Jointly minimizing $q$ over
$w$ in this representation is not a linear program because $qw_i$
contains products of variables.
\end{proposition}
\begin{proof}
Multiply each row-ratio inequality by its positive denominator $w_i$.
At fixed $q$, every resulting coefficient is fixed, so the inequalities
and box bounds are linear. If a weight vector is feasible at $q$, it
remains feasible at every larger rate because its coordinates are
positive. If $q$ is also variable, $qw_i$ is bilinear, which establishes
the final qualification.
\end{proof}

One may bisect a certified bracket of rates and solve the fixed-rate
feasibility problems. A returned numerical witness should still be
checked with exact arithmetic or validated error bounds before it is used
as a proof of contraction. The checker supplied here verifies a given
rational witness and does not claim to solve the optimization problem.
For
\[
 G=\begin{pmatrix}0&2\\1/8&0\end{pmatrix},\qquad
 w=\begin{pmatrix}4\\1\end{pmatrix},
 \qquad Gw=\tfrac12w,
\]
the unweighted row bound is two, while the weighted bound is $1/2$.
For any positive feasible weights, the two inequalities multiply to
$1/4\le q^2$, so a proposed rate $q=2/5$ is impossible. These exact
equalities and the rejection are included in the regression. Bounds on
$G$ must hold over the claimed domain; a few sampled Jacobians do not
establish them.

\section{Proper scoring must include its probability boundary}

The cross-entropy theorem in \citet{ToricAuditOracle2026} optimizes over
$q(Z)\in(0,1)$ but does not exclude deterministic conditional labels.
When $\eta(Z)=\mathbb P(Y=1\mid Z)$ equals zero or one, its stated
minimizer lies outside that open interval.

\begin{proposition}[Binary log loss with correct endpoint conventions]
\label{prop:oracle-logloss-endpoints}
For $\eta\in[0,1]$, extend
\[
 R_\eta(q)=-\eta\log_2q-(1-\eta)\log_2(1-q)
\]
to $q\in[0,1]$ by assigning zero to a zero-probability term and
$+\infty$ to a positive-probability event assigned probability zero.
Then $R_\eta$ is uniquely minimized at $q=\eta$. Over the restricted
domain $q\in(0,1)$ the minimizer exists precisely when $0<\eta<1$;
for $\eta=0$ or $1$, its infimum is zero and is not attained.
\end{proposition}
\begin{proof}
If $0<\eta<1$, both endpoints have infinite risk. In the interior,
$R_\eta'(q)=(q-\eta)/(q(1-q)\ln2)$, which is negative below
$\eta$ and positive above it. Thus the interior minimum is unique.
If $\eta=0$, the risk is $-\log_2(1-q)$, uniquely minimized at zero.
Every interior $q$ has positive risk and is improved by replacing it
with $q/2$; the risk tends to zero as $q\downarrow0$. If $\eta=1$,
the risk is $-\log_2q$, uniquely minimized at one. Every interior
$q$ is improved by $(1+q)/2$, and the risk tends to zero as
$q\uparrow1$. These cases prove all assertions.
\end{proof}

Applying this argument conditionally gives the Bayes statement for
measurable forecasts with the corresponding admissible range. A sigmoid
with finite logits produces only interior probabilities, so deterministic
targets need not admit a finite-logit optimum. The proper-scoring theorem
does not by itself prove empirical calibration of a restricted, finitely
trained oracle. Calibration and useful held-out control remain measured
properties, with the prefix and artifact accounting established earlier.

\section{Predictive synergy is not determined by geometric overlaps}

The last context-utility theorem of \citet{ToricAuditMemory2026} proposes
an additive information decomposition over independent sheaf components,
with overlap corrections determined by Cech data. The usual chain rule
for mutual information contains conditional statistical terms; independence
of source components alone does not remove predictive synergy.

\begin{proposition}[Independent disjoint components can carry one bit jointly]
\label{prop:context-xor-synergy}
Let $X,Y$ be independent fair bits, let $B=X\mathbin{\mathrm{xor}}Y$,
and let the additive statistic be $T=X+Y$. Then $B$ is determined by
$T$, but
\[
 I(B;X)=I(B;Y)=0,\qquad
 I(B;T)=I(B;X,Y)=1.
\]
The two source components may be stored on disjoint pieces with no
overlap restriction data.
\end{proposition}
\begin{proof}
The four pairs $(X,Y)$ have probability $1/4$. Conditional on either
value of either source bit, $B$ remains a fair bit, so both individual
mutual informations vanish. The target is a fair bit overall and equals
$T\bmod2$, so its conditional entropy given $T$ is zero. It is also
determined by the pair $(X,Y)$. Its unconditional entropy is one bit,
which proves the other equalities. The storage locations of the two
source bits impose no restriction on this joint probability law.
\end{proof}

For arbitrary finite variables and conditioning context $C$, the exact
decomposition is
\begin{align*}
 I(B;X,Y\mid C)
 &=I(B;X\mid C)+I(B;Y\mid X,C)\\
 &=I(B;X\mid C)+I(B;Y\mid C)
   +I(X;Y\mid B,C)-I(X;Y\mid C).
\end{align*}
Both identities follow by expanding conditional mutual information as
differences of conditional entropies and applying the entropy chain rule.
Even if $I(X;Y\mid C)=0$, conditioning on the target can make the
sources dependent. In the example, $I(X;Y\mid B)=1$, and there is
no source overlap whose absence would force this correction to vanish.
Additional probabilistic factorization assumptions are needed to replace
these terms by a particular geometric-overlap model. Compatibility of local
sections alone is insufficient. If $X,Y$ are already determined by the
complete conditioning prefix, all extra conditional information is zero,
as proved earlier; the example applies when the actual predictor has
retained only a smaller context.

The regression checks the complete rational probability table, all
pairwise-independence equalities, deterministic recovery from $T$, and
the conditional dependence given $B$. Its reported one-bit values follow
from these fair and deterministic distributions. The general entropy and
log-loss arguments in this section are written proofs; they are not added
to the count of Lean-checked declarations.

\section{Implications for the research program}

The repaired transport identity provides a shared finite contract for
surface cochains, quiver connectors, quantum instruments and learned
trajectory residuals. It can support a structure-preserving fluid model,
but a Navier--Stokes proof still needs the external construction's analytic
estimates. The corresponding GRH/Hilbert--P\'olya program still requires a
precise arithmetic spectral identification and operator-domain control;
finite self-adjoint matrices do not identify their spectra with zeta zeros.
A positive finite Hodge spectrum likewise does not establish a uniform
continuum Yang--Mills mass gap. Finally, a path-algebra certificate can
authenticate a structural computation only when an actual cryptographic
protocol binds it to a message and a threat model; difficult-looking
tropical equations are not, by themselves, hardness assumptions.

The exact regression script tests 288 generic rational paths, additional
rectangular paths, 128 flat cycles with noncommuting arrows, the prefix
counterexample, positive damping examples and the finite information
example. It also compares tropical dynamic programming against exhaustive
bounded-path enumeration and records the finite-mask and horizon
counterexamples. The new path-multiplicity checks also compare exact
products with an independent enumeration, including changing dimensions,
and retain a dense example attaining the bound. Fixed-rate small-gain
witnesses, the proper-scoring endpoint inequalities and the full XOR
probability table are checked separately. The logarithmic statements are ordinary mathematical
proofs above; the exact regression encodes their positive-input
counterexamples using the monotonicity of the logarithm, without treating
floating-point logarithms as formal certificates. The Lean log records
the exact theorem types and axiom dependencies for its seventeen selected
declarations. Those records state what has been checked and leave the
larger scientific and analytic obligations visible.

\chapter{Quantum--tropical fluid realizations: operators, control and substrates}
\label{ch:quantum-tropical-realizations}

\reviewstatus{This chapter develops an explicitly proposed realization of
ToricGT-style graph models, together with finite mathematical results and an
exact two-qubit feedback example. Its operator, tropical and control layers
have separate contracts. The physical literature supplies demonstrated
building blocks, not an experimental realization of the combined architecture.
Generative substrate design and the mappings to natural systems below are
research proposals. The new finite arguments are written proofs; only the
weighted-gain components identified in the formal type map have Lean coverage.
No new fluid-trained model, fabricated substrate or quantum speedup is reported.}

\section{What a realization must specify}

Here a \emph{quantum--tropical fluid realization} means a model with three
explicit ingredients: quantum states and admissible operations on them;
graph-based routing or a logarithmic asymptotic description; and dynamics
whose relation to a fluid or other physical substrate is quantitatively
specified. It is a proposed interface, rather than a claim that these three
descriptions are interchangeable. A useful realization must identify an
encoding, an evolution, a readout and the approximation error at each
interface.

The reviewed ToricGT implementation contributes graph tokens, endpoint
incidence, finite Weyl operators, a mixture of ordinary and max-plus
attention, and expert routing. The program and Tropical Quivers manuscripts
add proposed algebraic and dynamical constraints. Their exact revisions and
the corrections to selected statements are recorded in
Chapter~\ref{ch:toric-transport-audit}. We retain that pinned meaning of
``ToricGT-style''; the name alone imposes neither a hydrodynamic limit nor a
physical quantum representation.\cite{qtToricGTCode,qtToricGTProgram,ToricAuditQuivers2026}

A concrete specification can be written
\[
 (Q,\mathcal R,\mathcal A,\mathcal E,\Phi,\mathcal M,\pi,\mathcal P).
\]
Here $Q$ is a finite quiver, $\mathcal R$ its intended relations,
$\mathcal A$ the operator algebra, $\mathcal E$ an encoding,
$\Phi$ an admissible evolution, $\mathcal M$ a measurement instrument,
$\pi$ a classical or coherent control law, and $\mathcal P$ a physical
model with units, boundary conditions and parameter uncertainty. A compiler
must say which of these structures it preserves. The fluid discretization
and its continuum consistency belong to $\mathcal P$; the predictor's
training objective belongs to a further statistical specification.

For example, a gentle-algebra relation declaring a composition of two
arrows zero is meaningful for a representation by linear maps on vector
spaces. A composition of trace-preserving channels on a nonzero state space
cannot be the zero map: it still preserves trace. Thus a forbidden path may
be enforced in a classical routing graph, or represented by trace-decreasing
branches of an instrument, but cannot simply be imposed as a zero relation
between unconditional trace-preserving channels. Relations must be checked
in the category in which the proposed hardware actually operates.

\section{A finite operator algebra with conservative quantum transport}

\subsection{Classical position and quantum internal state}

Assign a finite-dimensional Hilbert space $\HH_i$ to each vertex and set
\[
 \mathcal A_Q=\bigoplus_{i\in V}\mathcal B(\HH_i).
\]
A state is a tuple $\rho=(\rho_i)$ with $\rho_i\succeq0$ and
$\sum_i\Tr\rho_i=1$. The mass at vertex $i$ is
$m_i=\Tr\rho_i$; when $m_i>0$, $\rho_i/m_i$ is its conditional internal
state. The direct sum describes a classical location register with quantum
internal degrees of freedom. Coherences between locations instead require
the full algebra on $\bigoplus_i\HH_i$, including off-diagonal blocks.
This choice changes the state space and therefore changes the physics.
Neither finite description is automatically a many-body quantum fluid:
a tensor-product or Fock-space model has additional particle and correlation
structure. An encoding into finitely many blocks must state which
observables it retains and bound truncation or leakage effects.

Let $P_{ij}\ge0$, $\sum_iP_{ij}=1$, and let
$\Phi_{ij}:\mathcal B(\HH_j)\to\mathcal B(\HH_i)$ be completely
positive trace-preserving maps for each allowed transition. Define
\begin{equation}\label{eq:qtf-transport}
 (\mathcal T\rho)_i=\sum_jP_{ij}\Phi_{ij}(\rho_j).
\end{equation}
This finite construction is related to open quantum walks, whose
operator-valued transition rules combine quantum internal states with
classical position records.\cite{qtsAttal2012}

\begin{proposition}[Conservative transport on a direct-sum algebra]
\label{prop:qtf-channel-transport}
Under the preceding assumptions, $\mathcal T$ is completely positive and
trace preserving. Its masses satisfy $m'=Pm$. For Hermitian inputs
$X=(X_j)$,
\[
 \sum_i\|(\mathcal TX)_i\|_1\le\sum_j\|X_j\|_1.
\]
In particular the total trace distance between two states cannot increase.
\end{proposition}
\begin{proof}
Write $\Phi_{ij}(X)=\sum_\alpha K_{ij\alpha}XK_{ij\alpha}^*$ with
$\sum_\alpha K_{ij\alpha}^*K_{ij\alpha}=I_{\HH_j}$. Regard each
$\sqrt{P_{ij}}K_{ij\alpha}$ as a block operator supported from $j$ to
$i$. Their Kraus sum defines a completely positive map, including after
tensoring with any finite ancillary matrix algebra. On input block $j$,
the sum of their adjoint products is
$\sum_iP_{ij}I_{\HH_j}=I_{\HH_j}$, which proves trace preservation.
Taking the trace of each output block gives $m'=Pm$.

For a Hermitian $X_j$, write its positive and negative parts as
$X_j=X_j^+-X_j^-$. Positivity and trace preservation imply
\[
 \|\Phi_{ij}(X_j)\|_1
 \le\Tr\Phi_{ij}(X_j^+)+\Tr\Phi_{ij}(X_j^-)
 =\|X_j\|_1.
\]
The triangle inequality and column sums of $P$ then give
\[
 \sum_i\Bigl\|\sum_jP_{ij}\Phi_{ij}(X_j)\Bigr\|_1
 \le\sum_{i,j}P_{ij}\|X_j\|_1
 =\sum_j\|X_j\|_1.
\]
\end{proof}

Complete positivity is an algebraic constraint on a map and its extensions
to entangled ancillary systems. A numerical test on a few positive input
matrices does not establish it. A Kraus representation with a checked
normalization, or an exactly positive Choi matrix with the correct partial
trace, gives a finite certificate. The Heisenberg dual $\mathcal T^*$ acts
linearly on observables and is unital. These statements concern the same
channel in two dual descriptions.

\subsection{The extra requirements imposed by a fluid interpretation}

A local transition matrix depending on a prescribed velocity field can
transport the masses $m_i$. Adding internal channels can represent
collisions, relaxation or a controlled dissipative feature map. This is a
finite transport model. To obtain a specified fluid equation one must also
identify physical density, momentum, energy, length and time scales; prove
the required conservation laws; specify constitutive assumptions; and
control the discretization and kinetic limits. A positive mass-conserving
map alone has no specified pressure law, viscosity or incompressibility.
The reduction contracts in the preceding chapters state the corresponding
closure and convergence obligations.

For a continuous-time finite quantum model, a standard admissible generator
has the form
\[
 \dot\rho=-i[H,\rho]+\sum_\alpha
 \left(L_\alpha\rho L_\alpha^*
 -\tfrac12\{L_\alpha^*L_\alpha,\rho\}\right),
 \qquad H=H^*.
\]
With finite matrices and time-independent coefficients, this generates a
completely positive trace-preserving semigroup.\cite{qtsLindblad1976}
Local choices of the $L_\alpha$ encode dissipative graph transitions;
off-diagonal Hamiltonian couplings can support coherent transport. A forward
Euler step need not itself be completely positive, even when the continuous
generator is admissible. A numerical realization must retain that
distinction or use a structure-preserving step.

Quantum fluids of light provide a particularly relevant physical family:
interactions, coherent driving and loss lead to rich collective dynamics,
with different descriptions appropriate to mean-field and correlated
regimes. Their equations supply a possible physical model for coupled
tokens; a generic optical neural-network nonlinearity does not automatically
reproduce the correlated quantum regime.\cite{qtsCarusottoCiuti2013,qtsAmo2009}

\section{Toric features, tropical routing and the information they discard}

On $\C^d$, define $X|k\rangle=|k+1\bmod d\rangle$ and
$Z|k\rangle=\omega^k|k\rangle$, where $\omega=e^{2\pi i/d}$. Then
$ZX=\omega XZ$. Products $Z^aX^b$ give finite Weyl features, while
expectations $\Tr(O\rho)$ give real features for Hermitian $O$.
A unitary change of basis transports both $O$ and $\rho$ and leaves their
pairing unchanged. A learned interface may therefore use these finite
operator features while retaining a classical graph controller. Additional
claims about a toric variety, a torus action or a BGG complex require their
own explicitly preserved structure.

The tropical part can enter as a routing optimizer or as a logarithmic
envelope. For $\tau>0$ and $m\ge1$ real scores $s_j$, define
\[
 L_\tau(s)=\tau\log\sum_{j=1}^m e^{s_j/\tau},\qquad
 p_j^\tau=\frac{e^{s_j/\tau}}{\sum_k e^{s_k/\tau}}.
\]

\begin{proposition}[Controlled tropical limit of a positive sum]
\label{prop:qtf-tropical-limit}
For $M=\max_js_j$,
\[
 M\le L_\tau(s)\le M+\tau\log m.
\]
If the maximizer $j_*$ is unique and
$s_{j_*}-s_j\ge\gamma>0$ for all $j\ne j_*$, then
\[
 p_{j_*}^\tau\ge
 \frac{1}{1+(m-1)e^{-\gamma/\tau}}.
\]
\end{proposition}
\begin{proof}
Factor $e^{M/\tau}$ from the sum. The remaining sum is between one and
$m$, giving the first bounds after taking its logarithm. Divide the
normalizing sum for $p_{j_*}^\tau$ by $e^{s_{j_*}/\tau}$.
Each other term is at most $e^{-\gamma/\tau}$, which gives the second
bound.
\end{proof}

This statement concerns a positive exponential sum. With several maximizers
the limiting probability is supported on that set; there is no uniquely
selected route without an extra rule. The reviewed tropical-attention code
uses scores that depend on dot products and applies a max-plus aggregation.
Even when the aggregation is piecewise affine in its supplied scores,
the complete map need not be piecewise affine in all original inputs.

For a product of nonnegative gain matrices, logarithms turn multiplication
along one path into addition. Taking the largest path gives a max-plus
envelope, with $\log0=-\infty$ and zero-weight paths absent. Ordinary
matrix multiplication also sums over paths; its norm
requires the path-multiplicity factor proved in
Chapter~\ref{ch:toric-transport-audit}. Signed or complex amplitudes add
phase cancellation as another issue. A tropical magnitude envelope can
upper-bound selected effects once these factors are supplied, but does not
retain the interference pattern or determine entanglement.

Min-plus arithmetic also arises directly in control. A deterministic
finite-horizon route problem has the Bellman recursion
\[
 V_n(i)=\min_{a\in\mathcal U(i)}
 \{c_n(i,a)+V_{n+1}(f(i,a))\}.
\]
This is a precise optimization layer that can select graph operations.
A stochastic, risk-sensitive or entropy-regularized controller has a
different recursion, whose transition probabilities and normalization must
be retained. A zero-temperature approximation can reduce exploration but
also creates switching boundaries. The common gain certificate below is
useful only if it covers every allowed branch and the effects of switching.

\section{Nonlinear dynamics on an operator state space}

\subsection{Where the nonlinearity lives}

An operator algebra is a linear space with a product and an adjoint. There
are several mathematically distinct ways to obtain nonlinear dynamics
associated with it. First, a classical controller $z$ can select a channel
or a generator:
\[
 \rho_{n+1}=\Phi_{\theta(z_n)}(\rho_n),\qquad
 z_{n+1}=f(z_n,y_n),\qquad
 y_n=\mathcal M_n(\rho_n).
\]
For a fixed classical setting, $\Phi_{\theta}$ is linear on density
operators. The coupled controlled system can be nonlinear. If $y_n$
denotes an ideal expectation rather than an actual measurement record,
the displayed recursion is an ensemble-level model and needs a readout
approximation before it describes hardware.
For actual measurement records, the conditional measurement backaction
must be included in the state update, or readout must use separately
prepared copies. It cannot be omitted while retaining a claim about the
same undisturbed quantum system.

Second, a measurement instrument has unnormalized branches
$\rho\mapsto K_y\rho K_y^*$. Each is linear and completely positive.
Conditioning on a nonzero-probability outcome divides by
$\Tr(K_y\rho K_y^*)$ and produces a nonlinear normalized update. A
physically complete account retains the probabilities and the classical
outcome register. Replacing that instrument with a deterministic normalized
map on an unknown state changes the operation.

Third, effective mean-field equations may contain a state-dependent
Hamiltonian, for example $\dot\rho=-i[H(\rho),\rho]$, or a cubic
order-parameter equation. Such equations approximate selected observables
or many-body regimes under stated assumptions. Their nonlinearity does not
by itself supply a fundamental nonlinear quantum channel. Driven polariton
and atomic condensate models make this distinction experimentally relevant.
\cite{qtsCarusottoCiuti2013,qtsBarontini2013}

These three cases give a practical interpretation of a nonlinear operator
perspective: learn and control dynamics \emph{on} a physically admissible
state space while retaining the linear completely positive operations
appropriate to each fixed experimental setting. Positivity, trace, symmetry,
causality and accessible controls become architectural constraints.

\subsection{Koopman lifting and its closure error}

For a deterministic state-space map $F:\mathcal X\to\mathcal X$, the
Koopman operator is $(Uf)(x)=f(F(x))$. It is linear in the observable $f$
even when $F$ is nonlinear. On an algebra of ordinary scalar observables it
also preserves products. For a stochastic transition, conditional
expectation instead gives a positive unital operator, generally without
product preservation. Neither construction is automatically the
Heisenberg dual of a channel on the noncommutative algebra $\mathcal A_Q$:
their domains and their notions of positivity must first be matched.

A finite dictionary $\psi(x)$ may satisfy
\[
 \psi(F(x))=K\psi(x)+r(x).
\]
An exact finite representation requires $r=0$ on the claimed domain.
Learning $K$ from trajectories ordinarily leaves a residual, and a small
training residual is not a uniform domain bound. Finite invariant
subspaces that also reconstruct the state impose substantial restrictions,
as the primary Koopman-control analysis explains.\cite{qtsBrunton2016}
The residual belongs in an error recurrence such as
Theorem~\ref{thm:weighted-gain-propagation}. This connects operator learning
to a checked control budget without declaring a finite closure exact.

\section{Control, observability and robustness}

\subsection{A common gain for coupled classical and quantum blocks}

Let $d_\rho$ measure quantum-state error, for example trace distance, and
let $d_z$ measure controller error. Suppose the actual closed-loop map
satisfies on an invariant domain
\[
 \begin{pmatrix}d_\rho'\\d_z'\end{pmatrix}
 \le
 G\begin{pmatrix}d_\rho\\d_z\end{pmatrix}
 +\begin{pmatrix}\eta_\rho\\\eta_z\end{pmatrix},\qquad
 G=\begin{pmatrix}a&b\\c&d\end{pmatrix}\ge0.
\]
For $q>\max(a,d)$ and $b,c>0$, the certificate
$Gw\le qw$ with $w_\rho,w_z>0$ is equivalent to
\[
 \frac{b}{q-a}\le\frac{w_\rho}{w_z}
 \le\frac{q-d}{c}.
\]
Indeed the first row rearranges to the left inequality and the second to
the right. Feasibility is therefore equivalent to
$bc\le(q-a)(q-d)$ at that rate. If $a,d<1$ and
$bc<(1-a)(1-d)$, continuity supplies a feasible $q<1$ sufficiently close
to one. Cases with zero coupling can be checked directly using $Gw\le qw$.
These formulas expose the cost of feedback: individually stable components
can still have an excessive loop gain.

Once the gains are justified, a common $w$ and $q$ give the full finite-time
budget of Theorem~\ref{thm:weighted-gain-propagation}. An optimizer may
propose $w$; the independent rational checker verifies the proposed
inequalities. Separate certificates can cover model discrepancy, channel
approximation, integration error and readout error as explicit residuals.
For a uniform weighted perturbation $\delta$ and $q<1$, the eventual
bound is $\delta/(1-q)$. A smaller $q$ improves this robustness bound,
while rapid contraction also suppresses memory of past inputs. Reservoir
design therefore needs a task-dependent tradeoff between memory, separation
of inputs and noise sensitivity.\cite{qtsFujiiNakajima2017}

\subsection{Accessible controls and finite readout}

For a closed finite quantum system, accessible Hamiltonians generate a Lie
algebra that constrains the reachable unitaries. Full unitary control still
preserves the eigenvalues of a density matrix. Changing its spectrum
requires measurement, dissipation or an enlarged system and subsequent
discarding. Thus the ability to fit a target matrix numerically does not
prove that a specific device can implement it.\cite{qtsAlbertiniDAlessandro2003}
For an open device, the allowed jump operators, rates, locality, noise and
bandwidth similarly constrain attainable channels.

Observability has a corresponding cost. One cannot obtain every expectation
of an unknown state from one undisturbed specimen. A hardware proposal must
specify repeated preparations, ensemble access, weak measurement or an
instrument with a state estimator. A statistical error guarantee over a
finite number of measurements can be converted into a deterministic
residual bound on the event that those estimates are accurate. The event's
confidence, the number of shots and the finite time horizon remain part of
the result; an unlimited-horizon guarantee does not follow automatically.

For example, suppose $M$ measurement settings are fixed in advance and
each uses $N$ independent shots with outcomes in $[-1,1]$. Hoeffding's
inequality and a union bound give simultaneous error at most $\varepsilon$
with probability at least $1-\delta$ when
\[
 N\ge\frac{2}{\varepsilon^2}\log\frac{2M}{\delta}.
\]
For one setting, the two-sided failure probability is at most
$2e^{-N\varepsilon^2/2}$; adding these $M$ bounds proves the assertion.
The confidence concerns sampling error for those fixed settings, not
calibration bias or arbitrary temporally correlated reservoir outputs.
Adaptive feedback requires a corresponding conditional sampling argument
and a specified finite horizon.\cite{qtsHoeffding1963}

Model predictive control can optimize a finite horizon subject to channel
admissibility, actuator limits and a certified invariant set. For a
classical energy model
\[
 \dot x=(J(x)-R(x))\nabla E(x)+B(x)u,\qquad
 J^T=-J,\quad R\succeq0,
\]
the chain rule gives
$\dot E=-\nabla E^TR\nabla E+y^Tu$ with
$y=B^T\nabla E$. This is a useful energy-balance constraint for a
fluid, elastic or electrical surrogate. A quantum dissipative model needs
its own thermodynamic assumptions: an arbitrary driven channel need not
decrease a chosen energy or increase entropy. The proposed controller
must preserve the balance appropriate to its physical model.

\section{An exact nonlinear quantum-state feedback example}

The following small construction instantiates the control contract without
assuming a gain bound from sampled derivatives. It uses a tuple of two
qubit density matrices; it does not encode entanglement between the two
blocks. Let $Z=\diag(1,-1)$,
\[
 z_j=\Tr(Z\rho_j),\quad
 B=\begin{pmatrix}0&2\\1/8&0\end{pmatrix},\quad
 b=\begin{pmatrix}1/4\\0\end{pmatrix},\quad
 u_i=\operatorname{clip}((Bz+b)_i,-1,1).
\]
Set $a=1/4$, $\tau(u)=(I+uZ)/2$, and define
\begin{equation}\label{eq:qtf-feedback}
 F_i(\rho)=a\rho_i+(1-a)\tau(u_i).
\end{equation}
The clipped controller is nonlinear but globally Lipschitz. Its expectations
are assumed known exactly for the ideal model. A device implementation
would estimate them through the readout process just described.

\begin{proposition}[A physical channel family and a contracting feedback law]
\label{prop:qtf-feedback-example}
For every fixed $u\in[-1,1]$,
\[
 \Phi_u(X)=aX+(1-a)\tau(u)\Tr X
\]
is completely positive and trace preserving. The ideal feedback map $F$
preserves pairs of density matrices and satisfies, with
$D_i=\tfrac12\|\rho_i-\sigma_i\|_1$,
\[
 D_i'\le\sum_jG_{ij}D_j,\qquad
 G=\begin{pmatrix}1/4&3/2\\3/32&1/4\end{pmatrix},\qquad
 G\begin{pmatrix}4\\1\end{pmatrix}
 =\tfrac58\begin{pmatrix}4\\1\end{pmatrix}.
\]
It therefore contracts the weighted maximum trace distance at rate $5/8$.
Its fixed point has Bloch vectors $(0,0,1/3)$ and $(0,0,1/24)$.
The feedback map itself is not affine in the pair of input states.
\end{proposition}
\begin{proof}
Since $|u|\le1$, $\tau(u)$ is a density matrix. In input--output order,
the unnormalized Choi matrix is
\[
 J(\Phi_u)=a|\Omega\rangle\langle\Omega|
 +(1-a)I\otimes\tau(u),\qquad
 |\Omega\rangle=|00\rangle+|11\rangle.
\]
Both terms are positive, and tracing over the output gives
$aI+(1-a)I=I$. Equivalently the replacement map has Kraus operators
$\sqrt{t_k}|k\rangle\langle l|$, where $t_k$ are the diagonal entries
of $\tau(u)$ and $k,l\in\{0,1\}$. Scaling these replacement operators
by $\sqrt{1-a}$ and adjoining $\sqrt a I$ gives the same channel.
Convexity then proves state-space
invariance for every chosen control.

Measurement duality gives $|\Delta z_j|\le2D_j$. Clipping is
one-Lipschitz and
$\tfrac12\|\tau(u)-\tau(v)\|_1=|u-v|/2$. Thus
\[
 D_i'\le aD_i+\frac{1-a}{2}|\Delta u_i|
 \le aD_i+(1-a)\sum_j|B_{ij}|D_j.
\]
Substitution gives $G$, and direct multiplication gives its certificate.
The state space is complete in this finite-dimensional metric, so the
strict contraction has at most one fixed point and converges to it.
At a fixed point the transverse Bloch coordinates vanish; the displayed
$z$ values solve $z=Bz+b$ and lie inside the clipping interval. They
therefore give existence as well as uniqueness.

For nonaffinity, hold the first input at $z_1=0$ and choose second inputs
with $z_2=0$ and $z_2=1$. The first output $z$ values are $3/16$ and
$3/4$. Their average is $15/32$. At the midpoint input $z_2=1/2$,
the clipped control is one and the first output is $3/4$. The difference
is $9/32$, so a single affine channel cannot equal this ideal feedback
map on all inputs.
\end{proof}

The fixed-control channels are physically admissible. The ideal feedback
recursion describes what happens when a controller has the required
expectations; it is not an instruction for obtaining them without readout
cost. The associated exact implementation uses rational Bloch coordinates
and rational Choi matrices. It retains a coherent initial state, tests an
actually saturated step, and checks trace-distance bounds by squaring
nonnegative quantities, avoiding floating-point square-root tolerances.
These replacement channels are generally nonunital:
$\Phi_u(I)=I+(1-a)uZ$. The entropy-monotonicity statement for the earlier
unital heat channel consequently does not transfer to this example.

There is also an exact tropical comparison. Weighted normalization gives
\[
 K=\diag(w)^{-1}G\diag(w)
 =\begin{pmatrix}1/4&3/8\\3/8&1/4\end{pmatrix}.
\]
Every row of $K^n$ sums to $(5/8)^n$, whereas the largest weight of one
length-$n$ path is $(3/8)^n$. The latter is attained by alternating the
vertices. The difference comes from adding the other paths, exactly as in
the path-multiplicity analysis. The proposed rate $3/5$ fails the supplied
weighted certificate and is retained as negative replay evidence.

Finally, if each measured expectation has deterministic error at most
$\varepsilon$, the extra weighted one-step discrepancy is at most
$3\varepsilon/16$: the two component bounds after division by $w_i$
are $3\varepsilon/16$ and $3\varepsilon/64$. Therefore the same
residual theorem gives the asymptotic envelope
\[
 \frac{3\varepsilon/16}{1-5/8}=\frac{\varepsilon}{2}.
\]
For $\varepsilon=1/100$, this is $1/200$. This is conditional error
control under the stated readout bound. It is not an experimentally
measured noise floor or a substitute for a finite-shot confidence analysis.

\section{Candidate engineered and natural substrates}

A substrate is relevant when its controlled degrees of freedom implement
the specified encoding, operation and readout with bounded error. Its
microscopic quantum origin alone does not establish useful coherent quantum
processing. The following examples identify both the usable mechanism and
the experiment needed to connect it to the proposed architecture.

\begin{longtable}{p{0.20\textwidth}p{0.34\textwidth}p{0.36\textwidth}}
\toprule
Substrate & Candidate role & Evidence and realization test\\\midrule
\endhead
Programmable photonic networks &
Interferometers implement coherent mode mixing; phase shifters and
detectors can supply a graph operator and readout. A separate controller
can update routing. &
The cited programmable-transport experiment used continuous-wave laser
illumination and intensity readout. A few-photon channel would need its
own nonclassical-state specification and validation, alongside loss and
control-drift measurements.
\cite{qtsHarris2017}\\[6pt]
Exciton--polariton fluids &
Coupled cavities or spatial modes supply driven, lossy, interacting
fields. They are candidates for nonlinear reservoirs and graph-coupled
fluid features. &
Quantum-fluid and superfluid phenomena motivate the mechanism. A
realization must identify its mean-field or correlated regime, its
input/output map, and the controller's error bounds.
\cite{qtsCarusottoCiuti2013,qtsAmo2009}\\[6pt]
Ultracold atomic gases &
Trapped modes or lattice sites provide quantum matter with controllable
couplings and, in selected settings, engineered local dissipation. &
Dissipative many-body dynamics have been studied experimentally. A
ToricGT mapping would additionally need accessible graph couplings,
repeated state preparation and an observable task with a calibrated
readout budget.\cite{qtsBarontini2013}\\[6pt]
Superconducting cavity circuits &
Microwave modes, nonlinear elements and engineered loss can implement
bosonic channels, memory and controlled dissipation. &
Engineered dissipation has confined dynamics to a quantum-state manifold.
The proposed graph realization must quantify truncation, leakage,
decoherence and the attainable channel family.\cite{qtsLeghtas2015}\\[6pt]
Magnonic and magnetic networks &
Spin-wave propagation and nonlinear responses can provide a compact
dynamical reservoir with graph-like couplings. &
Magnetic reservoir experiments support this computing mechanism. The
usual wave-regime demonstration establishes a physical reservoir;
quantum-state processing would need its own preparation and coherence
evidence.\cite{qtsNamiki2023}\\[6pt]
Soft bodies and fluid--structure systems &
Elastic motion, damping and hydrodynamic coupling can distribute
nonlinear memory across a body, with sensors and actuation supplying
the interface. &
Soft-body reservoir experiments demonstrate task-relevant dynamics.
Natural morphologies motivate graph constraints, but every generated
body needs a measured transfer map and robustness test.
\cite{qtsNakajima2015}\\[6pt]
Photosynthetic molecular networks &
Excitonic couplings, vibrations and environmental relaxation offer a
natural example of open-system energy transport. &
The functional role and character of observed coherence require care.
Electronic and vibrational contributions, disorder and open-system
conditions must be distinguished before proposing an information
processor.\cite{qtsDuan2017,qtsHarushDubi2021}\\\bottomrule
\end{longtable}

Natural vascular, porous and transport networks provide another source of
geometries and conservation constraints for a classical graph-fluid
realization. Their branching and feedback suggest a design grammar;
they do not supply a trained ToricGT model by identification alone.
Likewise, a crystal or a biological assembly discovered in nature can
inspire a coupling pattern without providing arbitrary programmable
control. The useful question is which functions the accessible dynamics
can implement and which invariants survive the readout.

There is no universal ordering of these substrates. Interferometric
systems favor controlled linear mixing; strongly interacting quantum
platforms favor selected many-body phenomena; dissipative mechanical and
magnetic reservoirs may provide economical nonlinear memory. A comparison
needs a common task, input bandwidth, output accuracy, reset protocol,
energy accounting and training budget. Device operation, control and
measurement must all enter the cost.

\section{De novo substrate design with generative models}

\subsection{Generate a physical candidate for a specified operator}

A generative model can propose a material or geometry conditioned on a
target response. Inorganic crystal generation and generative photonic
inverse design provide concrete examples of this approach. Their outputs
are candidates for a physical validation pipeline; the generative likelihood
is not an operator-realization certificate.\cite{qtsZeni2025,qtsJiangFan2019}

Let $s$ describe composition, geometry, connectivity and fabrication
parameters, and let $u$ describe accessible controls. A forward model
produces an effective evolution $\mathcal F_{s,u}$ and readout
$\mathcal M_s$. One can seek
\[
 \min_{s,u,w,q}\quad
 L_{\rm task}(s,u)+\lambda_E E(s,u)+\lambda_C C(s)
 +\lambda_R R_{\rm uncertainty}(s,u)
\]
subject to physical admissibility, an observable target interface, and
$G(s,u)w\le qw$ with $w>0$, $q<1$ when contraction is appropriate.
The objective terms need measured definitions: prediction error,
energy per usable output, fabrication or control cost, and a robustness
penalty. A generative prior can constrain plausible $s$ without replacing
the physical constraints.

The specification should begin with a target channel, scattering map,
response kernel or controlled fluid update, rather than with a visual
shape. For an optical device, relevant constraints include a realizable
material response, loss, bandwidth and fabrication tolerances. For a
magnetic, elastic or fluid device, stability, conservation laws, boundary
conditions and sensor coupling matter. For a proposed quantum device,
state preparation, available interactions and admissible measurement are
part of the target. Complete positivity constrains a reduced channel;
additional locality and controllability constraints determine whether the
chosen substrate can realize it.

\subsection{A falsifiable co-design loop}

The following sequence separates design exploration from validation.
\begin{enumerate}
\item Specify the graph task, allowed algebraic relations, encoding,
observables, error tolerance and resource budget. Decide whether the
target requires coherent quantum states or only a wave/dynamical response.
\item Generate candidate structures conditioned on those requirements.
Apply exact geometry, connectivity and composition constraints where
available, and retain the generator version and random seed.
\item Evaluate candidates with the appropriate physical solver and
uncertainty model. Reject violations of conservation, positivity,
stability or attainable controls; a good surrogate score cannot override
these checks.
\item Identify a reduced input--state--output model on a stated domain.
Measure its residual on held-out conditions and obtain a justified gain
envelope. If only sampled derivatives are available, record that empirical
scope rather than promoting it to a uniform bound.
\item Co-design the controller and positive weights. Replay finite
rational certificates independently, and propagate solver, readout and
fabrication errors through the common-gain budget.
\item Evaluate prospective fabrication variation, then fabricate or
assemble a selected candidate only under an actual experimental program.
Calibrate it, test closed-loop behavior and compare with the predicted
operator. A simulation-only result remains labeled as such.
\item Compare against conventional numerical solvers and simpler
physical or classical learning baselines on the same task. Include
state preparation, measurement, reset, calibration and controller costs.
\end{enumerate}

MatterGen establishes that conditional generative modeling can propose
inorganic crystals with target properties, including an experimentally
tested candidate. It does not supply the mapping from a desired quantum
channel to a fabricable quantum device.\cite{qtsZeni2025}
Generative photonic inverse design directly targets an optical response,
but its particular demonstrated response still needs an additional mapping
to this graph-fluid architecture.\cite{qtsJiangFan2019}
Those missing mappings are useful research questions: they specify what a
new project must actually demonstrate.

Natural systems can enter this loop as a prior on topology, material motifs
or dynamics. A learned grammar might favor sparse transport networks,
hierarchical elastic couplings or molecular geometries. The design must
then be evaluated with the same operator, control and readout contracts
as a wholly synthetic candidate. This makes biomimetic inspiration
testable without treating an analogy as a realization.

\section{Connections to the wider research program}

The construction offers a practical role for the project's order and
surface-algebra language: encode legal compositions and symmetries, select
observable features, and demand that a realization preserve specified
relations. Quantum information supplies admissible state transformations
and operational distances. Tropical mathematics supplies routing and
logarithmic bounds. Control theory connects these finite operations to
robust dynamical behavior. Each connection has an explicit map and an
error obligation.

For Hilbert--P\'olya research, a controllable finite operator family can
test spectral statistics or evaluate a proposed discretization. A finite
fit to zeros does not establish the required self-adjoint operator,
spectral identification or analytic limiting theorem. For Yang--Mills,
an engineered finite dissipative or Hamiltonian gap is a property of that
system and has no automatic implication for the continuum mass-gap
problem. For cryptography, a generated physical response may be a useful
device primitive, but a security claim additionally requires a threat
model, entropy or hardness assumptions, extraction and implementation
analysis. The existing cryptographic implementations and their security
scope remain separate evidence.

For Navier--Stokes, the useful next experiments are bounded: preservation
of discrete invariants, closure-error prediction, controllable manufactured
solutions and comparisons with conventional solvers at stated resolution.
The externally released forced singularity theorem retains its own
continuum hypotheses and independent-build status. The present finite
control construction neither supplies a stronger continuum theorem nor
removes the unresolved independent upstream build.

The immediate implementation milestone is the exact feedback example and
its independent gain replay. The next research milestone is a trained
graph-fluid surrogate with a documented invariant domain and held-out
error budget. A physical milestone requires one specified substrate,
calibrated encoding and readout, measured operator agreement and a fair
resource comparison.

\appendix
\chapter{Fluid verification, reproducibility and proof dependencies}

\section{Current verification boundary}

The final local companion has 121 selected declarations across seven modules,
including 16 reduction contracts, 17 transport declarations and six weighted-gain contracts. All seven were
freshly checked with Lean 4.24.0 and the pinned mathlib revision, with zero
warnings and only the disclosed standard foundational axioms. Exact types,
source hashes and logs are retained under \texttt{lean/final-verification}.
The earlier 82-declaration fluid-era records below remain identifiable
historical checks of the first four modules. Neither record completes the
separate upstream Navier--Stokes full build, which remains blocked by missing
rate-limited dependency artifacts.

\label{ch:fluid-verification}

\section{What was checked and what each check means}

The fluid extension has three distinct verification layers. The manuscript gives explicit proofs of the displayed finite statements and analytical interfaces. The local Lean package checks selected statements with precisely typed hypotheses. The executable examples compare independently implemented finite calculations. None of these layers is interchangeable with the complete upstream continuum development.

The additional file \path{extension/lean/FluidInterfaces.lean} contains 34 theorem declarations. It was elaborated successfully by Lean 4.24.0 against the same pinned mathlib revision as the earlier companion,
\begin{center}\small\texttt{f897ebcf72cd16f89ab4577d0c826cd14afaafc7}.\end{center}
The check returned exit code zero and no warnings. Every declaration has a printed type and an axiom report. The only reported foundations are \texttt{propext}, \texttt{Classical.choice} and \texttt{Quot.sound}; no admitted theorem or project-specific axiom is introduced. The verified source SHA-256 is
\begin{center}\scriptsize
\texttt{9ac752e3450fd07f9fb23cb4d0283ecee1b70d7a3188ece1cdcff5d8b66934b2}.
\end{center}
At that earlier checkpoint, all four modules were freshly rechecked together, with unchanged source hashes and zero warnings. The individual logs and declaration-by-declaration axiom comparisons are retained in \path{extension/review/navier-stokes/local-kernel-verification.json}. That earlier package contained 82 selected declarations: 25 original interfaces, seven quantum interfaces, sixteen diagram interfaces and these 34 fluid interfaces. This is a count of exact declarations, not a claim that 82 complete manuscript theorems, every numerical implementation, or the continuum singularity construction have been formalized here.

\begin{longtable}{p{.25\textwidth}p{.08\textwidth}p{.56\textwidth}}
\toprule Family & Count & Checked content and its boundary\\\midrule
\endhead
Finite energy & 11 & Finite real matrix work identities, pressure cancellation under the stated constraint, nonnegative factorized viscosity, and an actual time derivative of coordinate kinetic energy. No continuum regularity conclusion.\\
Orthogonal reduction & 5 & Kernel residual, orthogonality, Pythagorean energy split, energy nonincrease and best approximation for a symmetric idempotent matrix. No closure or trajectory error theorem.\\
Chain continuity & 2 & Boundary flux has zero subsequent boundary, and adding such a flux preserves the divergence. The complex equation is explicit.\\
Schur elimination & 1 & Equivalence of a block system and its eliminated system when the required block is a linear equivalence. No assumption that every block is invertible.\\
Discrete memory & 3 & All-step hidden-state expansion, resolved memory identity and the obstruction to an instantaneous closure. No decay assumption is inserted into the recurrence.\\
Forbidden composition & 1 & A strict lower block, then a block diagonal map, then another strict lower block has zero composite. This is one exact operator pattern from the Volterra discussion.\\
Quadratic order & 11 & The faithful order representation, addition and multiplication, norm determinant, matrix polynomial, determinant of every power, and absence of a commuting real linear complex structure. The topological covering interpretation is proved separately in the manuscript.\\
\bottomrule
\end{longtable}

The component map in \path{extension/review/navier-stokes/formalization-map.json} and the main claim register link these declarations to the statements they actually support. A matrix energy identity is a component of the finite ODE existence argument, not a Lean proof of that entire existence theorem. The exact discrete memory recurrence is not a Lean proof of the continuous integral formula. Likewise the determinant of $J^k$ is not, by itself, a formalized covering-space degree theorem.

\section{The separate upstream Navier--Stokes formalization}

The released project uses a different toolchain: Lean 4.34.0-rc2 and its own pinned mathlib dependency. Its repository is pinned at
\begin{center}\small\texttt{f9e8bc5b38b6e212696e8a30e3e91517af887bbd}.\end{center}
The audited final solution target has a local import closure of 609 modules. A comment- and string-aware source scan of this closure finds no admitted placeholders or project axiom declarations, and no challenge module is imported by that closure. This is useful source evidence but is not a kernel check.

A fresh upstream build was attempted with the exact verified compiler archive. Dependency-cache retrieval encountered HTTP 429 responses. The subsequent direct target check stopped because imported build artifacts were missing, before theorem-body elaboration. The full upstream kernel check and independent Comparator check therefore remain uncompleted in this work. The supplied \path{release-build-verification.json} records these outcomes, while \path{release-verify.py} supplies a pinned reproduction route. A failed dependency download neither verifies nor refutes the theorem. The local 34-declaration result is kept entirely separate from this upstream result.

The complete 166-page paper was read across the two mathematical reviewers, with every page attributed in \path{release-source-manifest.json}. That reading supports the exposition and the precise statement of the forced whole-space and periodic results. It does not substitute for the missing fresh formal build, nor does it turn a proof of forced breakdown into an unforced theorem.

\section{Independent finite computational checks}

The new program \path{extension/code/verify_fluid_interfaces.py} constructs exact integer cell complexes for six periodic square grids, with $L=3,\ldots,8$. It checks the zero boundary composite and the expected finite ranks. It independently evaluates symbolic Taylor--Green equations and dense complex-matrix circuit operations. Its machine-readable output is \path{extension/verification/fluids-python.json}.

The browser core is separately checked by \path{extension/site/verify-fluids.cjs}, which consumes independent Python reference data. The recorded comparison covers sixteen circuit inputs, 64 parameter-feature derivative cases and 8,865 complete toral-fibre points. The maximum amplitude discrepancy is about $2.22\times10^{-16}$; the maximum character-pullback discrepancy is about $1.01\times10^{-13}$; and the heat-channel probability discrepancy is about $3.47\times10^{-18}$. These are floating-point discrepancies, not exact-real proof bounds. The finite-difference momentum residual decreases from $3.03\times10^{-6}$ to $7.57\times10^{-7}$ to $1.89\times10^{-7}$ under successive halvings of the step, consistent with the method's second-order error on this manufactured solution.

The channel checks distinguish the live diagonal-state simulation from the separate coherent-state matrix test. Positivity, trace preservation and complete positivity follow analytically from the Kraus form in Chapter~\ref{ch:quantum-fluid}; finite Choi and density-matrix calculations test the implemented instances. The circuit is simulated classically and is untrained. The circuit register and the periodic position register are distinct systems. No NVIDIA acceleration is required or claimed for these small new examples; the earlier GPU surface-code and ToricGT experiments retain their separate recorded execution evidence.

\section{Reproduction commands}

From the repository root, the finite computations can be repeated with:
\begin{lstlisting}
python3 extension/code/verify_fluid_interfaces.py
node extension/site/verify-fluids.cjs
\end{lstlisting}
The first script requires the numerical and symbolic dependencies stated in its source. The second uses the generated reference data. The local formal check can be repeated using the installed pinned toolchain:
\begin{lstlisting}
python3 extension/review/navier-stokes/formalization-verify.py \
  --toolchain /absolute/path/to/lean-4.24.0-linux \
  --mathlib /absolute/path/to/mathlib4-v4.24.0
\end{lstlisting}
Alternatively, from a correctly resolved Lake package, run
\begin{lstlisting}
lake env lean FluidInterfaces.lean
\end{lstlisting}
The manuscript entry points share the same chapter sources. With the listed TeX packages installed:
\begin{lstlisting}
python3 extension/manuscript/build_pdf.py
python3 extension/manuscript/build_pdf.py \
  --main fluid-extension --min-pages 1
\end{lstlisting}
The full monograph keeps the 100-page lower gate. The standalone supplement is not padded to a target length. Both builds reject undefined references, missing characters and overfull boxes, and retain their actual compiler logs and source hashes. Visual review is recorded separately from those mechanical gates.

\section{The complete selected fluid formalization}

The following is the source actually checked, including its type and axiom audit block. Its finite types are part of the theorem statements; replacing them with function spaces would create a different theorem and new analytic obligations.

\lstinputlisting{../lean/FluidInterfaces.lean}

\backmatter
\bibliographystyle{plainnat}
\bibliography{references,ns-microscopic,ns-release,ns-structures,ns-spin,reduction-contracts,toric-audit,quantum-tropical-substrates}
\end{document}
